% Maîtrise des comportements interindividuels — SVT 1ère A — Cours signé OnBuch+
% Programme officiel MINESEC. Compiler avec : tectonic NA03-maitrise-comportements-interindividuels.tex
\newcommand{\DOCMATIERE}{SVT}
\newcommand{\DOCNIVEAU}{1ère A}
\newcommand{\DOCMODULE}{Module 1 : Le Monde Vivant}
\newcommand{\DOCLECON}{NA03}
\newcommand{\DOCTITRE}{Maîtrise des comportements interindividuels}
% OnBuch+ — préambule « cours signés » ENRICHI (compilé avec tectonic / XeLaTeX)
% Système visuel « Pop » d'OnBuch+ : fond crème (jamais blanc), bordures encre
% épaisses, ombres « dures » décalées, titres Archivo Black en capitales.
% Version MATHÉMATIQUES : graphes (pgfplots/tikz), boîtes pédagogiques
% (définition, propriété, méthode, attention, le savais-tu, exercice résolu,
% exercice, TP, bilan), page de garde et signature OnBuch+ ancrée en bas.
%
% Macros attendues dans main.tex AVANT \input{preamble} :
%   \DOCMATIERE  \DOCNIVEAU  \DOCMODULE  \DOCLECON (numéro)  \DOCTITRE
\documentclass[11pt]{article}

\usepackage{fontspec}
\usepackage[french]{babel}
\usepackage[a4paper,margin=2cm,top=3.4cm,bottom=2.8cm,headheight=1.4cm,footskip=1.3cm]{geometry}
\usepackage{amsmath,amssymb,amsfonts}
\usepackage{mathtools} % \xmapsto, \coloneqq... souvent utilisés par les modèles
\usepackage{cancel} % simplifications barrées (bilans de forces)
% \abs{z} (module, valeur absolue) et \norm{u} (norme), utilisés par les modèles
\providecommand{\abs}{}\let\abs\relax\DeclarePairedDelimiter{\abs}{\lvert}{\rvert}
\providecommand{\norm}{}\let\norm\relax\DeclarePairedDelimiter{\norm}{\lVert}{\rVert}
\usepackage[table]{xcolor} % [table] : fournit aussi \rowcolors (alternance de lignes)
\usepackage{pagecolor}
\usepackage{tikz}
\usetikzlibrary{arrows.meta,positioning,calc,shapes.geometric,shapes.misc,shapes.arrows,decorations.pathmorphing,decorations.pathreplacing,decorations.markings,patterns,shadows,mindmap,trees,fit,backgrounds,matrix,intersections,angles,quotes}
\usepackage{pgfplots}
\usepackage[version=4]{mhchem} % équations nucléaires \ce{^{235}_{92}U}
\usepackage[european,siunitx]{circuitikz} % schémas électriques
\pgfplotsset{compat=1.18}
\usepgfplotslibrary{fillbetween,groupplots} % aires entre courbes (calcul intégral)
\usepackage{enumitem}
\usepackage{fancyhdr}
% [most] charge toutes les bibliothèques tcolorbox (listings, minted, raster,
% documentation...) et gonfle inutilement la mémoire TeX sur un document long
% (~300 boîtes) : on ne charge que ce qui est réellement utilisé.
\usepackage[skins,breakable]{tcolorbox}
\usepackage{graphicx}
\usepackage{colortbl}
\usepackage{booktabs}
\usepackage{tabularx}
\usepackage{diagbox} % case d'en-tête diagonale des tableaux à double entrée
% Colonnes extensibles centrée / gauche / droite (C, L, R), souvent utilisées par les modèles
\newcolumntype{C}{>{\centering\arraybackslash}X}
\newcolumntype{L}{>{\raggedright\arraybackslash}X}
\newcolumntype{R}{>{\raggedleft\arraybackslash}X}
\usepackage{array}
\usepackage{multicol}
\usepackage{siunitx}
% parse-numbers=false : accepte aussi \SI{2,5 \times 10^{-3}}{...}
\sisetup{locale=FR,output-decimal-marker={,},group-minimum-digits=4,parse-numbers=false}
\DeclareSIUnit\ohm{\ensuremath{\symup{\Omega}}} % U+2126 absent de la police
\DeclareSIUnit\division{div} % réglages d'oscilloscope (ms/div, V/div)
\DeclareSIUnit\tr{tr} % tours (vitesse de rotation, ex. \SI{1500}{\tr\per\minute})
\DeclareSIUnit\ch{ch} % cheval-vapeur (puissance moteur, unité usuelle hors SI)
\DeclareSIUnit\franccfa{FCFA} % franc CFA (coûts, contexte camerounais)
\DeclareSIUnit\dioptre{\ensuremath{\delta}} % vergence d'une lentille (dioptrie)
\usepackage{qrcode}
% \degree : commande fréquemment utilisée par les modèles pour un angle ou une
% température en dehors de \SI{}{} (non fournie par les paquets chargés).
\providecommand{\degree}{^{\circ}}
% \qed (fin de démonstration), utilisé par les modèles sans amsthm
\providecommand{\qed}{\hfill$\square$}
% \barre : double barre des valeurs interdites dans un tableau de variations (tkz-tab)
\providecommand{\barre}{\ensuremath{\|}}
% environnement proof (amsthm non chargé) utilisé par les modèles
\ifdefined\proof\else\newenvironment{proof}[1][Démonstration]{\par\noindent\textit{#1.}\ }{\qed\par}\fi
\usepackage{lastpage}
\usepackage{hyperref}
\hypersetup{hidelinks}

% ============================================================
% Palette « Pop » OnBuch+ — mêmes hex que la classe P (plus_ui.dart)
% ============================================================
\definecolor{poporange}{HTML}{F59321}
\definecolor{popdark}{HTML}{E07A0C}
\definecolor{popcream}{HTML}{FFF6E8}
\definecolor{popink}{HTML}{1C1714}
\definecolor{poppurple}{HTML}{7A5AE0}
\definecolor{popgreen}{HTML}{2E9E6A}
\definecolor{poppink}{HTML}{E0457B}
\definecolor{popblue}{HTML}{2F7DE1}
\definecolor{popgold}{HTML}{C98A12}
\definecolor{popmuted}{HTML}{8A7F76}
\definecolor{popline}{HTML}{EADFD2}
\definecolor{popgray}{HTML}{8A7F76}
\colorlet{popgrey}{popgray}
% Alias tolérants (le modèle nomme parfois les couleurs différemment)
\colorlet{popwhite}{white}
\colorlet{popblack}{popink}
\colorlet{popred}{poppink}
\colorlet{popyellow}{popgold}
% Teintes claires pour fonds de tableaux / zones
\colorlet{popblueL}{popblue!10!white}
\colorlet{poporangeL}{poporange!14!white}
\colorlet{popgreenL}{popgreen!12!white}
\colorlet{poppurpleL}{poppurple!12!white}
\colorlet{poppinkL}{poppink!12!white}
\colorlet{popgoldL}{popgold!14!white}

\pagecolor{popcream}

% ============================================================
% Typographie
% ============================================================
\defaultfontfeatures{Ligatures=TeX}
\setmainfont{PlusJakartaSans-Medium.ttf}[Path=../../../fonts/,BoldFont=PlusJakartaSans-Bold.ttf]
% Maths en sans-serif (Fira Math) avec chiffres et lettres droites de Plus
% Jakarta, pour que les formules se fondent dans le texte.
\usepackage[math-style=ISO,bold-style=ISO]{unicode-math}
\setmathfont{FiraMath-Regular.otf}[Path=../../../fonts/]
\setmathfont{PlusJakartaSans-Medium.ttf}[Path=../../../fonts/,range={up/num,up/{Latin,latin},\mathup}]
\usepackage{icomma} % virgule décimale sans espace en mode math
% Caractères mathématiques Unicode absents des polices de texte (Plus Jakarta,
% Archivo Black) : rendus par leur équivalent mathématique, en texte comme en
% formule — sinon ils s'affichent en carré vide (ex. « Arithmétique dans ℤ »).
\usepackage{newunicodechar}
\newunicodechar{ℝ}{\ensuremath{\mathbb{R}}}
\newunicodechar{ℤ}{\ensuremath{\mathbb{Z}}}
\newunicodechar{ℂ}{\ensuremath{\mathbb{C}}}
\newunicodechar{ℕ}{\ensuremath{\mathbb{N}}}
\newunicodechar{ℚ}{\ensuremath{\mathbb{Q}}}
\newunicodechar{𝔻}{\ensuremath{\mathbb{D}}}
\newunicodechar{α}{\ensuremath{\alpha}}
\newunicodechar{β}{\ensuremath{\beta}}
\newunicodechar{γ}{\ensuremath{\gamma}}
\newunicodechar{δ}{\ensuremath{\delta}}
\newunicodechar{ε}{\ensuremath{\varepsilon}}
\newunicodechar{θ}{\ensuremath{\theta}}
\newunicodechar{λ}{\ensuremath{\lambda}}
\newunicodechar{μ}{\ensuremath{\mu}}
\newunicodechar{σ}{\ensuremath{\sigma}}
\newunicodechar{τ}{\ensuremath{\tau}}
\newunicodechar{φ}{\ensuremath{\varphi}}
\newunicodechar{ψ}{\ensuremath{\psi}}
\newunicodechar{ω}{\ensuremath{\omega}}
\newunicodechar{Σ}{\ensuremath{\Sigma}}
% majuscules produites par \MakeUppercase dans les titres (α-aminés -> Α)
\newunicodechar{Α}{\ensuremath{\alpha}}
\newunicodechar{Β}{\ensuremath{\beta}}
\newunicodechar{Γ}{\ensuremath{\Gamma}}
\newunicodechar{Δ}{\ensuremath{\Delta}}
\newunicodechar{Ε}{\ensuremath{\varepsilon}}
\newunicodechar{Θ}{\ensuremath{\Theta}}
\newunicodechar{Λ}{\ensuremath{\Lambda}}
\newunicodechar{Μ}{\ensuremath{\mu}}
\newunicodechar{Τ}{\ensuremath{\tau}}
\newunicodechar{Φ}{\ensuremath{\Phi}}
\newunicodechar{Ψ}{\ensuremath{\Psi}}
\newunicodechar{Ω}{\ensuremath{\Omega}}
\newunicodechar{↦}{\ensuremath{\mapsto}}
\newunicodechar{∈}{\ensuremath{\in}}
\newunicodechar{∉}{\ensuremath{\notin}}
\newunicodechar{∧}{\ensuremath{\wedge}}
\newunicodechar{⇔}{\ensuremath{\Leftrightarrow}}
\newunicodechar{⟺}{\ensuremath{\Longleftrightarrow}}
\newunicodechar{⇒}{\ensuremath{\Rightarrow}}
\newunicodechar{⟹}{\ensuremath{\Longrightarrow}}
\newunicodechar{∘}{\ensuremath{\circ}}
\newunicodechar{∪}{\ensuremath{\cup}}
\newunicodechar{∩}{\ensuremath{\cap}}
\newunicodechar{≡}{\ensuremath{\equiv}}
\newunicodechar{⊂}{\ensuremath{\subset}}
\newunicodechar{⊥}{\ensuremath{\perp}}
\newunicodechar{∃}{\ensuremath{\exists}}
\newunicodechar{∀}{\ensuremath{\forall}}
\newunicodechar{∛}{\ensuremath{\sqrt[3]{\,}}}
\newunicodechar{∜}{\ensuremath{\sqrt[4]{\,}}}
\newunicodechar{⃗}{\ensuremath{{}^{\to}}}
\newunicodechar{ⁿ}{\ifmmode^{n}\else\textsuperscript{\ensuremath{n}}\fi}
\newunicodechar{ˣ}{\ifmmode^{x}\else\textsuperscript{\ensuremath{x}}\fi}
\newunicodechar{ᵇ}{\ifmmode^{b}\else\textsuperscript{\ensuremath{b}}\fi}
\newunicodechar{ʳ}{\ifmmode^{r}\else\textsuperscript{\ensuremath{r}}\fi}
\newunicodechar{ᵘ}{\ifmmode^{u}\else\textsuperscript{\ensuremath{u}}\fi}
\newunicodechar{ʸ}{\ifmmode^{y}\else\textsuperscript{\ensuremath{y}}\fi}
\newunicodechar{ᵃ}{\ifmmode^{a}\else\textsuperscript{\ensuremath{a}}\fi}
\newunicodechar{ᵉ}{\ifmmode^{e}\else\textsuperscript{\ensuremath{e}}\fi}
\newunicodechar{ᵏ}{\ifmmode^{k}\else\textsuperscript{\ensuremath{k}}\fi}
\newunicodechar{ᵖ}{\ifmmode^{p}\else\textsuperscript{\ensuremath{p}}\fi}
\newunicodechar{ᴺ}{\ifmmode^{N}\else\textsuperscript{\ensuremath{N}}\fi}
\newunicodechar{ᶜ}{\ifmmode^{c}\else\textsuperscript{\ensuremath{c}}\fi}
\newunicodechar{ʰ}{\ifmmode^{h}\else\textsuperscript{\ensuremath{h}}\fi}
\newunicodechar{ᵠ}{\ifmmode^{q}\else\textsuperscript{\ensuremath{q}}\fi}
\newunicodechar{ₙ}{\ifmmode_{n}\else\textsubscript{\ensuremath{n}}\fi}
\newunicodechar{ₐ}{\ifmmode_{a}\else\textsubscript{\ensuremath{a}}\fi}
\newunicodechar{ᵢ}{\ifmmode_{i}\else\textsubscript{\ensuremath{i}}\fi}
\newunicodechar{ᵣ}{\ifmmode_{r}\else\textsubscript{\ensuremath{r}}\fi}
\newunicodechar{ₖ}{\ifmmode_{k}\else\textsubscript{\ensuremath{k}}\fi}
\newunicodechar{ᵦ}{\ifmmode_{\beta}\else\textsubscript{\ensuremath{\beta}}\fi}
\newunicodechar{ₓ}{\ifmmode_{x}\else\textsubscript{\ensuremath{x}}\fi}
\newfontfamily\popdisplay{ArchivoBlack-Regular.ttf}[Path=../../../fonts/]
\newfontfamily\poplogo{SpaceGrotesk-Bold.ttf}[Path=../../../fonts/]
\newfontfamily\popsemi{PlusJakartaSans-SemiBold.ttf}[Path=../../../fonts/]
\newfontfamily\popxbold{PlusJakartaSans-ExtraBold.ttf}[Path=../../../fonts/]
\newfontfamily\popxboldls{PlusJakartaSans-ExtraBold.ttf}[Path=../../../fonts/,LetterSpace=3.0]

\setlist[enumerate]{leftmargin=1.7em,itemsep=5pt,topsep=4pt}
\setlist[itemize]{leftmargin=1.7em,itemsep=4pt,topsep=4pt,label={\color{poporange}\textbullet}}
\setlength{\parindent}{0pt}
\setlength{\parskip}{5pt}
\renewcommand{\arraystretch}{1.25}

% pgfplots : style Pop par défaut
\pgfplotsset{
  every axis/.append style={axis line style={popink,line width=1pt},tick style={popink},
    grid style={popline,line width=0.5pt},label style={font=\small},tick label style={font=\footnotesize}},
  popcurve/.style={poporange,line width=1.8pt,no markers,smooth},
}
% Même style disponible dans un \draw TikZ simple (hors pgfplots)
\tikzset{popcurve/.style={poporange,line width=1.8pt}}
\tikzset{popgrid/.style={popline,very thin}}
\colorlet{popcurve}{poporange} % le nom du style est parfois utilisé comme couleur

% ============================================================
% Pill / squiggle
% ============================================================
\newcommand{\poppill}[3]{%
  \tikz[baseline=-0.55ex]\node[draw=popink,line width=1.2pt,rounded corners=2.6pt,%
    fill=#2,inner xsep=6.5pt,inner ysep=3pt]{%
    \popxboldls\fontsize{7}{7}\selectfont\color{#3}\MakeUppercase{#1}};%
}
\newcommand{\popsquiggle}[1]{%
  \tikz[baseline=-0.4ex]{
    \draw[#1,line width=2.6pt,line cap=round]
      plot [smooth] coordinates {(0,0.09) (0.34,0.01) (0.68,0.21) (1.02,0.01) (1.36,0.10)};
  }%
}
% Mot-clé surligné (marqueur orange sous le texte)
\newcommand{\cle}[1]{{\popxbold\color{popdark}#1}}

% ============================================================
% En-tête / pied
% ============================================================
\pagestyle{fancy}
\fancyhf{}
\renewcommand{\headrulewidth}{0pt}
\renewcommand{\footrulewidth}{0pt}
\lhead{\raisebox{-0.15ex}{\poplogo\fontsize{13}{13}\selectfont\color{popink}OnBuch\color{poporange}+}}
\rhead{\poppill{\DOCMATIERE\ \textbullet\ \DOCNIVEAU\ \textbullet\ Leçon \DOCLECON}{white}{popink}}
\lfoot{\popxbold\fontsize{7.6}{7.6}\selectfont\color{popmuted}\MakeUppercase{Cours signé OnBuch+}\ \textbullet\ {\popsemi\DOCTITRE}}
\rfoot{\tikz[baseline=-0.6ex]\node[rounded corners=8.5pt,fill=popink,minimum height=17pt,minimum width=17pt,inner xsep=5pt,inner sep=0]{\popxbold\fontsize{8}{8}\selectfont\color{popcream}\thepage\,/\,\pageref*{LastPage}};}

% ============================================================
% Titres
% ============================================================
\newcommand{\chaptitre}[2]{%
  \begin{tcolorbox}[enhanced,colback=poporange,colframe=popink,boxrule=2.6pt,arc=11pt,
      drop shadow=popink,left=15pt,right=15pt,top=12pt,bottom=11pt,before skip=3pt,after skip=18pt]
    \poppill{#2}{popink}{popcream}\par\vspace{9pt}
    {\raggedright\hyphenpenalty=10000\exhyphenpenalty=10000\popdisplay\fontsize{22}{24}\selectfont\color{popcream}\MakeUppercase{#1}\par}
  \end{tcolorbox}
}
% Section numérotée : pastille orange avec le numéro + titre Archivo
\newcounter{popsec}
\newcommand{\coursec}[1]{%
  \stepcounter{popsec}\setcounter{popsub}{0}%
  \par\vspace{16pt}\noindent
  \begin{minipage}{\linewidth}
  \tikz[baseline=-0.7ex]\node[draw=popink,line width=1.6pt,fill=poporange,rounded corners=4pt,minimum size=22pt,inner sep=0]{\popdisplay\fontsize{12}{12}\selectfont\color{popcream}\thepopsec};%
  \hspace{8pt}{\popdisplay\fontsize{15}{17}\selectfont\color{popink}\MakeUppercase{#1}}\par
  \vspace{4pt}\noindent{\color{poporange}\rule{36pt}{3.6pt}}
  \end{minipage}\par\nopagebreak\vspace{8pt}
}
\newcounter{popsub}
\newcommand{\courssub}[1]{%
  \stepcounter{popsub}%
  \par\vspace{9pt}\noindent{\popxbold\fontsize{11.5}{13}\selectfont\color{popink}{\color{poporange}\thepopsec.\thepopsub}\hspace{6pt}#1}\par\nopagebreak\vspace{4pt}
}

% ============================================================
% Boîtes pédagogiques — même recette : fond blanc, bordure épaisse colorée,
% ombre dure encre, badge plein. Titre optionnel [#1].
% ============================================================
\tcbset{popbase/.style={breakable,enhanced,colback=white,boxrule=2.3pt,arc=9pt,
  shadow={3pt}{-3pt}{0pt}{popink},left=14pt,right=14pt,top=11pt,bottom=11pt,before skip=11pt,after skip=15pt,
  fontupper=\popsemi\fontsize{10.6}{14.5}\selectfont\color{popink},
  fontlower=\popsemi\fontsize{10.6}{14.5}\selectfont\color{popink}}}
% Coche dessinée (le glyphe ✓ n'existe pas dans les polices du document)
\newcommand{\popcheck}[1]{\tikz[baseline=-0.5ex]\draw[#1,line width=1.6pt,line cap=round,line join=round](-0.13,0.0)--(-0.04,-0.09)--(0.13,0.1);}
\providecommand{\coche}{\popcheck{popgreen}}
\providecommand{\resultat}[1]{\par\smallskip\noindent\fcolorbox{popgreen}{popgreenL}{\parbox{\dimexpr\linewidth-2\fboxsep-2\fboxrule\relax}{\textbf{Résultat.} #1}}\par\smallskip}
\newcommand{\popboxhead}[3]{\poppill{#1}{#2}{white}\ifx\relax#3\relax\else\hspace{6pt}{\popxbold\fontsize{10.6}{12}\selectfont\color{popink}#3}\fi\par\vspace{7pt}}

\newtcolorbox{aretenir}[1][]{popbase,colframe=popgreen,before upper={\popboxhead{À retenir}{popgreen}{#1}}}
\newtcolorbox{exemplebox}[1][]{popbase,colframe=poporange,before upper={\popboxhead{Exemple}{poporange}{#1}}}
\newtcolorbox{definition}[1][]{popbase,colframe=popblue,before upper={\popboxhead{Définition}{popblue}{#1}}}
\newtcolorbox{propriete}[1][]{popbase,colframe=poppurple,before upper={\popboxhead{Propriété}{poppurple}{#1}}}
\newtcolorbox{methode}[1][]{popbase,colframe=popgold,before upper={\popboxhead{Méthode}{popgold}{#1}}}
\newtcolorbox{attention}[1][]{popbase,colframe=poppink,before upper={\popboxhead{Attention}{poppink}{#1}}}
\newtcolorbox{experience}[1][]{popbase,colframe=popblue,colback=popblueL,before upper={\popboxhead{Expérience}{popblue}{#1}}}
% « Le savais-tu ? » : carte violette pleine (contexte, culture, Cameroun)
\newtcolorbox{savaistu}[1][]{popbase,colback=poppurple,colframe=popink,
  fontupper=\popsemi\fontsize{10.4}{14.2}\selectfont\color{white},
  before upper={\poppill{Le savais-tu\,?}{white}{poppurple}\ifx\relax#1\relax\else\hspace{6pt}{\popxbold\color{white}#1}\fi\par\vspace{7pt}}}
% Exercice résolu : énoncé en haut, \tcblower puis la solution sur fond crème
\newcounter{popexo}\newcounter{popexr}
\newtcolorbox{exoresolu}[1][]{popbase,colframe=poporange,colbacklower=popcream,
  segmentation style={popink,line width=1.2pt,dashed},
  before upper={\stepcounter{popexr}\popboxhead{Exercice résolu \thepopexr}{poporange}{#1}},
  before lower={\poppill{Solution}{popink}{popcream}\par\vspace{6pt}}}
% Exercice (non résolu en place) — la correction est dans la partie Corrigés
\newtcolorbox{exercice}[1][]{popbase,colframe=popink,
  before upper={\stepcounter{popexo}\popboxhead{Exercice \thepopexo}{popink}{#1}}}
% Corrigé
\newtcolorbox{corrige}[1][]{popbase,colframe=popgreen,colback=popgreenL,
  before upper={\popboxhead{Corrigé}{popgreen}{#1}}}
% Objectifs / prérequis / bilan
\newtcolorbox{objectifs}{popbase,colframe=popink,colback=poporangeL,
  before upper={\popboxhead{Objectifs de la leçon}{popink}{}}}
\newtcolorbox{prerequis}{popbase,colframe=popink,colback=popblueL,
  before upper={\popboxhead{Prérequis}{popblue}{}}}
\newtcolorbox{bilan}{popbase,colframe=popink,colback=white,boxrule=2.8pt,
  before upper={\popboxhead{Fiche bilan}{poporange}{L'essentiel à connaître pour l'examen}}}
% Figure encadrée avec légende
\newtcolorbox{popfigure}[1][]{enhanced,breakable=false,colback=white,colframe=popline,boxrule=1.4pt,arc=8pt,
  left=10pt,right=10pt,top=10pt,bottom=8pt,before skip=10pt,after skip=12pt,halign=center,
  lower separated=false,halign lower=center,
  fontlower=\popsemi\fontsize{9}{11.5}\selectfont\color{popmuted},
  #1}
\newcommand{\legende}[1]{\par\vspace{6pt}{\popsemi\fontsize{9}{11.5}\selectfont\color{popmuted}#1}}

% ============================================================
% Signature OnBuch+
% ============================================================
\newcommand{\poplogoinline}[1]{{\poplogo\fontsize{#1}{#1}\selectfont\color{popink}OnBuch\color{poporange}+}}
\newcommand{\signatureonbuch}{%
  \begin{tcolorbox}[enhanced,colback=popink,colframe=popink,boxrule=2.2pt,arc=10pt,
      drop shadow=poporange,left=14pt,right=14pt,top=10pt,bottom=10pt,before skip=0pt,after skip=0pt]
    \tikz[baseline=-0.7ex]\node[circle,fill=popgreen,draw=popcream,line width=1.3pt,minimum size=22pt,inner sep=0]{\popcheck{white}};%
    \hspace{9pt}\begin{minipage}[c]{0.62\linewidth}
      {\popxbold\fontsize{12}{13}\selectfont\color{popcream}Signé\ \textbullet\ {\poplogo OnBuch\color{poporange}+}}\par\vspace{2pt}
      {\raggedright\popsemi\fontsize{8}{10.5}\selectfont\color{popline}\MakeUppercase{Équipe pédagogique OnBuch+}\\\MakeUppercase{Conforme au programme officiel MINESEC}\par}
    \end{minipage}\hfill
    {\poplogo\fontsize{22}{22}\selectfont\color{popcream}OnBuch\color{poporange}+}
  \end{tcolorbox}%
}

% ============================================================
% Page de garde
% #1 titre · #2 matière · #3 niveau · #4 module · #5 n° leçon · #6 durée ·
% #7 description / objectifs (texte ou itemize)
% ============================================================
\newcommand{\pagegarde}[7]{%
  \thispagestyle{empty}
  \newgeometry{a4paper,margin=1.7cm,top=1.6cm,bottom=1.4cm}%
  \begin{tikzpicture}[remember picture,overlay]
    \fill[poporange!18] ([xshift=-3.2cm,yshift=-3.6cm]current page.north east) circle (4.6cm);
    \fill[poppurple!14] ([xshift=2.4cm,yshift=7.6cm]current page.south west) circle (3.8cm);
  \end{tikzpicture}%
  {\poplogo\fontsize{30}{30}\selectfont\color{popink}OnBuch\color{poporange}+}\hfill
  \raisebox{0.6ex}{\poppill{Cours signé \textbullet\ Premium}{popink}{popcream}}\par
  \vspace{1pt}\popsquiggle{poppurple}\par
  \vspace{8pt}
  \begin{tcolorbox}[enhanced,colback=poporange,colframe=popink,boxrule=2.8pt,arc=13pt,
      drop shadow=popink,left=18pt,right=18pt,top=12pt,bottom=13pt,after skip=10pt]
    \poppill{#2 \textbullet\ #3}{popink}{popcream}\hspace{5pt}\poppill{#4}{white}{popink}\par\vspace{8pt}
    {\popdisplay\fontsize{14}{14}\selectfont\color{popink}LEÇON #5}\par\vspace{4pt}
    {\raggedright\hyphenpenalty=10000\exhyphenpenalty=10000\popdisplay\fontsize{25}{27}\selectfont\color{popcream}\MakeUppercase{#1}\par}
    \vspace{8pt}
    {\popxbold\fontsize{9.5}{9.5}\selectfont\color{popink}\MakeUppercase{Durée conseillée : #6}}
  \end{tcolorbox}
  \begin{tcolorbox}[enhanced,colback=white,colframe=popink,boxrule=2.2pt,arc=10pt,
      drop shadow=popink,left=16pt,right=16pt,top=10pt,bottom=10pt,after skip=8pt]
    \poppill{Dans cette leçon}{poppurple}{white}\par\vspace{5pt}
    \popsemi\fontsize{9.3}{12.6}\selectfont\color{popink}#7
  \end{tcolorbox}
  \noindent\begin{minipage}[t]{0.487\linewidth}
    \begin{tcolorbox}[enhanced,colback=popgreen,colframe=popink,boxrule=2.2pt,arc=10pt,
        drop shadow=popink,left=11pt,right=11pt,top=10pt,bottom=10pt,height=3.1cm,valign=center]
      \tcbox[colback=white,colframe=popink,boxrule=1.3pt,arc=5pt,boxsep=0pt,nobeforeafter,
        left=3pt,right=3pt,top=3pt,bottom=3pt,valign=center]{\qrcode[height=1.9cm]{https://play.google.com/store/apps/details?id=com.luvvix.onbuch&hl=fr}}%
      \hspace{8pt}\begin{minipage}[c]{0.5\linewidth}
        {\popdisplay\fontsize{11}{12.5}\selectfont\color{white}\MakeUppercase{Télécharge OnBuch}}\par\vspace{4pt}
        {\raggedright\popsemi\fontsize{8.4}{11}\selectfont\color{white}Gratuit \textbullet\ Google Play\\Tuteur IA, annales, concours, résultats\par}
      \end{minipage}
    \end{tcolorbox}
  \end{minipage}\hfill
  \begin{minipage}[t]{0.487\linewidth}
    \begin{tcolorbox}[enhanced,colback=poppurple,colframe=popink,boxrule=2.2pt,arc=10pt,
        drop shadow=popink,left=11pt,right=11pt,top=10pt,bottom=10pt,height=3.1cm,valign=center]
      \tikz[baseline=-0.7ex]\node[circle,fill=white,draw=popink,line width=1.3pt,minimum size=1.6cm,inner sep=0]{\popdisplay\fontsize{24}{24}\selectfont\color{poppurple}+};%
      \hspace{8pt}\begin{minipage}[c]{0.6\linewidth}
        {\popdisplay\fontsize{11}{12.5}\selectfont\color{white}\MakeUppercase{Passe à OnBuch+}}\par\vspace{4pt}
        {\raggedright\popsemi\fontsize{8.4}{11}\selectfont\color{white}Tous les cours signés, Tuteur IA illimité, fiches PDF — onglet OnBuch+ de l'app\par}
      \end{minipage}
    \end{tcolorbox}
  \end{minipage}
  \vspace{8pt}
  \popcontenu
  \vfill
  \signatureonbuch
  \restoregeometry
  \newpage
}

% Rangée « Ce cours contient » (page de garde)
\newcommand{\poptile}[3]{% #1 couleur #2 gros texte #3 légende
  \begin{tcolorbox}[enhanced,colback=#1,colframe=popink,boxrule=2pt,arc=9pt,shadow={2.5pt}{-2.5pt}{0pt}{popink},
      left=6pt,right=6pt,top=9pt,bottom=9pt,halign=center,valign=center,height=1.75cm,width=\linewidth,nobeforeafter]
    {\popdisplay\fontsize{12.5}{13}\selectfont\color{white}\MakeUppercase{#2}}\par\vspace{3pt}
    {\centering\popsemi\fontsize{7.8}{9.5}\selectfont\color{white}#3\par}
  \end{tcolorbox}}
\newcommand{\popcontenu}{%
  \par\noindent{\popxboldls\fontsize{7.5}{7.5}\selectfont\color{popmuted}CE COURS SIGNÉ CONTIENT}\par\vspace{7pt}
  \noindent\begin{minipage}[t]{0.235\linewidth}\poptile{popblue}{Cours}{complet, illustré, pas à pas}\end{minipage}\hfill
  \begin{minipage}[t]{0.235\linewidth}\poptile{poporange}{Méthodes}{et exercices résolus}\end{minipage}\hfill
  \begin{minipage}[t]{0.235\linewidth}\poptile{poppink}{Exercices}{type examen + corrigés}\end{minipage}\hfill
  \begin{minipage}[t]{0.235\linewidth}\poptile{popgreen}{Fiche bilan}{l'essentiel à retenir}\end{minipage}\par}

% Sommaire
\newtcolorbox{sommaire}{enhanced,colback=white,colframe=popink,boxrule=2.2pt,arc=10pt,
  drop shadow=popink,left=18pt,right=18pt,top=13pt,bottom=13pt,before skip=6pt,after skip=20pt,
  before upper={\poppill{Sommaire}{poppurple}{white}\par\vspace{9pt}\popsemi\fontsize{10.6}{14.5}\selectfont\color{popink}}}

% Signature de fin de cours — poussée tout en bas de la dernière page
\newcommand{\findecours}{%
  \par\vfill\nopagebreak
  \begin{center}\popsquiggle{poporange}\end{center}\vspace{4pt}
  \signatureonbuch
}
\AtBeginDocument{\def\llbracket{\mathopen{[\mkern-3.5mu[}}\def\rrbracket{\mathclose{]\mkern-3.5mu]}}} % intervalles d'entiers (glyphe ⟦ absent de la police)
% Glyphes absents de Fira Math : redéfinis à partir de glyphes disponibles
\AtBeginDocument{%
  \def\setminus{\mathbin{\backslash}}\let\smallsetminus\setminus
  \def\vdots{\mathord{\vbox{\baselineskip3.2pt\lineskiplimit0pt\kern4pt\hbox{.}\hbox{.}\hbox{.}}}}%
  \def\ddots{\mathinner{\mkern1mu\raise6.5pt\hbox{.}\mkern2mu\raise3.6pt\hbox{.}\mkern2mu\raise0.7pt\hbox{.}\mkern1mu}}%
  \def\triangle{\mathord{\scalebox{0.9}{$\symup{\Delta}$}}}%
  \def\nabla{\mathord{\raisebox{\depth}{\scalebox{1}[-1]{$\symup{\Delta}$}}}}%
  \def\checkmark{\popcheck{popdark}}%
  \def\star{\mathbin{\ast}}\def\bigstar{\mathord{\ast}}%
  \def\diamond{\mathbin{\scalebox{0.6}{$\Diamond$}}}%
  \def\bigcup{\mathop{\mathchoice{\raisebox{-0.3em}{\scalebox{1.7}{$\cup$}}}{\scalebox{1.25}{$\cup$}}{\cup}{\cup}}\displaylimits}%
  \def\bigcap{\mathop{\mathchoice{\raisebox{-0.3em}{\scalebox{1.7}{$\cap$}}}{\scalebox{1.25}{$\cap$}}{\cap}{\cap}}\displaylimits}%
  \def\lesseqqgtr{\mathrel{\lessgtr}}\def\bigoplus{\mathop{\oplus}\displaylimits}%
}
\providecommand{\ssi}{si et seulement si} % abréviation fréquente
\AtBeginDocument{\providecommand{\wideparen}{\overparen}} % arc de cercle

%% FIGFIX-BEGIN v5 — correctifs globaux des figures (docs/onbuch-generation/tools/figfix)
\usepackage{adjustbox}\usepackage{etoolbox}\tcbuselibrary{raster}
% toute figure TikZ/pgfplots plus large que la ligne est réduite à la largeur disponible
\BeforeBeginEnvironment{tikzpicture}{\begin{adjustbox}{max width=\linewidth}}
\AfterEndEnvironment{tikzpicture}{\end{adjustbox}}
% légendes pgfplots lisibles par-dessus les courbes
\pgfplotsset{every axis/.append style={legend style={fill=white,fill opacity=0.92,text opacity=1,draw=black!15,inner sep=3pt}}}
% étiquettes d'arêtes (syntaxe edge["..."]) avec fond blanc
\tikzset{every edge quotes/.append style={fill=white,inner sep=1.2pt}}
%% FIGFIX-END
\begin{document}

\pagegarde{Maîtrise des comportements interindividuels}{SVT}{1ère A}{Module 1 : Le Monde Vivant}{NA03}{3 h}{Cette leçon étudie les fondements physiologiques (systèmes nerveux et hormonal) des comportements et des relations interindividuelles : perception, émotion et communication. Elle conduit l'élève à expliquer des comportements sociaux, à analyser des situations de vie courante et à développer le respect de la diversité des comportements humains.
\vspace{4pt}
\begin{multicols}{2}\small
\begin{itemize}
\item À la fin de cette leçon, je sais définir les notions de comportement, de comportement interindividuel, de perception, d'émotion et de communication.
\item À la fin de cette leçon, je sais décrire le rôle des récepteurs sensoriels et des aires cérébrales dans la perception d'autrui.
\item À la fin de cette leçon, je sais expliquer les mécanismes nerveux (réflexe, système limbique, cortex) d'une émotion et d'une réaction comportementale.
\item À la fin de cette leçon, je sais expliquer le rôle des hormones (adrénaline, ocytocine, testostérone, cortisol) dans les comportements sociaux.
\item À la fin de cette leçon, je sais identifier les différentes formes de communication (verbale, non verbale) et leur support physiologique.
\item À la fin de cette leçon, je sais analyser une situation de vie courante (conflit, solidarité, stress) en mobilisant ses bases physiologiques.
\end{itemize}\end{multicols}}

\begin{sommaire}
\begin{multicols}{2}
\begin{enumerate}
\item Comportements et relations interindividuels : généralités
\item La perception d'autrui : support nerveux
\item Les émotions : bases nerveuses et hormonales
\item La communication interindividuelle
\item Régulation hormonale des comportements sociaux
\item Diversité des comportements humains et respect d'autrui
\item Activité d'intégration
\item Méthodes et exercices résolus
\item Exercices
\item Corrigés des exercices
\item Fiche bilan
\end{enumerate}
\end{multicols}
\end{sommaire}

\begin{objectifs}
\begin{itemize}
\item À la fin de cette leçon, je sais définir les notions de comportement, de comportement interindividuel, de perception, d'émotion et de communication.
\item À la fin de cette leçon, je sais décrire le rôle des récepteurs sensoriels et des aires cérébrales dans la perception d'autrui.
\item À la fin de cette leçon, je sais expliquer les mécanismes nerveux (réflexe, système limbique, cortex) d'une émotion et d'une réaction comportementale.
\item À la fin de cette leçon, je sais expliquer le rôle des hormones (adrénaline, ocytocine, testostérone, cortisol) dans les comportements sociaux.
\item À la fin de cette leçon, je sais identifier les différentes formes de communication (verbale, non verbale) et leur support physiologique.
\item À la fin de cette leçon, je sais analyser une situation de vie courante (conflit, solidarité, stress) en mobilisant ses bases physiologiques.
\item À la fin de cette leçon, je sais citer les facteurs biologiques et psychosociaux expliquant la diversité des comportements humains.
\item À la fin de cette leçon, je sais adopter et justifier une attitude de respect face à la diversité des comportements.
\end{itemize}
\end{objectifs}

\begin{prerequis}
\begin{itemize}
\item Organisation générale du système nerveux (cerveau, moelle épinière, nerfs) vue en classe de Seconde
\item Notion de réflexe et d'arc réflexe (message nerveux, neurone)
\item Notion d'hormone et de glande endocrine (Seconde)
\item Les sens et les récepteurs sensoriels (œil, oreille, peau) vus en Seconde
\item Notion de milieu intérieur et de régulation de l'organisme
\end{itemize}
\end{prerequis}

\newpage

\coursec{Comportements et relations interindividuels : généralités}

\textbf{Situation d'accroche.} C'est samedi après-midi au quartier. Le match de football oppose deux équipes de jeunes. Soudain, un joueur insulté par un adversaire sent son cœur s'accélérer, ses poings se serrer, son visage chauffer. Avant même d'avoir réfléchi, il répond violemment : la bagarre éclate et gagne tout le terrain. Quelques minutes plus tard, calmé, il regrette son geste : « Je ne sais pas ce qui m'a pris... »

Cette scène, que chacun a déjà vécue ou observée, pose des questions fondamentales : pourquoi notre corps réagit-il ainsi \emph{avant même} que nous ayons réfléchi ? Quels mécanismes biologiques commandent nos émotions, notre façon de percevoir autrui et de communiquer ? Pourquoi le même événement (une insulte) ne provoque-t-il pas la même réaction chez tout le monde ?

\textbf{Problématique.} Qu'est-ce qu'un comportement interindividuel et sur quels mécanismes physiologiques repose-t-il ? Comprendre ces mécanismes, c'est faire le premier pas vers la \cle{maîtrise de soi} et le respect d'autrui, conditions de la cohésion sociale dans nos quartiers, nos écoles et nos villages.

\courssub{Qu'est-ce qu'un comportement ?}

\textbf{Observation préalable.} Place ta main au-dessus d'une flamme : tu la retires instantanément. Entends ton nom appelé dans la cour : tu te retournes. Vois un serpent sur le sentier : tu recules. Dans chaque cas, ton organisme a \emph{réagi} à quelque chose qui vient du milieu. C'est exactement ce qu'on appelle un comportement.

\begin{definition}[Comportement]
Un \cle{comportement} est l'ensemble des \textbf{réactions observables} d'un individu face à une \textbf{stimulation} (ou \cle{stimulus}) provenant du milieu extérieur ou de son propre organisme (milieu intérieur). Une stimulation est toute modification du milieu capable d'être détectée par l'organisme : lumière, son, odeur, contact, mais aussi faim, soif ou douleur interne.
\end{definition}

Trois remarques importantes pour bien comprendre cette définition :

\begin{itemize}
\item Le comportement est \textbf{observable} : on peut le voir, le mesurer, le décrire (un geste, une parole, une fuite, une mimique). Une pensée secrète n'est pas un comportement tant qu'elle ne se traduit par rien de visible.
\item Le comportement est une \textbf{réponse} : il fait toujours suite à une stimulation, même si celle-ci est parfois difficile à identifier (par exemple une baisse de sucre dans le sang qui déclenche la recherche de nourriture).
\item Le comportement met en jeu \textbf{tout l'organisme} : les organes des sens qui reçoivent la stimulation, le système nerveux qui l'analyse, les muscles et les glandes qui exécutent la réponse.
\end{itemize}

\begin{exemplebox}[Des comportements de la vie quotidienne]
Retirer sa main d'une marmite brûlante, cligner des yeux quand un grain de sable approche, saluer un ancien dans la rue, courir vers la mangeoire comme les poules du village quand on jette le maïs : tous ces actes sont des comportements, car ce sont des réactions observables à des stimulations.
\end{exemplebox}

\courssub{Le comportement interindividuel}

\textbf{Observation préalable.} Compare deux situations : un élève marche seul dans la cour, puis le même élève aperçoit son meilleur ami. Immédiatement, son allure change : il sourit, accélère le pas, lève la main. La \emph{présence d'un autre individu} a modifié son comportement. C'est ce cas particulier, essentiel pour la vie en société, qui nous intéresse dans cette leçon.

\begin{definition}[Comportement interindividuel]
Un \cle{comportement interindividuel} est un comportement \textbf{déclenché ou modifié par la présence ou l'action d'un autre individu} de la même espèce. La stimulation n'est plus un objet ou un phénomène physique, mais un \textbf{être vivant semblable} : son regard, sa voix, ses gestes, sa simple présence.
\end{definition}

Les comportements interindividuels sont extrêmement variés. On peut les classer en grandes catégories selon leur effet sur la relation :

\begin{itemize}
\item L'\cle{attraction} : comportements qui rapprochent les individus (sourire, salutation, invitation à jouer, parade amoureuse) ;
\item L'\cle{agression} : comportements d'attaque ou de menace (insulte, coup, regard hostile, intimidation) ;
\item La \cle{coopération} : comportements d'entraide en vue d'un but commun (travail collectif du champ, tontine, construction d'une case) ;
\item L'\cle{évitement} : comportements de fuite ou d'éloignement (changer de trottoir, se taire, quitter le groupe).
\end{itemize}

\begin{exemplebox}[Des comportements interindividuels bien camerounais]
\begin{itemize}
\item \textbf{Les salutations codifiées} : dans de nombreuses communautés, le cadet salue l'aîné avec des gestes précis (inclination, poignée de main à deux mains, paroles rituelles). Ce comportement, déclenché par la rencontre de l'aîné, exprime le respect et maintient l'harmonie sociale.
\item \textbf{L'entraide villageoise} : lors d'une \emph{njangi} (tontine) ou d'un travail collectif de défrichage, chacun coopère en attendant d'être aidé à son tour. La présence des autres stimule l'effort et la solidarité.
\item \textbf{Les réactions lors d'un match} : but encaissé, faute sifflée, insulte d'un adversaire... les joueurs et les supporters réagissent instantanément les uns aux autres, parfois jusqu'au débordement.
\item \textbf{Les conflits communautaires} : disputes de bornage de champs, rivalités de quartier : des comportements d'agression et d'évitement peuvent s'enchaîner et s'amplifier entre groupes, montrant combien la maîtrise des comportements est un enjeu social majeur.
\end{itemize}
\end{exemplebox}

\courssub{Les trois composantes d'un comportement interindividuel}

\textbf{Démarche d'investigation.} Reprenons la scène du match et décomposons-la comme le ferait un scientifique.

\begin{itemize}
\item \textbf{Observation} : le joueur insulté a d'abord \emph{entendu et compris} l'insulte ; puis il a \emph{ressenti} quelque chose (cœur qui s'accélère, chaleur au visage) ; enfin il a \emph{agi} (riposte violente).
\item \textbf{Hypothèse} : tout comportement interindividuel se déroule en trois temps : recevoir l'information sur l'autre, ressentir une émotion, produire une réponse.
\item \textbf{Vérification} : cette séquence se retrouve dans tous les exemples analysés (salutation, coopération, fuite). De plus, on peut bloquer chaque étape séparément : un joueur sourd ne perçoit pas l'insulte et ne réagit pas ; un joueur très calme perçoit l'insulte mais ne ressent pas de colère intense et ne riposte pas.
\item \textbf{Conclusion} : le comportement interindividuel repose sur trois composantes enchaînées : la \textbf{perception}, l'\textbf{émotion} et la \textbf{communication}.
\end{itemize}

\begin{definition}[Les trois piliers du comportement interindividuel]
\begin{itemize}
\item La \cle{perception} : ensemble des mécanismes par lesquels les organes des sens (yeux, oreilles, peau...) et le cerveau \textbf{reçoivent et interprètent} les informations venant d'autrui (visage, voix, gestes, paroles).
\item L'\cle{émotion} : \textbf{réaction affective} intense et brève (joie, colère, peur, tristesse), accompagnée de modifications corporelles (accélération cardiaque, rougeur, sueur), qui prépare l'organisme à agir.
\item La \cle{communication} : ensemble des signaux \textbf{émis vers autrui} (paroles, gestes, mimiques, attitudes) qui constituent la réponse comportementale et deviennent à leur tour des stimulations pour l'autre.
\end{itemize}
\end{definition}

Ces trois composantes forment une \textbf{chaîne fonctionnelle} : la réponse émise par un individu devient le stimulus perçu par l'autre, ce qui explique que les comportements s'enchaînent et parfois s'amplifient (une insulte en appelle une autre).

\begin{popfigure}
\begin{center}
\begin{tikzpicture}[
  node distance=1.5cm and 1.2cm,
  box/.style={draw=popblue, fill=popblueL, rounded corners=3pt, minimum height=1.1cm, text width=2.8cm, align=center, font=\small\bfseries, text=popdark},
  sub/.style={font=\scriptsize, text=popmuted, align=center, text width=3.0cm},
  fleche/.style={-{Stealth[length=3mm]}, very thick, poporange},
  loop/.style={-{Stealth[length=3mm]}, thick, dashed, poppurple, bend left=45}
]
\node[box] (stim) {Stimulus social};
\node[sub, below=2mm of stim] {paroles, gestes, regard d'autrui};
\node[box, right=of stim] (perc) {Perception};
\node[sub, below=2mm of perc] {organes des sens + cerveau interprètent};
\node[box, right=of perc] (emo) {Émotion};
\node[sub, below=2mm of emo] {colère, joie, peur + réactions corporelles};
\node[box, right=of emo] (rep) {Réponse comportementale};
\node[sub, below=2mm of rep] {communication : mots, gestes, mimiques};
\draw[fleche] (stim) -- (perc);
\draw[fleche] (perc) -- (emo);
\draw[fleche] (emo) -- (rep);
\draw[loop] (rep.north) to node[midway, above, font=\scriptsize\itshape, text=poppurple, align=center] {la réponse devient\\stimulus pour l'autre} (stim.north);
\end{tikzpicture}
\end{center}
\legende{Chaîne fonctionnelle du comportement interindividuel : le stimulus social est perçu, provoque une émotion, qui débouche sur une réponse. La flèche violette en pointillés montre que la réponse d'un individu devient à son tour un stimulus pour l'autre : c'est le principe de l'interaction.}
\end{popfigure}

\begin{exemplebox}[Analyse d'une situation : l'élève bousculé dans la cour]
À la récréation, un élève est violemment bousculé par un camarade qui court. Décomposons sa réaction :
\begin{itemize}
\item \textbf{Perception} : ses yeux voient le camarade arriver, sa peau ressent le choc, son cerveau interprète le geste (« il m'a bousculé exprès ! ») ;
\item \textbf{Émotion} : une colère soudaine monte ; son cœur bat plus vite, ses muscles se tendent, son visage rougit ;
\item \textbf{Communication (réponse)} : il crie « Hé ! Fais attention ! » et pousse le camarade à son tour.
\end{itemize}
Remarque décisive : si le même élève avait perçu le geste comme \emph{accidentel} (le camarade trébuche et s'excuse), l'émotion ressentie aurait été faible et la réponse toute différente (« ce n'est rien »). Cela montre que c'est l'\textbf{interprétation} faite par le cerveau, autant que le stimulus lui-même, qui oriente le comportement.
\end{exemplebox}

\begin{attention}[Le stimulus n'est pas la seule cause de la réponse]
Deux élèves recevant la même insulte peuvent réagir très différemment (rire, indifférence, ou bagarre). Le comportement dépend de la perception et de l'interprétation, donc de l'expérience, de la culture et de l'état du moment (fatigue, stress). Ne réduis jamais un comportement à un simple « réflexe mécanique » dans tes rédactions : la chaîne perception $\rightarrow$ émotion $\rightarrow$ réponse comporte toujours une part d'interprétation.
\end{attention}

\courssub{Comportements innés et comportements acquis}

\textbf{Observation préalable.} Un nouveau-né camerounais, comme un nouveau-né japonais ou brésilien, sourit dès les premières semaines et pleure quand il a faim : personne ne le lui a appris. En revanche, un enfant de l'Ouest salue en s'inclinant selon le code de sa chefferie, tandis qu'un enfant du Nord adopte d'autres gestes : il a fallu qu'on le lui montre et qu'il s'exerce. Il existe donc deux grandes sources de comportements.

\begin{definition}[Inné et acquis]
\begin{itemize}
\item Un comportement \cle{inné} est un comportement \textbf{génétiquement programmé}, présent dès la naissance (ou apparaissant spontanément au cours du développement), \textbf{sans apprentissage}. Il est à peu près identique chez tous les individus de l'espèce.
\item Un comportement \cle{acquis} est un comportement \textbf{appris par l'expérience, l'imitation et l'éducation}, au sein d'une culture donnée. Il varie d'un individu et d'une société à l'autre.
\end{itemize}
\end{definition}

\begin{center}
\begin{tabularx}{\linewidth}{|l|X|X|}
\hline
\rowcolor{popblueL} \textbf{Critère} & \textbf{Comportement inné} & \textbf{Comportement acquis} \\
\hline
Origine & Programme génétique (hérédité) & Expérience, imitation, éducation \\
\hline
Apprentissage & Aucun & Nécessaire \\
\hline
Variabilité & Identique dans toute l'espèce & Variable selon les cultures et les individus \\
\hline
Exemples & Sourire du nouveau-né, succion, pleurs de faim, réflexe de retrait & Salutations codifiées, langue parlée, règles de politesse, jeux, techniques agricoles \\
\hline
\end{tabularx}
\end{center}

\begin{attention}[Erreur fréquente au Probatoire]
Ne confonds pas \textbf{inné} et \textbf{acquis} ! Pièges classiques :
\begin{itemize}
\item Le sourire est \textbf{inné} (il existe chez tous les nouveau-nés du monde, même aveugles de naissance), mais les \emph{circonstances} où l'on sourit et les règles de politesse qui l'encadrent sont \textbf{acquises}.
\item Marcher sur deux jambes est une capacité innée ; la démarche dansée lors d'une cérémonie traditionnelle est acquise.
\item La peur devant un danger brutal est innée ; avoir peur d'un fétiche particulier est acquis par la culture.
\end{itemize}
Retiens la question-test : « Faut-il un apprentissage pour que ce comportement apparaisse ? » Si oui $\rightarrow$ acquis ; si non $\rightarrow$ inné. Attention aussi : beaucoup de comportements humains sont \textbf{mixtes} (une base innée façonnée par la culture).
\end{attention}

\begin{savaistu}[Le sourire, un langage universel inscrit dans nos gènes]
Des études menées sur des nouveau-nés de tous les continents ont montré que le sourire apparaît dès les premières semaines de vie, y compris chez les bébés nés aveugles qui n'ont jamais pu voir un visage sourire. C'est donc bien un comportement \textbf{inné}, universel dans l'espèce humaine. Mieux : le sourire d'un bébé déclenche presque automatiquement l'attendrissement et les soins des adultes qui l'entourent. Les biologistes y voient un comportement interindividuel hérité de l'évolution : en « accrochant » affectivement ses parents, le nourrisson, totalement dépendant, augmente ses chances de survie. La nature a donc programmé nos premiers outils de communication sociale avant même notre naissance !
\end{savaistu}

\courssub{Pourquoi maîtriser ses comportements interindividuels ?}

Revenons une dernière fois à notre joueur de football. Son émotion de colère était normale et innée ; c'est sa \emph{réponse} (la violence) qui a posé problème. Or la réponse, elle, peut être modulée par l'éducation, l'expérience et la volonté : c'est là tout l'enjeu de la \cle{maîtrise des comportements}.

\begin{aretenir}[Intérêt de la maîtrise des comportements interindividuels]
\begin{itemize}
\item \textbf{Cohésion sociale} : des comportements d'attraction et de coopération maîtrisés (salutations, entraide, dialogue) renforcent les liens dans la famille, l'école, le quartier et le village.
\item \textbf{Prévention des conflits} : comprendre que la colère est une réaction physiologique normale mais que la riposte violente est un choix permet de désamorcer les disputes avant qu'elles ne dégénèrent (bagarres de stade, conflits communautaires).
\item \textbf{Respect de la diversité} : les comportements acquis variant selon les cultures, connaître leurs bases biologiques aide à ne pas juger ni rejeter autrui pour des différences qui relèvent de l'apprentissage culturel.
\end{itemize}
\end{aretenir}

\begin{exoresolu}[Analyser un comportement interindividuel]
\textbf{Énoncé.} Lors de la fête de la jeunesse dans un lycée de Yaoundé, un élève voit un camarade trébucher et tomber devant toute la classe. Certains élèves éclatent de rire ; un autre se précipite pour l'aider à se relever. (1) Identifie le stimulus social dans cette situation. (2) Pour chacun des deux groupes d'élèves, décris les trois composantes du comportement interindividuel. (3) Indique, en justifiant, si les comportements observés sont plutôt innés ou acquis.
\tcblower
\textbf{Solution détaillée.}
\begin{enumerate}
\item \textbf{Le stimulus social} est la chute du camarade (vision d'un individu qui tombe), événement impliquant un autre individu : il s'agit bien d'une situation interindividuelle.
\item \textbf{Chez les élèves qui rient} : perception = ils voient la chute et l'interprètent comme comique ; émotion = amusement ; communication (réponse) = éclats de rire, un comportement d'attraction entre rieurs mais qui peut blesser le camarade. \textbf{Chez l'élève qui aide} : perception = il voit la chute et l'interprète comme une détresse (« il s'est peut-être fait mal ») ; émotion = inquiétude, empathie ; communication = il se lève, court, tend la main et rassure le camarade : comportement de coopération.
\item La capacité de rire et l'empathie ont une \textbf{base innée} (tous les humains en sont dotés dès la naissance). Mais la \emph{façon} de réagir (se moquer ou porter secours) est largement \textbf{acquise} : elle dépend de l'éducation reçue, de l'imitation des adultes et des valeurs du groupe. Cet exemple illustre le caractère mixte des comportements humains et montre que l'éducation peut orienter nos réponses vers l'entraide plutôt que vers la moquerie.
\end{enumerate}
\end{exoresolu}

\begin{exercice}[À toi de jouer]
Dans ton village, un jeune refusant de saluer les anciens selon le rituel traditionnel provoque la colère des notables. (1) Analyse la situation en identifiant le stimulus social, puis les trois composantes du comportement des notables. (2) Le rituel de salutation est-il un comportement inné ou acquis ? Justifie. (3) Propose une attitude permettant d'éviter le conflit, en t'appuyant sur la notion de diversité culturelle des comportements acquis.
\end{exercice}

\begin{corrige}[Exercice 1]
(1) Stimulus social : le passage du jeune sans saluer. Perception : les notables voient et entendent, puis interprètent ce silence comme un manque de respect. Émotion : colère, sentiment d'offense. Communication : reproches, réprimandes. (2) Le rituel de salutation est un comportement \textbf{acquis} : il faut l'apprendre, il varie d'une culture à l'autre et n'existe pas chez le nouveau-né. (3) Attitude de prévention : le dialogue. Les notables peuvent expliquer calmement la signification du rituel au jeune (qui peut l'ignorer, surtout s'il a grandi en ville), et le jeune peut exprimer son respect par un autre geste accepté. Comprendre que les codes sont acquis et variables évite de transformer une différence culturelle en conflit.
\end{corrige}

\coursec{La perception d'autrui : support nerveux}

Après avoir défini les grands types de comportements interindividuels dans la section précédente, nous devons comprendre comment notre système nerveux nous permet de « lire » autrui. Avant d'interagir, nous percevons : un visage qui se ferme, une voix qui s'élève, une main tendue. Comment ces signaux physiques deviennent-ils des informations sociales exploitables par notre cerveau ? C'est l'objet de cette section.

\courssub{Les organes des sens : portes d'entrée de l'information sociale}

\emph{Observation.} Dans la cour de récréation, tu distingues instantanément un camarade qui rit de celui qui boude, rien qu'en le regardant. Au téléphone, tu devines l'humeur de ton interlocuteur à son ton de voix, sans le voir. Un contact sur l'épaule peut te rassurer ou t'alerter selon la pression exercée.

\begin{definition}[Perception]
La \cle{perception} est l'ensemble des processus par lesquels le système nerveux \textbf{reçoit} des informations sensorielles (stimuli physiques ou chimiques), les \textbf{traite} et les \textbf{interprète} pour construire une représentation consciente du monde extérieur, y compris d'autrui. Elle ne se limite pas à la sensation (détection brute), elle inclut l'attribution de sens.
\end{definition}

\begin{propriete}[Organes des sens impliqués dans la perception sociale]
La perception d'autrui mobilise principalement quatre modalités sensorielles :
\begin{itemize}
    \item \textbf{La vision (œil) :} analyse des expressions faciales (sourcils froncés, sourire, regard fuyant), des postures corporelles, des gestes, de la distance interpersonnelle (proxémie). C'est le canal dominant chez l'humain (environ 80\% de l'information sociale).
    \item \textbf{L'audition (oreille) :} décodage du langage verbal (mots) et surtout \emph{paraverbal} : intonation, débit, timbre, volume, pauses. Une même phrase prononcée sur deux tons différents change radicalement de sens.
    \item \textbf{Le toucher (peau) :} perçoit le contact physique (poignée de main, accolade, tape amicale, agression), la température, la texture, la pression. Essentiel pour l'attachement, la réconfortation ou la menace.
    \item \textbf{L'odorat (nez) :} détecte les odeurs corporelles, les phéromones putatives, l'hygiène, l'environnement olfactif de l'autre. Rôle subtil mais réel dans l'attraction, la reconnaissance mère-enfant, l'évitement de la maladie.
\end{itemize}
\end{propriete}

\begin{exemplebox}[Reconnaître la colère chez un camarade]
\textbf{Situation :} Ton ami Pierre entre en classe, il claque la porte, jette son sac violemment, évite ton regard, ses sourcils sont froncés, sa mâchoire serrée.
\begin{itemize}
    \item \textbf{Vision :} Tu perçois les micro-expressions faciales (froncement des sourcils = muscle corrugator supercilii, serrage de la mâchoire = masséters), la posture raide, le geste brusque.
    \item \textbf{Audition :} Le claquement de porte, le bruit sec du sac, peut-être un souffle bruyant ou un ton sec s'il parle.
    \item \textbf{Toucher (si contact) :} Main froide, moite ou poigne ferme si tu lui serres la main.
    \item \textbf{Intégration :} Ton cerveau assemble ces indices multimodaux en une fraction de seconde : « Pierre est en colère ».
\end{itemize}
\end{exemplebox}

\courssub{Le trajet du message sensoriel : du récepteur au cortex}

\emph{Problème :} Comment l'onde lumineuse réfléchie par le visage de Pierre ou l'onde sonore de sa voix devient-elle l'idée « Pierre est en colère » dans ton esprit ?

\begin{methode}[Décrire le trajet d'un message sensoriel (4 étapes)]
Pour tout message sensoriel (visuel, auditif, tactile, olfactif), le parcours suit un schéma invariant :
\begin{enumerate}
    \item \textbf{Réception (Transduction) :} Au niveau du \cle{récepteur sensoriel} (cônes/bâtonnets de la rétine, cellules ciliées de la cochlée, corpuscules de Meissner/Pacini dans la peau, neurones olfactifs), l'énergie du stimulus (lumière, son, pression, molécule odorante) est convertie en \textbf{potentiel d'action} (influx nerveux). C'est le codage neuronal.
    \item \textbf{Conduction :} L'influx nerveux voyage le long des \cle{nerfs sensitifs} (nerf optique II, nerf auditif VIII, nerfs spinaux pour le toucher, nerf olfactif I) vers le système nerveux central. Vitesse : jusqu'à 100 m/s pour les fibres myélinisées (toucher, vision), plus lent pour les fibres non myélinisées (douleur diffuse, odorat).
    \item \textbf{Traitement sous-cortical (Relais obligatoire) :} Pour la vision, l'audition et le toucher (mais \textbf{pas} l'odorat), le message fait obligatoirement escale au \cle{thalamus} (noyaux géniculés latéraux pour la vision, médiaux pour l'audition, noyau postéro-latéral ventral pour le toucher). Le thalamus filtre, amplifie et dirige l'information vers le cortex. C'est une « porte » vers la conscience.
    \item \textbf{Traitement cortical et Interprétation :} L'information atteint le \cle{cortex cérébral} primaire spécifique (cortex visuel V1, auditif A1, somesthésique S1) pour l'analyse des caractéristiques élémentaires (formes, fréquences, localisation), puis les \cle{aires d'association} (multimodales) pour la reconnaissance d'objets, de visages, de voix, d'intentions (ex : aire fusiforme pour les visages, aire de Wernicke pour le langage, jonction temporo-pariétale pour la théorie de l'esprit).
\end{enumerate}
\end{methode}

\begin{popfigure}
\begin{tikzpicture}[
    node distance=1.8cm and 2.5cm,
    box/.style={rectangle, draw=popdark, thick, fill=popcream, rounded corners, minimum width=2.8cm, minimum height=1.2cm, align=center, font=\small},
    arrow/.style={-Stealth, thick, popdark},
    labelnode/.style={font=\footnotesize, text=popmuted, align=center}
]
% Nodes
\node[box, align=center] (recepteur){Récepteur sensoriel\\ (Rétine, Cochlée, Peau, Épithélium olfactif)};
\node[box, right=of recepteur, align=center] (nerf){Nerf sensitif\\ (Optique, Auditif, Spinaux, Olfactif)};
\node[box, right=of nerf, align=center] (thalamus){Thalamus\\ (Relais, Filtrage, Attention)};
\node[box, right=of thalamus, align=center] (cortex1){Cortex primaire\\ (V1, A1, S1)\\(Analyse élémentaire)};
\node[box, right=of cortex1, align=center] (assoc){Aires d'association\\ (Multimodales)\\(Reconnaissance, Sens)};

% Arrows
\draw[arrow] (recepteur) -- (nerf) node[midway, above, labelnode, align=center]{Transduction \\ Potentiel d'action};
\draw[arrow] (nerf) -- (thalamus) node[midway, above, labelnode, align=center]{Conduction \\ (jusqu'à 100 m/s)};
\draw[arrow] (thalamus) -- (cortex1) node[midway, above, labelnode] {Relais thalamocortical};
\draw[arrow] (cortex1) -- (assoc) node[midway, above, labelnode, align=center]{Intégration \\ Interprétation};

% Exception olfaction
\node[labelnode, below=0.3cm of thalamus] (note) {\textbf{Note :} L'olfaction contourne le thalamus directement vers le cortex piriforme et le système limbique (émotion/mémoire).};
\draw[poporange, dashed, thick] (recepteur.south) -- ++(0,-0.8) -| (note.north);
\end{tikzpicture}
\legende{Schéma du trajet général du message sensoriel (vision, audition, toucher) du récepteur périphérique aux aires d'association corticales. Le thalamus est un relais obligatoire (sauf olfaction).}
\end{popfigure}

\courssub{Les aires cérébrales impliquées : de la sensation à la signification}

\begin{propriete}[Localisation fonctionnelle des aires sensorielles primaires et associatives]
\begin{center}
\begin{tabularx}{\linewidth}{|l|X|X|}\hline
\rowcolor{popblueL}
\textbf{Modalité} & \textbf{Cortex primaire (Analyse élémentaire)} & \textbf{Aires d'association (Interprétation, Sens)} \\ \hline
\textbf{Vision} & Cortex visuel primaire (V1, aire 17 de Brodmann) : détection contours, orientation, mouvement, couleur. & \textbf{Aire fusiforme (FFA) :} reconnaissance des visages.\\
& & \textbf{Cortex temporal supérieur (STS) :} perception des regards, mouvements biologiques, intentions.\\
& & \textbf{Aires pariétales :} localisation spatiale, attention. \\ \hline
\textbf{Audition} & Cortex auditif primaire (A1, aires 41/42) : analyse fréquences, intensité, localisation sonore. & \textbf{Aire de Wernicke (aire 22) :} compréhension du langage verbal.\\
& & \textbf{Cortex temporal antérieur :} reconnaissance des voix, identité du locuteur.\\
& & \textbf{Aires prosodiques (hémisphère droit souvent) :} analyse de l'intonation, émotion dans la voix. \\ \hline
\textbf{Toucher} & Cortex somesthésique primaire (S1, aires 3,1,2) : localisation, intensité, texture (carte homonculus). & \textbf{Cortex pariétal postérieur (aires 5, 7) :} reconnaissance d'objets au toucher (stéréognosie), intégration visuotactile.\\
& & \textbf{Insula postérieure :} toucher affectif (caresses, contact social), interoception. \\ \hline
\textbf{Odorat} & Cortex piriforme (paleocortex) : identification qualité odorante. & \textbf{Cortex orbito-frontal :} hédonie (plaisir/dégoût), prise de décision.\\
& (Pas de relais thalamique direct) & \textbf{Amygdale, Hippocampe :} mémoire émotionnelle, reconnaissance sociale. \\ \hline
\end{tabularx}
\end{center}
\end{propriete}

\begin{experience}[Expérience de pensée : La lésion du cortex visuel vs aire fusiforme]
\textbf{Protocole mental (basé sur des cas cliniques réels) :}
\begin{enumerate}
    \item \textbf{Patient A :} Lésion bilatérale du cortex visuel primaire (V1). \textbf{Résultat :} Cécité corticale. Le patient ne « voit » rien consciemment. Pourtant, s'il doit deviner l'expression d'un visage présenté dans son champ aveugle, il réussit au-dessus du hasard (\emph{blindsight} / vision aveugle). \textbf{Interprétation :} Une voie sous-corticale (colliculus supérieur $\to$ pulvinar thalamique $\to$ amygdale/aire fusiforme) traite l'émotion sans conscience.
    \item \textbf{Patient B :} Lésion sélective de l'aire fusiforme droite (prosopagnosie). \textbf{Résultat :} Le patient voit parfaitement les traits (yeux, nez, bouche) mais ne reconnaît \textbf{pas} le visage de sa mère, ni le sien dans le miroir. Il reconnaît la voix instantanément. \textbf{Interprétation :} L'aire fusiforme est spécialisée dans la \textbf{reconnaissance faciale holistique} (configurale), étape cruciale de l'identification d'autrui.
\end{enumerate}
\textbf{Conclusion :} La perception d'autrui repose sur une hiérarchie : détection (V1) $\to$ analyse structurale (aires occipito-temporales) $\to$ reconnaissance d'identité/émotion (aire fusiforme, STS) $\to$ attribution d'intentions (jonction temporo-pariétale, cortex préfrontal).
\end{experience}

\courssub{Le cortex préfrontal : juge des intentions et décideur social}

\begin{popfigure}
\begin{tikzpicture}[
    scale=0.9,
    every node/.style={font=\small},
    brain/.style={draw=popdark, thick, fill=popcream},
    label/.style={font=\footnotesize, text=popdark, align=center}
]
% Coupe sagittale simplifiée (hémisphère gauche)
% Contour cerveau
\draw[brain, line width=1.5pt] (0,0) .. controls (0,3) and (3,5) .. (6,5) .. controls (8,5) and (9,3) .. (9,0) .. controls (9,-1) and (7,-2) .. (5,-2) .. controls (3,-2) and (1,-1) .. (0,0);
% Corps calleux
\draw[fill=popblue!20, draw=popblue, thick] (2,1.5) .. controls (2,1) and (4,1) .. (4,1.5) .. controls (4,2) and (2,2) .. (2,1.5);
% Cervelet / Tronc (esquisse)
\draw[brain, line width=1pt] (5.5,-2) .. controls (6,-2.5) and (7,-2.5) .. (7.5,-2) .. controls (7.5,-1.5) and (6.5,-1.5) .. (5.5,-1.5) -- cycle;

% Zones fonctionnelles (remplissage semi-transparent + étiquettes)
% Cortex Préfrontal (CPF)
\fill[popblue!30] (0.5,1.5) .. controls (0.5,3.5) and (3,4.5) .. (4.5,4) .. controls (4.5,2.5) and (3,1.5) .. (0.5,1.5);
\node[label, popblue, align=center] at (2.2,3.5){\textbf{Cortex préfrontal (CPF)}\\(Dorsolatéral, Ventromédian, Orbitofrontal)\\\textit{Planification, Décision, Inhibition, Théorie de l'esprit}};

% Système Limbique (Amygdale, Hippocampe)
\fill[poporange!30] (4.5,1.5) .. controls (4.5,2.5) and (6,3) .. (6.5,2.5) .. controls (6.5,1) and (5.5,0.5) .. (4.5,1.5);
\node[label, poporange, align=center] at (5.5,2.2){\textbf{Système limbique}\\(Amygdale, Hippocampe)\\\textit{Émotion, Mémoire, Valeur}};

% Hypothalamus (sous thalamus)
\fill[popgreen!30] (4.8,0.8) circle (0.4);
\node[label, popgreen, align=center] at (4.8,0.2){\textbf{Hypothalamus}\\\textit{Autonome, Endocrine}};

% Thalamus
\fill[poppurple!30] (4,1.8) ellipse (0.5 and 0.35);
\node[label, poppurple, align=center] at (4,1.8){\textbf{Thalamus}\\\textit{Relais sensoriel}};

% Cortex Sensoriel/Moteur (Sillon central ~ x=5.5)
\draw[dashed, popmuted] (5.5,5) -- (5.5,-1.5);
\node[label, popmuted] at (6.5,4) {Cortex moteur};
\node[label, popmuted] at (6.5,1) {Cortex somesthésique};

% Cortex Visuel (Occipital ~ x=1)
\fill[poppink!30] (0.5,4.5) .. controls (0.5,5) and (1.5,5) .. (1.5,4.5) -- cycle;
\node[label, poppink] at (1,4.8) {Cortex visuel (V1)};

% Flèches de connexion (simplifiées)
\draw[->, thick, popdark, bend left=20] (2.5,3.2) to node[above, label] {Contrôle top-down} (5.2,2.5);
\draw[->, thick, popdark, bend right=20] (5.2,2.0) to node[below, label] {Signal émotionnel} (2.5,2.5);
\draw[->, thick, popdark] (4,1.8) -- (2.5,3.0); % Thalamus -> CPF
\draw[->, thick, popdark] (4.8,0.8) -- (5.5,0.5); % Hypothalamus -> Autonome
\end{tikzpicture}
\legende{Coupe sagittale médiane schématique du cerveau humain localisant les structures clés de la perception sociale : Cortex préfrontal (bleu, contrôle exécutif, mentalisation), Système limbique (orange, émotion), Thalamus (violet, relais), Hypothalamus (vert, sortie neuroendocrine). Les flèches illustrent les boucles de régulation émotionnelle (CPF $\leftrightarrow$ Amygdale).}
\end{popfigure}

\begin{propriete}[Rôle du cortex préfrontal (CPF) dans la perception sociale]
Le CPF (notamment dorsolatéral, ventromédian et orbitofrontal) ne « perçoit » pas les stimuli bruts, il \textbf{interprète} et \textbf{régule} :
\begin{itemize}
    \item \textbf{Mentalisation / Théorie de l'esprit (ToM) :} Capacité d'inférer les états mentaux d'autrui (croyances, désirs, intentions) distincts des siens. Implique la \textbf{jonction temporo-pariétale (JTP)} et le \textbf{CPF médian}. Exemple : « Pierre claque la porte \emph{parce qu'il pense} que je l'ai trahi, \emph{pas} parce qu'il est méchant par nature ».
    \item \textbf{Évaluation de la valeur sociale :} Le CPF ventromédian/orbitofrontal associe une valeur (récompense/risque) aux stimuli sociaux (visage de confiance, signe de menace). Guide l'approche ou l'évitement.
    \item \textbf{Régulation émotionnelle (Contrôle top-down) :} Le CPF inhibe l'amygdale (peur/colère) via des projections GABAergiques. Permet de ne pas réagir impulsivement à une insulte, de relativiser une critique.
    \item \textbf{Prise de décision sociale :} Intègre normes morales, conséquences à long terme, réputation. Choix : riposter, dialoguer, fuir, pardonner.
\end{itemize}
\end{propriete}

\begin{aretenir}[Synthèse : Hiérarchie du traitement de l'information sociale]
\begin{enumerate}
    \item \textbf{Niveau 1 (Réflexe/Automatique) :} Thalamus $\to$ Amygdale $\to$ Réponse physiologique (cœur, sueur, fuite/attaque). Très rapide (quelques dizaines de ms), non conscient, « sale » (faux positifs). \textit{Ex: Sursaut face à un visage surgissant.}
    \item \textbf{Niveau 2 (Reconnaissance/Catégorisation) :} Cortex primaire $\to$ Aires d'association (Fusiforme, STS) $\to$ Identification : « Visage humain », « Colère », « Pierre ». Rapide ($\sim$ 100-170 ms), largement automatique.
    \item \textbf{Niveau 3 (Mentalisation/Contrôle) :} CPF + JTP $\to$ Attribution d'intention, régulation, décision adaptée. Lent ($\sim$ 300-600 ms), conscient, flexible, contextuel. \textit{Ex: « Pierre est stressé par l'examen, je ne le prends pas pour moi, je lui parle calmement. »}
\end{enumerate}
\end{aretenir}

\courssub{La perception est une construction, pas une copie : influence de l'expérience, de la culture et du contexte}

\begin{attention}[Piège majeur : La perception n'est pas une photographie objective]
\textbf{Erreur fréquente :} Croire que « je vois ce qui est là ». \textbf{Réalité scientifique :} Le cerveau \textbf{prédit} la réalité à partir de modèles internes (attentes, apprentissages antérieurs) et ne corrige ces prédictions qu'avec l'erreur de prédiction (signal sensoriel résiduel). C'est le principe du \textbf{cerveau bayésien} (inférence active).
\begin{itemize}
    \item \textbf{Expérience / Apprentissage :} Un éleveur de bétail au Nord-Cameroun distingue des nuances de comportement chez un taureau (menace imminente vs simple agitation) qu'un citadin ne verra pas. Un médecin urgentiste à l'Hôpital Central de Yaoundé perçoit des signes de détresse vitale invisibles pour un novice.
    \item \textbf{Culture :} Le contact visuel direct est perçu comme \textbf{honnêteté/confiance} dans les cultures occidentales, mais comme \textbf{irrespect/provocation} face à un aîné dans de nombreuses cultures camerounaises (Bamiléké, Bassa, Peul, etc.). Le même stimulus (regard dans les yeux) $\to$ interprétation opposée $\to$ comportement différent (soutien du regard vs baisser les yeux).
    \item \textbf{Contexte immédiat :} Un camarade qui crie « Tu es nul ! » sur un terrain de foot (contexte jeu/banter) vs en classe devant le professeur (contexte humiliation). Même stimulus auditif, perception radicalement différente (amitié vs agression).
    \item \textbf{État interne (Biais) :} Fatigue, faim, stress, humeur dépressive, préjugés (stéréotypes ethniques, de genre, sociaux) agissent comme des \textbf{priors} (antécédents bayésiens) qui déforment l'interprétation. Une personne anxieuse perçoit « menace » là où il n'y a que « neutralité ».
\end{itemize}
\end{attention}

\begin{experience}[Illusion perceptive : Le visage creux (Mask illusion) / McGurk]
\textbf{Matériel :} Masque creux (intérieur d'un masque blanc) éclairé de face, ou vidéo de l'effet McGurk (son « ba » + vidéo « ga » $\to$ perception « da »).
\textbf{Protocole :}
\begin{enumerate}
    \item Observer le masque creux qui tourne. Malgré la connaissance qu'il est creux, le cerveau \textbf{impose} la perception d'un visage convexe (normal) car le \emph{prior} « les visages sont convexes » est massivement plus fort que le signal visuel (ombrages inversés).
    \item Regarder la vidéo McGurk : fermer les yeux $\to$ on entend « ba » ; ouvrir les yeux $\to$ on entend « da » (fusion visuo-auditive).
\end{enumerate}
\textbf{Interprétation :} La perception est une \textbf{inférence multimodale}. Le cerveau combine vue et ouïe pour choisir l'hypothèse la plus probable. La connaissance consciente ne « corrige » pas l'illusion : le traitement perceptif bas niveau est \textbf{encapsulé} (impenétrable cognitivement).
\textbf{Conclusion :} Nous ne percevons pas le monde « tel qu'il est », mais « tel qu'il est utile de le percevoir pour survivre et interagir » selon notre histoire phylogénétique et ontogénétique.
\end{experience}

\begin{savaistu}[Les neurones miroirs : une base neuronale de l'empathie ?]
Découverts chez le singe (équipe de Giacomo Rizzolatti, Parme, années 90) dans la zone F5 (homologue de l'aire de Broca chez l'homme), ces neurones s'activent \textbf{à la fois} quand le singe \textbf{exécute} une action (saisir une arachide) \textbf{et} quand il \textbf{observe} un congénère (ou l'expérimentateur) exécuter la \textbf{même action}.
\begin{itemize}
    \item \textbf{Chez l'homme :} Système miroir identifié dans le cortex pariétal inférieur, cortex préfrontal ventral (aire 44/45), cortex moteur supplémentaire, insula.
    \item \textbf{Fonction hypothétique :} Simulation interne de l'action d'autrui $\to$ compréhension immédiate de l'intention motrice (« il veut saisir ») sans raisonnement explicite. Base du \textbf{contagion émotionnelle} (voir la douleur active notre propre matrice de la douleur : insula, cingulaire antérieur).
    \item \textbf{Nuance cruciale :} Les neurones miroirs ne suffisent pas à l'empathie complexe (prise de perspective, souci de l'autre). Ils fournissent une \textbf{résonance corporelle} automatique. L'empathie cognitive (Théorie de l'Esprit) requiert le CPF et la JTP. L'empathie affective requiert la régulation (ne pas être submergé).
    \item \textbf{Contexte camerounais :} Lors des danses traditionnelles (ex: \emph{Njock} chez les Bassa, \emph{Tso} chez les Bamiléké), la synchronisation motrice et émotionnelle du groupe repose largement sur ce système de résonance motrice partagée, renforçant la cohésion sociale.
\end{itemize}
\end{savaistu}

\courssub{Conséquence majeure : Subjectivité de la perception interindividuelle}

\begin{exoresolu}[Analyse d'un malentendu interculturel]
\textbf{Énoncé :} Lors d'une réunion parents-professeurs dans un lycée de Douala, M. Kouam (professeur, formation occidentale) demande à M. Essomba (père d'élève, tradition bassa) : « Votre fils a-t-il des difficultés ? » M. Essomba baisse les yeux, ne répond pas immédiatement, puis dit doucement : « Il fait de son mieux. » M. Kouam note : « Père peu concerné, évite le regard, minimise les problèmes. » M. Essomba pense : « Ce professeur est agressif, il me fixe, il ne respecte pas mon âge, il veut m'humilier publiquement. »

\textbf{Analyse neuro-cognitive et culturelle :}
\begin{enumerate}
    \item \textbf{Stimulus identique :} Regard direct soutenu de M. Kouam $\to$ Rétine $\to$ V1 $\to$ Aire fusiforme $\to$ Amygdale (évaluation menace/valeur).
    \item \textbf{Traitement divergent (Niveau 3 - CPF/Culture) :}
    \begin{itemize}
        \item \textbf{M. Kouam (Prior : « Regard direct = Honnêteté, Engagement, Écoute active ») :} Interprétation : « Je suis ouvert, j'écoute. » $\to$ Comportement : Maintien du regard, posture face à face.
        \item \textbf{M. Essomba (Prior : « Regard direct sur un aîné/supérieur = Défi, Irrespect, Agression ») :} Interprétation : « Il m'agresse. » $\to$ Activation amygdale (peur/colère) $\to$ Inhibition CPF (soumission/respect) $\to$ Comportement : Baisse des yeux (soumission), réponse évasive (protection face), ton bas (apaisement).
    \end{itemize}
    \item \textbf{Boucle de rétroaction :} Le regard baissé de M. Essomba confirme le biais de M. Kouam (« évite le regard = culpabilité/désintérêt »). La posture fixe de M. Kouam confirme le biais de M. Essomba (« regard fixe = agression »).
\end{enumerate}
\textbf{Conclusion :} Une même réalité physique (photons réfléchis par les yeux de M. Kouam) a généré deux perceptions sociales antagonistes, sources de conflit. La \textbf{compétence interculturelle} consiste à prendre conscience de ses propres \emph{priors} culturels et à les inhiber (CPF) pour réinterpréter le comportement d'autrui dans \textbf{son} cadre de référence.
\end{exoresolu}

\begin{exemplebox}[Application concrète : Gestion d'un conflit en classe]
\textbf{Situation :} Deux élèves, Amina et Boris, se disputent violemment. L'enseignant intervient.
\begin{itemize}
    \item \textbf{Perception naïve (Niveau 1/2) :} L'enseignant voit « bagarre », entend « cris ». Amygdale $\to$ Stress $\to$ Réaction punitive immédiate (« Silence ! Tous les deux au bureau ! »).
    \item \textbf{Perception experte (Niveau 3 - Mentalisation) :} L'enseignant inhibe sa réaction (CPF dorsolatéral). Il observe : Amina a les poings serrés, regard fuyant, larmes (peur/humiliation). Boris a le menton haut, regard fixe, sourire en coin (domination/provocation). Il infère : Amina a subi une insulte raciste/sexiste (contexte), Boris teste les limites.
    \item \textbf{Action adaptée :} Séparer sans humilier. Parler à Amina en privé (sécurité, validation). Sanctionner Boris pour le fond (insulte) pas la forme (bagarre), avec explication éducative. Restaurer le climat de classe.
\end{itemize}
\textbf{Lien avec le programme :} C'est l'application directe du savoir-faire « Expliquer un comportement interindividuel à partir de son support physiologique » et « Adopter une attitude de respect face à la diversité des comportements ».
\end{exemplebox}

\begin{exercice}[Entraînement au Probatoire]
\textbf{Exercice 1.} Un élève présente une lésion sélective de l'aire fusiforme droite (prosopagnosie) suite à un AVC. Il ne reconnaît plus les visages de sa famille, mais reconnaît leurs voix. Expliquez, en utilisant les termes \textbf{cortex visuel primaire}, \textbf{aire fusiforme}, \textbf{aires d'association}, \textbf{cortex auditif}, pourquoi la reconnaissance vocale est préservée alors que la reconnaissance faciale est abolie.

\textbf{Exercice 2.} Lors d'un entretien d'embauche, le candidat évite le regard du recruteur. Le recruteur note « Manque de franchise ». Le candidat vient d'une culture où baisser les yeux face à un aîné est une marque de respect.
\begin{enumerate}
    \item Identifiez l'organe des sens et le trajet neural principal impliqués dans la perception du regard.
    \item Expliquez pourquoi le même stimulus (regard évité) conduit à deux interprétations opposées en mobilisant les concepts de \textbf{prior bayésien / expérience culturelle} et de \textbf{cortex préfrontal}.
    \item Proposez une attitude professionnelle du recruteur pour éviter ce biais (savoir-être).
\end{enumerate}
\end{exercice}


\begin{corrige}[Exercice 1]
\textbf{1. Voie visuelle (déficiente) :} Lumière (visage) $\to$ Rétine $\to$ Nerf optique $\to$ Thalamus (Noyau géniculé latéral) $\to$ \textbf{Cortex visuel primaire (V1)} : analyse intacte des traits (yeux, nez, bouche présents). $\to$ Projection vers \textbf{Aire fusiforme (FFA)} : \textbf{Lésion ici} $\to$ Impossible d'intégrer les traits en une configuration holistic « visage connu » $\to$ Pas de signal de familiarité $\to$ Pas de reconnaissance d'identité.
\textbf{2. Voie auditive (intacte) :} Son (voix) $\to$ Oreille $\to$ Nerf auditif $\to$ Thalamus (Noyau géniculaire médial) $\to$ \textbf{Cortex auditif primaire (A1)} $\to$ \textbf{Aires d'association temporales antérieures} (reconnaissance voix) $\to$ Reconnaissance identité préservée.
\textbf{Conclusion :} La lésion est sélective à un module d'association visuel (FFA), les modules primaires (V1, A1) et l'association auditive sont intacts.
\end{corrige}

\begin{corrige}[Exercice 2]
\textbf{1.} Organe : \textbf{Œil} (récepteurs : cônes/bâtonnets). Trajet : Rétine $\to$ Nerf optique $\to$ Thalamus (NGL) $\to$ Cortex visuel V1 $\to$ Aire fusiforme (détection regard) $\to$ STS (direction regard) $\to$ Amygdale/CPF (interprétation sociale).
\textbf{2.} \textbf{Recruteur (Culture « Occidentale/Entreprise ») :} Prior $P(\text{Regard évité} | \text{Malhonnêteté/Timidité})$ élevé. CPF valide cette hypothèse $\to$ Interprétation « Manque de franchise ».
\textbf{Candidat (Culture « Traditionnelle/Respect ») :} Prior $P(\text{Regard direct} | \text{Irrespect/Agression})$ élevé. Pour éviter l'agression perçue, il active le contrôle (CPF) pour inhiber le regard direct $\to$ Comportement « Baisser les yeux » = Respect.
\textbf{Divergence :} Même entrée sensorielle $\to$ Modèles internes (priors) culturels différents $\to$ Inférences bayésiennes opposées.
\textbf{3.} Attitude professionnelle (Savoir-être) : \textbf{Suspendre son jugement} (métacognition CPF). \textbf{Questionner} plutôt qu'inférer : « Je remarque que vous baissez les yeux, est-ce une marque de respect dans votre culture ? » $\to$ Mise à jour du prior (apprentissage) $\to$ Perception révisée $\to$ Équité.
\end{corrige}

\coursec{Les émotions : bases nerveuses et hormonales}

Après avoir vu comment notre système nerveux construit la perception d'autrui (section précédente), nous comprenons que cette perception n'est jamais neutre : elle est immédiatement teintée d'affectivité. Une même scène — un camarade qui rit, un enseignant qui élève la voix, un inconnu qui s'approche — peut déclencher la joie, la peur ou la colère selon l'histoire personnelle et l'état physiologique de l'observateur. C'est cette coloration affective, brève et intense, que l'on appelle \cle{émotion}. Elle est le moteur de nos réactions sociales les plus immédiates : fuite, approche, agression, protection. Comprendre ses bases physiologiques, c'est accéder aux leviers de la maîtrise de soi, compétence centrale pour des relations interindividuelles apaisées.

\courssub{Définition et manifestations corporelles de l'émotion}

\begin{definition}[L'émotion : réaction affective globale]
Une \textbf{émotion} est une réaction affective brève, intense, déclenchée par un stimulus (externe ou interne), qui s'accompagne de modifications physiologiques visibles (rythme cardiaque, respiration, sueur, rougeur, tension musculaire) et d'une expérience subjective (sentiment de peur, de joie, de colère, de dégoût, de surprise, de tristesse). Elle diffère de l'humeur (plus diffuse, plus durable) et du sentiment (construction cognitive plus stable).
\end{definition}

\begin{experience}[Observation : mon corps réagit avant que je ne décide]
\textbf{Matériel :} Un chronomètre, un tensiomètre électronique, un oxymètre de pouls (ou simplement la prise du pouls radial). \\
\textbf{Protocole :}
\begin{enumerate}
  \item Mesurer au repos (assis, calme) : fréquence cardiaque (FC), tension artérielle (TA), fréquence respiratoire (FR).
  \item Provoquer une émotion brève et intense : bruit soudain (ballon qui éclate), image surprenante, annonce inattendue (« Vous avez une note de 20/20 ! » ou « Le proviseur veut vous voir »).
  \item Mesurer à nouveau FC, TA, FR immédiatement, puis à 1 min, 3 min, 5 min.
\end{enumerate}
\textbf{Observations types :} FC passe de 70 à 130 battements/min en quelques secondes ; TA augmente (ex. 110/70 $\rightarrow$ 140/90 mmHg) ; FR s'accélère ; mains moites, pâleur puis rougeur au visage. Retour progressif à la normale en 5--10 min.
\textbf{Interprétation :} L'émotion mobilise l'organisme tout entier via le système nerveux autonome et les hormones. C'est une \emph{réaction d'urgence} préparant l'action (fuite ou combat).
\end{experience}

\courssub{Le système limbique : centre cérébral des émotions}

\begin{definition}[Système limbique et structures clés]
Le \textbf{système limbique} (ou \emph{cerveau limbique}) est un ensemble de structures cérébrales profondes, de part et d'autre du thalamus, impliquées dans les émotions, la mémoire et la motivation. Les trois structures majeures pour l'émotion sont :
\begin{itemize}
  \item L'\textbf{amygdale} (corps amygdalien) : détection de la menace, déclenchement de la peur et de l'agressivité ; apprentissage de la peur (conditionnement).
  \item L'\textbf{hippocampe} : contextualisation de l'émotion (lieu, temps, souvenirs associés) ; lien émotion-mémoire.
  \item L'\textbf{hypothalamus} : centre de commande de la réponse corporelle. Il active le \textbf{système nerveux orthosympathique} (dont la médullosurrénale $\rightarrow$ adrénaline) et l'\textbf{axe hypothalamo-hypophyso-surrénalien} (corticosurrénale $\rightarrow$ cortisol).
\end{itemize}
\end{definition}

\begin{popfigure}
\begin{tikzpicture}[
  scale=0.95,
  every node/.style={font=\small},
  brain/.style={draw=popdark, fill=popblueL, rounded corners=8pt, thick},
  label/.style={draw=popdark, fill=white, rounded corners=4pt, inner sep=3pt, text width=2.2cm, align=center},
  arrow/.style={-Stealth, thick, popdark}
]
% Hémisphère droit (vue latérale simplifiée)
\draw[brain] (0,0) ellipse (4cm and 2.5cm); 

% Structures profondes (positionnées schématiquement)
\node[brain, label] (amyg) at (-2.5,-0.5) {Amydale};
\node[brain, label] (hippo) at (-1.5,0.5) {Hippocampe};
\node[brain, label] (hypo) at (-0.5,-1.2) {Hypothalamus};

% Structures corticales
\node[brain, label] (pref) at (1.8,1.5) {Cortex préfrontal};
\node[brain, label] (sens) at (0.5,0.5) {Cortex sensoriel};

% Thalamus (relais central, profond)
\node[draw=popdark, fill=poporangeL, circle, inner sep=2pt, label=right:Thalamus] (thal) at (0.2,0) {};

% Flèches voie rapide (thalamus -> amygdale)
\draw[arrow, poporange, line width=2pt] (thal.east) .. controls (1.5,-0.5) .. (amyg.south east);
\node[poporange, font=\small\bfseries] at (1.5,-1.2) {Voie rapide};

% Flèches voie lente (thalamus -> cortex sensoriel -> préfrontal -> amygdale)
\draw[arrow, popblue, line width=2pt] (thal.north east) .. controls (0.5,1.2) .. (sens.south west);
\draw[arrow, popblue, line width=2pt] (sens.north east) .. controls (1.2,1.5) .. (pref.south west);
\draw[arrow, popblue, line width=2pt] (pref.south west) .. controls (0,-0.5) .. (amyg.north west);
\node[popblue, font=\small\bfseries] at (1.8,2.0) {Voie lente (via cortex)};

% Flèches sortie hypothalamus -> corps
\draw[arrow, popgreen, line width=2pt] (hypo.south) -- (-0.5,-2.8) node[below, label, text width=3cm] {Système nerveux autonome \\ Glandes surrénales (adrénaline, cortisol)};
\end{tikzpicture}
\legende{Schéma fonctionnel du système limbique et des deux voies de traitement de l'information émotionnelle. La voie rapide (orange) : thalamus $\rightarrow$ amygdale (réaction immédiate, non consciente). La voie lente (bleue) : thalamus $\rightarrow$ cortex sensoriel $\rightarrow$ cortex préfrontal $\rightarrow$ amygdale (évaluation consciente, régulation). L'hypothalamus (vert) commande la réponse corporelle via le SNA et l'axe HHS.}
\end{popfigure}

\courssub{Les deux voies de l'information émotionnelle : rapidité vs contrôle}

\begin{propriete}[Double trajet de l'information émotionnelle (admis)]
Lorsqu'un stimulus émotionnel (ex. visage menaçant) atteint le thalamus, l'information suit \textbf{deux voies parallèles} :
\begin{enumerate}
  \item \textbf{Voie rapide (sous-corticale) :} Thalamus $\rightarrow$ Amygdale. Temps : \textbf{< 100 ms}. Pas de traitement cortical conscient. Déclenche la réponse physiologique (peur, fuite) \emph{avant} la conscience de l'émotion. Base de la réaction impulsive.
  \item \textbf{Voie lente (corticale) :} Thalamus $\rightarrow$ Cortex sensoriel $\rightarrow$ Cortex préfrontal $\rightarrow$ Amygdale. Temps : \textbf{300--500 ms}. Permet l'analyse fine, la contextualisation (hippocampe), la régulation (cortex préfrontal) et la décision adaptée.
\end{enumerate}
La voie rapide assure la survie immédiate ; la voie lente permet l'adaptation sociale.
\end{propriete}

\begin{savaistu}[L'amygdale, sentinelle ultra-rapide]
Des études en imagerie (IRMf) et chez des patients épileptiques (électrodes intracrâniennes) montrent que l'amygdale réagit à un visage exprimant la peur ou la colère en \textbf{moins de 100 millisecondes}, bien avant que le sujet n'en ait conscience (temps de prise de conscience $\approx$ 300--500 ms). C'est pourquoi on peut « sauter » en entendant un bruit fort avant même de savoir de quoi il s'agit. Au Cameroun, cette réactivité a sauvé des vies face aux serpents ou aux véhicules surgissant dans les rues de Douala ou Yaoundé.
\end{savaistu}

\courssub{Rôle du système nerveux autonome et de l'adrénaline}

\begin{definition}[Système nerveux autonome (SNA) et réponse émotionnelle]
Le \textbf{SNA} assure l'innervation involontaire des organes (cœur, vaisseaux, poumons, tube digestif, glandes). Il a deux divisions antagonistes :
\begin{itemize}
  \item \textbf{Système orthosympathique (sympathique) :} « Accélérateur ». Activé par l'amygdale via l'hypothalamus. Libère de la \textbf{noradrénaline} aux synapses neuro-effectoriennes et stimule la médullosurrénale à sécréter de l'\textbf{adrénaline} dans le sang. Effets : $\uparrow$ FC, $\uparrow$ TA, $\uparrow$ FR, dilatation des bronches, dilatation des pupilles (mydriase), transpiration, inhibition digestion, mobilisation glucose (glycogénolyse hépatique). Prépare à \textbf{l'action : fuite ou combat}.
  \item \textbf{Système parasympathique (vagal) :} « Frein ». Dominant au repos. Ralentit le cœur, stimule la digestion, constriction pupillaire. Permet le retour au calme après l'émotion.
\end{itemize}
\end{definition}

\begin{exemplebox}[Exemple chiffré : la frayeur du motard à Bamenda]
Un motard-taxi (« bendskin ») roule tranquillement (FC 70 bpm). Une voiture lui coupe brutalement la route. En 3 secondes : amygdale activée $\rightarrow$ hypothalamus $\rightarrow$ orthosympathique + adrénaline. Son cœur grimpe à \textbf{145 bpm}, TA 150/95 mmHg, mains moites, mydriase (vision en tunnel), jambes tremblantes (glucose musculaire). Il freine et insulte le conducteur (colère = combat). 5 minutes plus tard, parasympathique reprend le dessus : FC redescend à 85 bpm, il reprend son souffle. Sans cette montée d'adrénaline, son temps de réaction aurait été trop long pour éviter l'accident.
\end{exemplebox}

\begin{popfigure}
\begin{tikzpicture}
\begin{axis}[
  width=\linewidth,
  height=8cm,
  xlabel={Temps (minutes)},
  ylabel={Fréquence cardiaque (bpm)},
  xmin=0, xmax=10,
  ymin=50, ymax=160,
  xtick={0,1,2,3,4,5,6,7,8,9,10},
  ytick={60,80,100,120,140,160},
  grid=major,
  grid style={popline},
  axis line style={popdark},
  tick style={popdark},
  legend style={at={(0.98,0.98)}, anchor=north east, fill=white, draw=popdark},
  popcurve/.style={thick, poporange, mark=*}
]
% Données simulées : repos, pic à 1 min, redescente exponentielle
\addplot[popcurve] coordinates {
  (0,70)
  (0.5,72)
  (1,145)
  (1.5,138)
  (2,125)
  (3,105)
  (4,92)
  (5,84)
  (6,78)
  (7,74)
  (8,72)
  (9,71)
  (10,70)
};
\addlegendentry{FC observée}
% Ligne de base
\addplot[dashed, popmuted, thick] coordinates {(0,70) (10,70)};
\addlegendentry{FC de base (70 bpm)}
% Zones
\fill[poporange!15] (0.5,50) rectangle (2,160);
\node[poporange!80!black, rotate=90, font=\small] at (1.25,155) {Phase orthosympathique};
\fill[popgreen!15] (2,50) rectangle (10,160);
\node[popgreen!80!black, rotate=90, font=\small] at (6,155) {Récupération parasympathique};
\end{axis}
\end{tikzpicture}
\legende{Évolution de la fréquence cardiaque avant, pendant et après une frayeur intense (simulation réaliste). Pic à 145 bpm vers 1 min (effet adrénaline + orthosympathique), puis retour exponentiel vers la base sous contrôle parasympathique.}
\end{popfigure}

\courssub{Le cortex préfrontal : régulation et maîtrise de soi}

\begin{propriete}[Rôle inhibiteur du cortex préfrontal (admis)]
Le \textbf{cortex préfrontal (CPF)}, notamment sa partie ventromédiane et dorsolatérale, reçoit l'information via la voie lente. Il :
\begin{itemize}
  \item Évalue le stimulus (menace réelle ? fausse alerte ? contexte social ?).
  \item Inhibe l'activité de l'amygdale via des projections GABAergiques (frein chimique).
  \item Planifie une réponse adaptée (paroles, geste mesuré, recul) au lieu de la réaction brute.
\end{itemize}
C'est la base physiologique de la \textbf{maîtrise de soi} (contrôle des impulsions). Un CPF immature (adolescent) ou lésé (traumatisme, addiction, stress chronique) $\rightarrow$ difficultés de régulation émotionnelle $\rightarrow$ comportements impulsifs, agressivité, anxiété.
\end{propriete}

\begin{methode}[Expliquer un comportement impulsif en 3 temps]
Pour analyser une réaction impulsive (ex. gifle, fuite, cri, insulte) :
\begin{enumerate}
  \item \textbf{Déclencheur amygdalien :} Identifier le stimulus perçu comme menaçant (visage, mot, bruit) $\rightarrow$ activation voie rapide thalamus-amygdale $\rightarrow$ réponse émotionnelle brute (peur/colère) en < 100 ms.
  \item \textbf{Réponse physiologique orthosympathique + adrénaline :} Hypothalamus $\rightarrow$ SNA orthosympathique + médullosurrénale $\rightarrow$ adrénaline dans le sang $\rightarrow$ corps prêt à l'action (FC $\uparrow$, force musculaire $\uparrow$, inhibition réflexion).
  \item \textbf{Régulation tardive (ou absente) par le cortex préfrontal :} Voie lente arrive 300--500 ms plus tard. Si CPF fonctionnel et entraîné : inhibition de l'amygdale, choix d'une réponse adaptée. Si CPF immature, fatigué, sous alcool, ou surcharge cognitive : inhibition défaillante $\rightarrow$ passage à l'acte impulsif.
\end{enumerate}
\end{methode}

\courssub{Lien émotion -- comportement social : exemples fondamentaux}

\begin{center}
\begin{tabularx}{\linewidth}{|l|X|X|X|}
\hline
\rowcolor{popblueL}
\textbf{Émotion de base} & \textbf{Déclencheur typique (amygdale)} & \textbf{Réponse physiologique (orthosympathique + adrénaline)} & \textbf{Comportement social adaptatif (si CPF régule)} \\
\hline
\cle{Peur} & Menace pour l'intégrité (danger, inconnu, hauteur) & Fuite physiologique : FC $\uparrow\uparrow$, sang vers muscles, vigilance max & \textbf{Fuite} ou \textbf{protection} (se cacher, appeler au secours, éviter le danger) \\
\hline
\cle{Colère} & Obstacle à un but, injustice, agression, territoire violé & Combat physiologique : FC $\uparrow$, TA $\uparrow\uparrow$, force, visage rouge & \textbf{Affirmation} (dire non, négocier, poser limite) ou \textbf{agressivité} si non régulée \\
\hline
\cle{Joie} & Réussite, rencontre aimée, jeu, récompense & Ouverture : FC modérée $\uparrow$, respiration ample, ocytocine, dopamine & \textbf{Rapprochement} : sourire, partage, étreinte, coopération \\
\hline
\cle{Tristesse} & Perte, séparation, échec & Repli : FC $\downarrow$, énergie basse, larmes (parasympathique) & \textbf{Recherche de soutien} : pleurer appelle l'empathie, renforce liens \\
\hline
\cle{Dégoût} & Aliment toxique, odeur nauséabonde, transgression morale & Rejet : nausée, révulsion, FC variable & \textbf{Éloignement} : recracher, s'éloigner, poser tabou social \\
\hline
\end{tabularx}
\end{center}

\begin{attention}[Piège fréquent : confondre émotion et comportement]
\begin{itemize}
  \item \textbf{Erreur :} « La colère, c'est mal. » $\rightarrow$ Non. La colère est une émotion \emph{utile} (signale une limite franchie). C'est le \textbf{comportement agressif} (frapper, insulter) qui peut être inadapté.
  \item \textbf{Erreur :} « Il ne faut pas avoir peur. » $\rightarrow$ La peur protège. L'absence de peur (lésion amygdale) conduit à des prises de risques mortelles.
  \item \textbf{À retenir :} L'émotion est un \textbf{signal d'information} sur nos besoins et l'environnement. La compétence sociale, c'est \textbf{accueillir l'émotion} (voie rapide) \textbf{ET choisir le comportement} (voie lente + CPF).
\end{itemize}
\end{attention}

\courssub{Exercice résolu : analyse d'une situation de vie courante}

\begin{exoresolu}[La dispute au marché de Nkolbisson]
\textbf{Énoncé :} Marie vend des tomates. Un client prend une tomate, la mord, la repose en disant « Elle n'est pas mûre ! » sans payer. Marie ressent une bouffée de chaleur, son cœur bat la chamade, elle a envie de crier et de jeter la tomate à la figure du client. Elle inspire profondément, compte jusqu'à cinq, et dit calmement : « Monsieur, la tomate est mûre, elle coûte 100 F. Si vous l'avez mordue, vous la prenez. » Analysez cette situation en mobilisant les notions du cours.

\textbf{Solution détaillée :}
\begin{enumerate}
  \item \textbf{Stimulus déclencheur :} Action du client (vol + dépréciation du travail) perçue comme menace pour l'intégrité (respect, revenu) et injustice $\rightarrow$ information visuelle/auditive vers thalamus.
  \item \textbf{Voie rapide (amygdale) :} En < 100 ms, amygdale identifie « agression/vol » $\rightarrow$ active hypothalamus $\rightarrow$ décharge orthosympathique + adrénaline. \textbf{Manifestations :} bouffée de chaleur (vasodilatation cutanée), tachycardie (FC $\uparrow$), tension musculaire (envie de jeter/crier = combat). C'est la \textbf{colère} brute.
  \item \textbf{Voie lente (cortex préfrontal) :} Information arrive au cortex sensoriel $\rightarrow$ CPF (300--500 ms). CPF contextualise : « C'est un client, pas un ennemi physique. Crier nuit à mon commerce. La loi m'autorise à exiger paiement. » $\rightarrow$ CPF inhibe l'amygdale (GABA) $\rightarrow$ réduit la réponse orthosympathique.
  \item \textbf{Comportement choisi :} Respiration profonde (active le vague/parasympathique $\rightarrow$ frein cardiaque), comptage (détourne l'attention, donne du temps au CPF), réponse verbale assertive (non agressive, claire, ferme). \textbf{Résultat :} Conflit désamorcé, revenu préservé, dignité gardée.
\end{enumerate}
\textbf{Conclusion :} Marie a utilisé sa régulation corticale pour transformer une impulsion de combat en une action sociale adaptée. C'est l'exemple type de la maîtrise de soi reposant sur l'équilibre amygdale / cortex préfrontal.
\end{exoresolu}

\begin{exercice}[À vous de jouer : le silence en classe]
Paul, élève de Première A, doit passer au tableau pour résoudre un exercice. En se levant, il trébuche sur son banc, la classe éclate de rire. Paul sent son visage brûler, son cœur battre très fort, il a envie de fuir la salle. Il reste pourtant au tableau, rouge, tremblant, et résout l'exercice.
\begin{enumerate}
  \item Nommez l'émotion ressentie par Paul et la structure cérébrale qui l'a déclenchée en premier.
  \item Décrivez les deux voies nerveuses empruntées par l'information « rire de la classe ».
  \item Expliquez physiologiquement pourquoi Paul a pu rester au tableau malgré l'envie de fuir.
  \item Quel rôle a joué le système nerveux parasympathique pendant qu'il résolvait l'exercice ?
\end{enumerate}
\end{exercice}

\begin{corrige}[Exercice : le silence en classe]
\begin{enumerate}
  \item \textbf{Émotion :} Honte / peur sociale (évaluation négative par le groupe). \textbf{Structure :} Amygdale (détection de la menace pour l'estime de soi / statut social).
  \item \textbf{Voie rapide :} Thalamus $\rightarrow$ Amygdale (réaction immédiate : rougeur, tachycardie, envie de fuite). \textbf{Voie lente :} Thalamus $\rightarrow$ Cortex visuel/auditif $\rightarrow$ Cortex préfrontal (analyse : « Ce sont mes camarades, ce n'est pas un danger vital, je dois finir l'exercice »).
  \item Paul est resté grâce à la \textbf{régulation par le cortex préfrontal} : inhibition de la réponse de fuite commandée par l'amygdale, choix d'un comportement orienté vers l'objectif (réussir l'exercice, ne pas montrer sa faiblesse). Le CPF a modulé l'activité amygdalienne.
  \item Pendant la résolution, le \textbf{parasympathique} (nerf vague) a été activé (respiration plus calme, concentration) pour contrebalancer l'orthosympathique : FC diminue progressivement, trembling diminue, retour progressif à l'homéostasie.
\end{enumerate}
\end{corrige}

\begin{aretenir}[Synthèse : les émotions, moteurs de la vie sociale]
\begin{itemize}
  \item L'émotion est une réaction brève, globale (corps + cerveau), adaptative.
  \item \textbf{Amygdale} = détection ultra-rapide de la valence émotionnelle (voie rapide).
  \item \textbf{Hypothalamus} = chef d'orchestre de la réponse corporelle (SNA + adrénaline/cortisol).
  \item \textbf{Cortex préfrontal} = régulateur, frein, décideur social (voie lente).
  \item \textbf{Maîtrise de soi} = capacité du CPF à inhiber l'amygdale pour choisir un comportement adapté au contexte.
  \item Aucune émotion n'est « négative » en soi ; c'est l'absence de régulation qui génère des conduites à risque.
\end{itemize}
\end{aretenir}

\coursec{La communication interindividuelle}

Après avoir exploré les fondements nerveux et hormonaux des émotions, nous comprenons maintenant que l'état interne d'un individu (peur, joie, colère) ne reste pas enfermé : il s'exprime, se partage, se négocie. C'est tout l'objet de la communication. Imagine une scène au marché Mokolo à Yaoundé : une marchande de tomates crie ses prix (verbal), mais son regard fuyant et ses mains croisées (non verbal) trahissent sa fatigue ou son agacement. Un acheteur perçoit ce message contradictoire et adapte sa négociation. Comment cet échange d'informations s'organise-t-il biologiquement ? Quels sont les supports cérébraux du langage, cette capacité spécifiquement humaine ? Et pourquoi le « non-verbal » pèse-t-il souvent plus lourd que les mots ?

\courssub{Définition et modèle général de la communication}

\begin{definition}[Communication interindividuelle]
La \cle{communication interindividuelle} est l'ensemble des processus par lesquels un individu (l'\cle{émetteur}) transmet une information (le \cle{message}) à un autre individu (le \cle{récepteur}) via un \cle{canal} sensoriel (visuel, auditif, tactile, chimique), en vue de modifier le comportement ou l'état cognitif du récepteur. Tout acte de communication suppose un \cle{code} partagé (langue, gestuelle culturelle) et une \cle{rétroaction} (feedback) permettant à l'émetteur de vérifier la bonne réception.
\end{definition}

\begin{popfigure}
\begin{tikzpicture}[
    node distance=2.5cm and 3cm,
    box/.style={rectangle, draw=popdark, thick, fill=popcream, rounded corners, minimum width=2.8cm, minimum height=1.2cm, align=center, font=\small},
    arrow/.style={-Stealth, thick, popblue},
    backarrow/.style={-Stealth, thick, poporange, dashed, bend left=40}
]
\node[box, align=center] (emetteur){\textbf{ÉMETTEUR}\\(Intention, encodage)};
\node[box, right=of emetteur, align=center] (message){\textbf{MESSAGE}\\(Verbal / Non verbal / Chimique)};
\node[box, right=of message, align=center] (recepteur){\textbf{RÉCEPTEUR}\\(Décodage, interprétation)};
\node[box, below=1.5cm of message, align=center] (canal){\textbf{CANAL}\\(Auditif, Visuel, Tactile, Olfactif)};
\node[box, below=1.5cm of emetteur, align=center] (code){\textbf{CODE PARTAGÉ}\\(Langue, Culture, Contexte)};
\node[box, below=1.5cm of recepteur, align=center] (contexte){\textbf{CONTEXTE}\\(Situation, Relation, Culture)};

\draw[arrow] (emetteur) -- node[above, font=\tiny, sloped] {encode} (message);
\draw[arrow] (message) -- node[above, font=\tiny, sloped] {transmet} (recepteur);
\draw[arrow] (canal.north) -- node[right, font=\tiny] {support} (message.south);
\draw[arrow] (code.north east) -- node[above, font=\tiny] {règles} (emetteur.south west);
\draw[arrow] (code.north west) -- node[above, font=\tiny] {règles} (recepteur.south east);
\draw[arrow] (contexte.north) -- node[left, font=\tiny] {influence} (recepteur.south);
\draw[arrow] (contexte.north east) -- node[right, font=\tiny] {influence} (emetteur.south west);

% Rétroaction (feedback)
\draw[backarrow] (recepteur.north west) to node[above, font=\tiny, poporange] {Rétroaction (Feedback)} (emetteur.north east);
\end{tikzpicture}
\legende{Schéma du modèle de la communication interindividuelle (inspiré de Jakobson). La rétroaction (flèche orange pointillée) permet à l'émetteur d'ajuster son message en temps réel.}
\end{popfigure}

\begin{methode}[Analyser une situation de communication]
Pour décrypter toute interaction, suis ces 5 étapes :
\begin{enumerate}
\item \textbf{Identifier l'émetteur et le récepteur :} Qui parle ? À qui ? Quelle est leur relation (symétrique, complémentaire) ?
\item \textbf{Isoler le message :} Quel est le contenu \emph{verbal} (mots) et \emph{non verbal} (gestes, ton, regard) ?
\item \textbf{Préciser le(s) canal(aux) :} Auditif (parole), Visuel (gestes, écrit), Tactile (contact), Chimique (odeurs).
\item \textbf{Repérer le code et le contexte :} Langue commune ? Codes culturels implicites ? Lieu, moment, enjeux sociaux ?
\item \textbf{Analyser la rétroaction :} Le récepteur répond-il (verbal/non verbal) ? L'émetteur ajuste-t-il son discours ?
\end{enumerate}
\end{methode}

\begin{exemplebox}[Analyse d'une scène au lycée]
\textbf{Situation :} Un élève (Émetteur) demande à son professeur (Récepteur) : « Monsieur, j'ai fini le devoir. »\par
\begin{itemize}
\item \textbf{Verbal :} Phrase polie, syntaxe correcte.
\item \textbf{Non verbal :} L'élève évite le regard, voix basse, épaules rentrées, main cachant la feuille.
\item \textbf{Canaux :} Auditif (mots), Visuel (posture, regard).
\item \textbf{Code/Contexte :} Relation autorité/soumission ; devoir difficile ; peur de la note.
\item \textbf{Rétroaction :} Le professeur regarde la feuille, voit des réponses incomplètes, soupire (non verbal) et dit : « Vraiment ? Revois l'exercice 3. » (Verbal + ajustement).
\end{itemize}
\textbf{Interprétation :} Le message verbal ment ; le non verbal révèle la vérité (tricherie/inachèvement). Le professeur décode la contradiction grâce à la congruence canaux/contexte.
\end{exemplebox}

\courssub{La communication verbale : le langage articulé et ses aires cérébrales}

Le langage articulé est la forme la plus sophistiquée de communication verbale. Il repose sur une double articulation (phonèmes $\to$ morphèmes $\to$ phrases) et une latéralisation cérébrale majeure.

\begin{propriete}[Localisation cérébrale du langage (admise, fondée sur l'observation clinique des aphasies)]
Chez plus de 90\% des droitiers et 70\% des gauchers, les fonctions langagières sont \cle{latéralisées dans l'hémisphère gauche}. Deux aires corticales spécialisées, reliées par le \cle{faisceau arqué} (faisceau de substance blanche), jouent des rôles distincts mais complémentaires :
\begin{itemize}
\item \textbf{L'aire de Broca (aires 44 et 45 de Brodmann) :} Située dans le \textbf{gyrus frontal inférieur} (pars opercularis et triangularis). Elle est le centre \textbf{moteur du langage} : programmation des séquences articulatoires, syntaxe, grammaire, production de la parole fluide.
\item \textbf{L'aire de Wernicke (aire 22 de Brodmann) :} Située dans le \textbf{gyrus temporal supérieur} postérieur. Elle est le centre \textbf{sensoriel du langage} : compréhension du langage oral et écrit, sémantique, sélection des mots.
\end{itemize}
Une lésion de l'aire de Broca cause une \cle{aphasie de Broca} (non-fluençante, agrammatique, compréhension préservée). Une lésion de l'aire de Wernicke cause une \cle{aphasie de Wernicke} (fluençante, paraphasique, compréhension altérée).
\end{propriete}

\begin{popfigure}
\begin{tikzpicture}[
    brain/.style={draw=popdark, thick, fill=popblueL, opacity=0.3},
    area/.style={draw=popblue, very thick, fill=popblue!20, rounded corners=3pt},
    label/.style={font=\small, text=popdark, align=center},
    arrow/.style={-Stealth, thick, poporange}
]
% Hémisphère gauche stylisé (vue latérale)
\draw[brain] (0,0) .. controls (0,3) and (4,4) .. (6,3.5) .. controls (7,3) and (7,1) .. (6,0.5) .. controls (5,0) and (3,-0.5) .. (1,-0.2) .. controls (0.5,0) and (0,0.5) .. (0,0);
% Sulcus central (Rolando)
\draw[dashed, popmuted] (3.5,3.2) .. controls (3.8,2) and (3.5,1) .. (3.2,0.2);
% Silvien (Sylvius)
\draw[dashed, popmuted] (2.5,2.5) .. controls (4,1.8) and (5,1.5) .. (5.5,1);

% Aire de Broca (Frontal inférieur)
\draw[area] (1.5,1.8) ellipse (1.2cm and 0.6cm);
\node[label, popblue, align=center] at (1.5, 1.8){\textbf{Aire de Broca}\\(Production, Syntaxe)};

% Aire de Wernicke (Temporal sup postérieur)
\draw[area] (4.2, 1.3) ellipse (1.1cm and 0.5cm);
\node[label, popblue, align=center] at (4.2, 1.3){\textbf{Aire de Wernicke}\\(Compréhension, Sémantique)};

% Faisceau arqué
\draw[arrow, dashed] (2.5, 1.6) .. controls (3.2, 2.2) and (3.8, 2.0) .. (3.5, 1.5);
\node[label, poporange] at (3.3, 2.3) {Faisceau arqué};

% Légendes anatomiques
\node[label, left] at (0.5, 3) {Frontal};
\node[label, above] at (3.5, 3.5) {Pariétal};
\node[label, right] at (6, 1.5) {Temporal};
\node[label, below] at (3, -0.3) {Occipital (arrière)};

% Flèches fonctions
\draw[arrow] (0.2, 1.8) -- node[left, label, align=center]{Motricité\\bouche/langue} (1.0, 1.8);
\draw[arrow] (5.3, 1.3) -- node[right, label, align=center]{Auditif\\Compréhension} (5.8, 1.3);
\end{tikzpicture}
\legende{Carte simplifiée de l'hémisphère gauche (vue latérale) montrant les aires de Broca (frontal) et de Wernicke (temporal) reliées par le faisceau arqué. La latéralisation à gauche est la règle majoritaire.}
\end{popfigure}

\begin{experience}[Preuve clinique : L'aphasie de Broca (Observation historique de Paul Broca, 1861)]
\textbf{Problème :} Le langage est-il localisé ou diffus dans le cerveau ?\par
\textbf{Observation :} Leborgne (« Tan »), patient aphasique ne prononçant que « tan », comprend tout mais ne peut parler. Autopsie : lésion du gyrus frontal inférieur gauche.\par
\textbf{Hypothèse :} Cette zone contrôle la production articulatoire du langage.\par
\textbf{Vérification :} Broca opère 8 cas similaires $\to$ même lésion, même symptôme (aphasie motrice).\par
\textbf{Conclusion :} L'aire de Broca est nécessaire à la \textbf{production} du langage articulé (programmation motrice). C'est la première preuve scientifique de la localisation des fonctions supérieures.
\end{experience}

\begin{exoresolu}[Différencier aphasie de Broca et aphasie de Wernicke]
\textbf{Énoncé :} Deux patients, M. A et M. B, présentent des troubles du langage après un AVC hémisphère gauche. M. A parle peu, phrases télégraphiques (« Moi... partir... maison »), comprend les consignes, est frustré. M. B parle abondamment, flux rapide, phrases grammaticales mais vides de sens (« Le table mange le temps »), ne comprend pas « Levez le bras », n'est pas conscient de son trouble. Localisez les lésions et nommez les aphasies.\par
\tcblower
\textbf{Solution détaillée :}
\begin{center}
\begin{tabularx}{\linewidth}{|l|X|X|}\hline
\textbf{Critère} & \textbf{M. A} & \textbf{M. B} \\ \hline
\textbf{Fluence} & Non fluente (effort, arrêt) & Fluente (logorrhée) \\ \hline
\textbf{Compréhension} & \textbf{Préservée} & \textbf{Altérée} \\ \hline
\textbf{Syntaxe} & Agrammatisme (télégraphique) & Paraphasies (mots wrong), néologismes \\ \hline
\textbf{Conscience du trouble} & Oui (anxiété) & Non (anosognosie) \\ \hline
\textbf{Aire lésée} & \textbf{Aire de Broca} (Frontal inf. G) & \textbf{Aire de Wernicke} (Temporal sup. post.) \\ \hline
\textbf{Type d'aphasie} & \textbf{Aphasie de Broca (motrice)} & \textbf{Aphasie de Wernicke (sensorielle)} \\ \hline
\end{tabularx}
\end{center}
\textbf{Conclusion :} La double dissociation (production vs compréhension) confirme la spécialisation fonctionnelle des deux aires.
\end{exoresolu}

\courssub{La communication non verbale : le corps qui parle}

Si le verbal transmet l'information « froide » (données, logique), le non verbal transmet l'information « chaude » (émotions, attitudes, relation). Il est \cle{multicanal} et souvent \cle{inconscient}.

\begin{definition}[Communication non verbale]
Ensemble des signaux corporels et paralinguistiques accompagnant ou remplaçant la parole : \textbf{kinésique} (gestes, mimiques, postures, regard), \textbf{paraverbal} (ton, débit, volume, timbre, silences), \textbf{proxémique} (distance interpersonnelle), \textbf{haptique} (contact physique), \textbf{apparence} (vêtements, ornements).
\end{definition}

\begin{center}
\begin{tabularx}{\linewidth}{|l|X|X|}\hline
\rowcolor{popblueL}
\textbf{Composante} & \textbf{Éléments} & \textbf{Fonction principale} \\ \hline
\textbf{Kinésique} & Mimiques (sourire, froncement), Gestes (illustrateurs, régulateurs, emblèmes), Posture (ouverte/fermée), Regard (contact/fuite) & Exprimer émotions, réguler tour de parole, illustrer propos \\ \hline
\textbf{Paraverbal} & Intensité, Hauteur, Débit, Timbre, Pauses, Rires, Soupirs & Modifier sens phrase (ironie, doute), signaler état affectif \\ \hline
\textbf{Proxémique} & Distance intime ($<0,45$m), personnelle ($0,45-1,2$m), sociale ($1,2-3,6$m), publique ($>3,6$m) & Définir intimité, statut, agressivité ou séduction \\ \hline
\textbf{Haptique} & Poignée de main, tape épaule, baiser, caresse & Créer lien, apaiser, dominer, réconforter \\ \hline
\end{tabularx}
\end{center}

\begin{attention}[Piège majeur : Réduire la communication aux mots]
\begin{itemize}
\item \textbf{Mythe des « 7\%-38\%-55\% » (Mehrabian) :} Souvent cité à tort comme règle universelle. En réalité, ces pourcentages ne valent \textbf{que pour la communication des sentiments/attitudes quand verbal et non verbal sont incohérents}.
\item \textbf{Règle de congruence :} Si verbal et non verbal sont \textbf{congruents}, le verbal domine l'information factuelle. S'ils sont \textbf{incongruents}, le récepteur croit \textbf{le non verbal} (jugé plus difficile à contrôler volontairement).
\item \textbf{Erreur au Probatoire :} Dire « le non verbal compte pour 93\% ». Faux. Dire « en cas de contradiction, le non verbal l'emporte sur le verbal pour l'interprétation de l'attitude ». Vrai.
\end{itemize}
\end{attention}

\begin{exemplebox}[Contradiction verbal / non verbal : le « Oui » qui veut dire « Non »]
\textbf{Contexte :} Un chef de service demande à son subalterne : « Tu es d'accord avec cette nouvelle procédure ? »\par
\textbf{Verbal :} « Oui, chef, tout à fait d'accord. » (Parole soumise).\par
\textbf{Non verbal :} Sourcil froncé, lèvres pincées, regard baissé, bras croisés, ton monocorde, silence de 3 secondes avant de répondre.\par
\textbf{Analyse :} Incongruence totale. Le non verbal signale \textbf{hostilité, désaccord, soumission forcée}. Le chef perçoit le « vrai » message (le non verbal) et anticipe une application sabotée.
\end{exemplebox}

\begin{savaistu}[Codes culturels camerounais : le corps parle en pays bamiléké et au-delà]
Au Cameroun, la communication non verbale est un langage codifié, variant selon les ethnies et le contexte urbain/rural :
\begin{itemize}
\item \textbf{Le claquement de doigts (\"snapping\") :} Souvent utilisé pour appeler un serveur, un taxi, ou marquer l'accord/la rapidité. En zone bamiléké, un claquement sec peut signifier « Compris » ou « Vite ».
\item \textbf{Le hochement de tête arrière (menton relevé) :} Signifie souvent « Non » ou « Je ne sais pas » (contrairement au hochement avant « Oui » standard). Fréquent dans les Grassfields.
\item \textbf{Le regard :} Baisser les yeux face à un aîné = \textbf{Respect} (pas mensonge ni timidité pathologique). Soutenir le regard = Défi ou égalité.
\item \textbf{La main gauche :} Traditionnellement « impure » (hygiène). Donner/recevoir de la main gauche = \textbf{Insulte} ou manque d'éducation. On tend la main droite, parfois soutenue par la gauche au poignet (marque de grand respect).
\item \textbf{La distance (Proxémique) :} Plus réduite qu'en Europe dans les files d'attente ou transports (contact corporel toléré), mais très codifiée face aux autorités/ainés.
\end{itemize}
\textbf{Compétence interculturelle :} Un soignant ou un enseignant doit connaître ces codes pour ne pas juger « suspect » un élève qui ne le regarde pas, ou « agressif » un patient qui claque des doigts.
\end{savaistu}

\courssub{Communication chimique chez l'Homme : la piste des phéromones}

\begin{definition}[Phéromone (définition étologique stricte)]
Substance chimique sécrétée par un individu, reçue par un congénère de la même espèce, qui déclenche une \textbf{réponse comportementale ou physiologique stéréotypée, innée et involontaire} (ex: attraction sexuelle, alarme, marquage territorial).
\end{definition}

\begin{attention}[Prudence scientifique : L'homme n'est pas la souris]
\begin{itemize}
\item \textbf{Organe voméro-nasal (OVN) / Jacobson :} Chez l'homme, il est \textbf{vestigial, non fonctionnel} (pas de connexion nerveuse vers le cerveau chez l'adulte). La détection chimique se fait via l'\textbf{épithélium olfactif principal}.
\item \textbf{Aucune phéromone humaine identifiée formellement :} Malgré des décennies de recherche, \textbf{aucune molécule} ne remplit \textbf{tous} les critères de la définition étologique (réponse innée, stéréotypée, involontaire, reproductible).
\item \textbf{Preuves d'effets chimiosensoriels (non phéromonaux) :}
\begin{itemize}
\item \textbf{Synchronisation menstruelle (Effet McClintock, 1971) :} Controversée, non reproduite de façon robuste ; biais méthodologiques.
\item \textbf{Reconnaissance odeur mère/nourrisson :} Apprise, pas innée (expérience précoce).
\item \textbf{Préférence MHC (Complexe Majeur d'Histocompatibilité) :} Les femmes préfèrent l'odeur de hommes MHC-dissemblables (diversité immunitaire). Effet modulé par pilule contraceptive. \textbf{C'est une préférence olfactive apprise/évolutive, pas un réflexe phéromonal.}
\item \textbf{Stéroïdes androgènes (androstadiénone) / œstratétraénol :} Modulent l'humeur, l'attention, le cortisol, l'activation hypothalamo-hypophysaire \textbf{selon le contexte et le sexe du receveur}. Effets \textbf{modulateurs} (priming), pas déclencheurs de comportement fixe.
\end{itemize}
\end{itemize}
\textbf{Conclusion pédagogique :} L'homme utilise des \textbf{signaux chimiques} (odeurs corporelles, signature olfactive) pour la communication sociale (reconnaissance, attraction, émotions), mais \textbf{pas de phéromones} au sens éthologique strict. Parlez de \textbf{chimiosignaux} ou \textbf{semiochimiques}.
\end{attention}

\courssub{Troubles de la communication : l'aphasie, fenêtre sur le cerveau}

L'étude des aphasies (perte partielle ou totale du langage suite à lésion cérébrale acquise) est la preuve historique majeure de la localisation fonctionnelle.

\begin{experience}[Démarche d'investigation : Localiser une lésion à partir du symptôme (Examen neurologique)]
\textbf{Situation :} Patient droitier, AVC ischémique territoire artère cérébrale moyenne (ACM) gauche.\par
\textbf{Observation (Examen du langage) :}
\begin{enumerate}
\item \textbf{Parole spontanée :} Fluente, paraphasies phonémiques (« taple » pour « table »), néologismes, logorrhée. $\to$ \textbf{Production intacte}.
\item \textbf{Répétition :} Impossible (« Répétez : "Le chat mange la souris" » $\to$ « Le... gat... mange... »).
\item \textbf{Compréhension :} Altérée (ne pointe pas l'objet « montre » sur commande verbale ; confusion sémantique).
\item \textbf{Nommage :} Anomie sévère (ne trouve pas le mot « montre »).
\item \textbf{Lecture/Écriture :} Altérées (alexie, agraphie).
\end{enumerate}
\textbf{Interprétation :} Syndrome de \textbf{Wernicke} (compréhension + répétition + nommage atteints ; fluence conservée).\par
\textbf{Localisation :} Lésion du gyrus temporal supérieur postérieur gauche (Aire de Wernicke + aires 39/40).\par
\textbf{Confirmation :} IRM cérébrale $\to$ Hypersignal T2/FLAIR dans le territoire temporal postérieur ACM gauche.
\end{experience}

\begin{propriete}[Classification clinique des aphasies (Schéma de Wernicke-Geschwind simplifié)]
\begin{center}
\begin{tabularx}{\linewidth}{|l|X|X|X|}\hline
\rowcolor{popblueL}
\textbf{Type} & \textbf{Zone lésée (Hémisphère G)} & \textbf{Parole} & \textbf{Compréhension} \\ \hline
\textbf{Broca (Motrice)} & Frontal inf. (Aires 44/45) & Non fluente, agrammatique, effort & \textbf{Bonne} \\ \hline
\textbf{Wernicke (Sensorielle)} & Temporal sup. post. (Aire 22) & Fluente, paraphasique, jargon & \textbf{Mauvaise} \\ \hline
\textbf{Conduction} & Faisceau arqué (connexion) & Fluente, \textbf{Répétition impossible} & Bonne \\ \hline
\textbf{Globale} & Frontal + Temporal (vaste ACM) & Mutisme ou stéréotypies & Nulle \\ \hline
\end{tabularx}
\end{center}
\textbf{Note :} La \textbf{répétition} est le test discriminant clé : conservée dans Broca/Wernicke pures, perdue dans Conduction/Globale.
\end{propriete}

\courssub{Communication, cohésion sociale et analyse critique : rumeurs et réseaux sociaux}

La communication n'est pas qu'un échange dyadique (A $\leftrightarrow$ B) ; elle structure les groupes. Les \cle{réseaux sociaux numériques} (Facebook, WhatsApp, TikTok, X) transforment radicalement la dynamique.

\begin{exemplebox}[Analyse critique : La rumeur « Vaccin stérilisant » au Cameroun (2021-2023)]
\textbf{Situation :} Message vocal WhatsApp (émetteur anonyme, « tante infirmière ») : « Le vaccin COVID/HPV rend stérile. Ne laissez pas vacciner vos filles. »\par
\textbf{Analyse par le modèle de la communication :}
\begin{itemize}
\item \textbf{Émetteur :} Anonyme, autorité usurpée (blouse blanche), pas de source vérifiable.
\item \textbf{Message :} Contenu émotionnel fort (peur stérilité), binaire (bon/mauvais), absence de données chiffrées.
\item \textbf{Canal :} WhatsApp (chiffré, viral, hors contrôle éditorial), oral (accessible analphabètes), partage pair-à-pair (confiance sociale).
\item \textbf{Récepteurs :} Parents, leaders d'opinion locaux, influenceurs. Contexte : méfiance historique (essais cliniques passés), faible littératie sanitaire.
\item \textbf{Rétroaction :} Commentaires « Merci tante », transferts massifs, refus vaccination dans les écoles.
\item \textbf{Effet de groupe :} \textbf{Chambre d'écho} (algorithme renforce croyance), \textbf{biais de confirmation} (on croit ce qui arrange), \textbf{poliarisation} (pro/anti vaccin = clans ennemis).
\end{itemize}
\textbf{Conséquence sanitaire :} Chute couverture vaccinale HPV (cancer col utérus) et COVID, résurgence rougeole.
\textbf{Leçon :} La communication numérique désintermédiée (sans journaliste/médecin vérificateur) amplifie l'émotion au détriment de la preuve. L'éducation aux médias (EMI) est une urgence de santé publique.
\end{exemplebox}

\begin{methode}[Analyser un message viral (réseaux sociaux) : Esprit critique]
\begin{enumerate}
\item \textbf{Source :} Qui émet ? Compte vérifié ? Anonyme ? Expertise réelle ? Conflit d'intérêt ?
\item \textbf{Contenu :} Fait vérifiable (chiffres, source primaire) ou Opinion/Émotion ? Preuves citées (liens vers études, rapports officiels) ?
\item \textbf{Forme :} Titre racoleur (clickbait) ? Image hors contexte (reverse image search : TinEye, Google Images) ? Montage vidéo ?
\item \textbf{Contexte :} Date ? Lieu ? Événement d'actualité qui instrumentalise l'info ?
\item \textbf{Corroboration :} L'info est-elle reprise par médias de référence (CRTV, Le Monde, OMS, MinSanté) ? \textbf{Si non $\to$ Ne pas partager.}
\end{enumerate}
\end{methode}

\begin{aretenir}[Synthèse : La communication interindividuelle]
\begin{itemize}
\item C'est un \textbf{processus dynamique, circulaire} (émetteur $\leftrightarrow$ récepteur) et \textbf{multicanal} (verbal, non verbal, chimique).
\item Le \textbf{langage verbal} repose sur des \textbf{aires cérébrales spécialisées et latéralisées} (Broca = production/syntaxe ; Wernicke = compréhension/sémantique ; Faisceau arqué = connexion). Preuve : double dissociation des aphasies.
\item La \textbf{communication non verbale} (kinésique, paraverbal, proxémique) porte l'affect, la relation, la culture. \textbf{En cas de contradiction avec le verbal, c'est elle qui est crue.}
\item Les \textbf{phéromones} n'existent pas chez l'homme au sens strict ; on parle de \textbf{chimiosignaux} (odeurs corporelles) à effets modulateurs, traités par l'olfaction principale.
\item Les \textbf{réseaux sociaux} transforment la communication de masse : viralité, anonymat, chambres d'écho $\to$ risque de désinformation massive. L'\textbf{analyse critique} (source, preuve, corroboration) est une compétence vitale.
\end{itemize}
\end{aretenir}

\begin{exercice}[Entraînement au Probatoire]
\begin{enumerate}
\item Un patient présente une aphasie non fluente, agrammatique, avec compréhension conservée et hémiparésie droite. Nommez le syndrome, l'aire lésée, l'artère probablement responsable.
\item Lors d'un entretien d'embauche, le candidat dit : « Je suis très à l'aise pour parler en public » (verbal), mais il transpire, évite le regard, bafouille, voix tremblante (non verbal). Le recruteur retient-il le verbal ou le non verbal ? Justifiez par le concept de congruence.
\item Citez deux arguments scientifiques contre l'existence de phéromones humaines au sens éthologique strict.
\item Proposez une stratégie de communication pour déconstruire une rumeur sanitaire sur WhatsApp dans un quartier de Douala (émetteur, message, canal, relais de confiance).
\end{enumerate}
\end{exercice}

\begin{corrige}[Exercice n°1]
\begin{enumerate}
\item \textbf{Syndrome :} Aphasie de Broca (motrice). \textbf{Aire :} Aire de Broca (gyrus frontal inf. G, aires 44/45). \textbf{Artère :} Branche supérieure de l'artère cérébrale moyenne (ACM) gauche (artère rolandique / pré-rolandique).
\item \textbf{Non verbal.} Incongruence marquée (verbal = confiance ; non verbal = anxiété intense). Le non verbal est involontaire, difficile à contrôler, perçu comme plus authentique. Le recruteur infère : « Il ment ou manque de confiance réelle ».
\item 1) Absence d'organe voméro-nasal fonctionnel (pas de voie nerveuse dédiée). 2) Aucune molécule n'induit une réponse comportementale stéréotypée, innée, involontaire et reproductible (critères de Karlson \& Lüscher). Les effets observés sont modulatoires, contextuels, appris.
\item Stratégie : \textbf{Émetteurs de confiance} (chef quartier, infirmier chef de poste, leader religieux, influenceur local connu) $\to$ \textbf{Message} court, vidéo locale (langue/culture), fait + émotion positive (protection enfants) + appel à l'action (lieu/heure vaccin) $\to$ \textbf{Canal} WhatsApp (transferts), radio locale, causerie au marché $\to$ \textbf{Rétroaction} : numéro vert questions, témoignage vaccinés du quartier. Ne pas censurer : \textbf{pré-bunker} (inoculation cognitive) en expliquant la tactique du mensonge viral.
\end{enumerate}
\end{corrige}

\coursec{Régulation hormonale des comportements sociaux}

Après avoir étudié les modalités de la communication interindividuelle — langage verbal, mimiques, postures, phéromones — nous comprenons que l'échange d'informations entre individus repose sur des signaux extérieurs. Mais qu'en est-il de la \emph{réponse} de l'organisme à ces signaux ? Pourquoi une même remarque peut-elle déclencher une colère explosive chez l'un et une indifférence polie chez l'autre ? La réponse réside pour une grande part dans notre \textbf{milieu intérieur}, régi par le système hormonal. Les hormones ne « commandent » pas nos actes comme un marionnettiste, mais elles \textbf{modulent la réactivité} de notre cerveau, baissant ou relevant le seuil de déclenchement de certains comportements. Explorons cette chimie subtile qui colore nos vies sociales.

\courssub{Rappel : le système endocrinien, un réseau de messagers chimiques}

\begin{definition}[Hormone, glande endocrine, organe cible]
Une \cle{hormone} est une substance chimique (protéine, peptide, stéroïde ou dérivée d'acide aminé) sécrétée par une \cle{glande endocrine} (glande sans canal excréteur) directement dans le sang. Elle est transportée par la circulation sanguine vers des \cle{organes cibles} (ou cellules cibles) qui possèdent des \textbf{récepteurs spécifiques} à cette hormone. La fixation hormone-récepteur déclenche une cascade de réactions intracellulaires modifiant l'activité de la cellule cible.
\end{definition}

\begin{exemplebox}[Exemple concret : l'insuline, modèle de régulation hormonale]
Après un repas riche en glucides (riz, manioc, banane plantain), la glycémie augmente. Les cellules \cle{bêta} des \cle{îlots de Langerhans} (pancréas, glande endocrine) détectent cette hausse et sécrètent de l'\cle{insuline} dans le sang. L'insuline rejoint le foie, le muscle et le tissu adipeux (organes cibles) via la circulation. Elle se fixe sur ses récepteurs membranaires, déclenchant l'entrée du glucose dans les cellules et sa conversion en glycogène ou en graisses. La glycémie redescend : l'homéostasie est restaurée.
\end{exemplebox}

\begin{attention}[Piège fréquent : confondre glande endocrine et glande exocrine]
Le pancréas est une glande \emph{mixte} : sa partie endocrine (îlots de Langerhans) sécrète insuline et glucagon dans le sang ; sa partie exocrine sécrète le suc pancréatique dans le duodénum via le canal de Wirsung. Seule la partie endocrine nous intéresse ici pour la régulation comportementale.
\end{attention}

\courssub{Les principales hormones modulatrices des comportements sociaux}

\begin{aretenir}[Principe général : modulation, non déterminisme]
Une hormone ne « crée » pas un comportement à elle seule. Elle \textbf{modifie la probabilité et l'intensité} de sa survenue en agissant sur la réactivité des circuits neuronaux (seuil de déclenchement, amplification du signal). Le contexte psychosocial, l'apprentissage, la culture et l'histoire personnelle restent les déterminants majeurs de l'expression finale du comportement.
\end{aretenir}

\textbf{Adrénaline et noradrénaline : les hormones de l'urgence et de l'affrontement.}\par

Secrétées par la \cle{médullosurrénale} (partie centrale des glandes surrénales) sous le contrôle direct du système nerveux \cle{sympathique}, l'\cle{adrénaline} (épinéphrine) et la \cle{noradrénaline} (norépinéphrine) préparent l'organisme à la réaction \textbf{« lutte ou fuite »} (\emph{fight or flight}).

\begin{experience}[Mise en évidence du rôle de l'adrénaline dans la réaction d'urgence]
\textbf{Observation :} Face à un danger soudain (ex. : traversée brutale d'une route à Yaoundé, rencontre avec un animal sauvage près du Mont Cameroun), le cœur s'emballe, la respiration s'accélère, les mains tremblent, la vigilance est maximale.
\textbf{Hypothèse :} Une hormone libérée rapidement dans le sang active ces réponses.
\textbf{Expérience (données classiques) :} Injection d'adrénaline chez l'homme → tachycardie, hypertension, bronchodilatation, hyperglycémie, mydriase, inhibition digestion. Injection de bêta-bloquants (antagonistes) → atténuation de ces signes.
\textbf{Interprétation :} L'adrénaline mobilise l'énergie (glycogène → glucose), optimise l'apport O\textsubscript{2} (cœur, poumons) et recentre l'attention sur la menace.
\textbf{Conclusion :} L'adrénaline est l'hormone clé de la réponse aiguë au stress physique ou psychologique, facilitant les comportements d'agressivité (lutte) ou de fuite.
\end{experience}

\begin{exemplebox}[Contexte camerounais : le stress du conducteur de « bendy »]
Au carrefour « Carrefour Wouri » à Douala, un taxi-man freine brutalement pour éviter un piéton. Son système sympathique décharge massivement de la noradrénaline (neurotransmetteur) et ses surrénales déversent de l'adrénaline dans le sang. En 2 à 3 secondes, son rythme cardiaque passe de 70 à 130 battements/min, ses pupilles se dilatent pour mieux voir, son foie libère du glucose pour alimenter ses muscles. Il a la « pêche » pour klaxonner, crier ou manœuvrer. Sans cette montée hormonale, sa réaction serait trop lente.
\end{exemplebox}

\textbf{Cortisol : l'hormone du stress prolongé.}\par

Si le stress persiste (examens du Probatoire, conflits familiaux chroniques, insécurité alimentaire), l'axe \cle{hypothalamo-hypophyso-surrénalien} (axe HHS ou axe corticotrope) prend le relais. L'hypothalamus sécrète la \cle{CRH} (corticolibérine) → l'hypophyse antérieure sécrète l'\cle{ACTH} (corticotrophine) → la \cle{zone fasciculée} du cortex surrénalien sécrète le \cle{cortisol} (glucocorticoïde principal chez l'homme).

\begin{propriete}[Effets du cortisol sur le comportement (admis)]
Le cortisol maintient l'état d'alerte sur la durée (heures à jours) : maintien de la glycémie, modulation de l'immunité, effet sur l'humeur. \textbf{Stress chronique (cortisol constamment élevé)} → irritabilité, troubles du sommeil, repli social, difficultés de concentration, anxiété, risque dépressif. À l'inverse, un déficit (maladie d'Addison) entraîne fatigue extrême, apathie, hypotension.
\end{propriete}

\begin{exemplebox}[Exemple chiffré : cinétique adrénaline vs cortisol]
\begin{itemize}
    \item \textbf{Adrénaline} : pic plasmatique en \textbf{2--3 minutes} après le stress, demi-vie \textbf{2--3 minutes}. Action brève, intense.
    \item \textbf{Cortisol} : pic plasmatique en \textbf{20--30 minutes}, demi-vie \textbf{60--90 minutes}. Action lente, durable.
\end{itemize}
C'est pourquoi on « tremble encore » une heure après une frayeur (cortisol), alors que le pic de panique (adrénaline) est passé depuis longtemps.
\end{exemplebox}

\textbf{Ocytocine : l'hormone de l'attachement et de la confiance.}\par

Synthétisée dans les noyaux \cle{paraventriculaire} et \cle{supraoptique} de l'\cle{hypothalamus}, stockée et libérée par la \cle{neurohypophyse} (hypophyse postérieure). Elle agit comme hormone (sang) et comme neuromodulateur (cerveau).

\begin{savaistu}[L'ocytocine, « hormone de l'attachement »]
Son taux sanguin augmente significativement lors de \textbf{contacts affectueux} : étreintes, caresses, relations sexuelles (orgasme), allaitement (réflexe d'éjection du lait), accouchement (contractions utérines). Chez la mère, elle facilite le \textbf{lien mère-enfant} dès la naissance. Chez le père, le taux augmente aussi avec l'investissement parental (portage, jeux). Des études montrent qu'une administration intranasale d'ocytocine augmente la \textbf{confiance} envers autrui dans des jeux économiques et améliore la \textbf{reconnaissance des émotions} faciales.
\end{savaistu}

\begin{exemplebox}[Situation concrète : la « kangourou » en néonatologie]
Au Centre Mère-Enfant de la Fondation Chantal Biya (Yaoundé), la méthode « mère kangourou » (contact peau-à-peau prolongé) est pratiquée pour les prématurés. Ce contact stimule la libération d'ocytocine chez la mère \textbf{et} le bébé : réduction du stress (cortisol ↓), stabilisation température/fréquence cardiaque, renforcement du lien d'attachement, meilleure lactation. L'ocytocine est le ciment biologique de l'empathie et de la cohésion sociale.
\end{exemplebox}

\textbf{Stéroïdes gonadiques : testostérone et œstrogènes.}\par

Secrétés principalement par les \cle{testicules} (testostérone) et les \cle{ovaires} (œstradiol, progestérone), sous contrôle de l'axe \cle{hypothalamo-hypophyso-gonadique} (GnRH → LH/FSH). Ils organisent le cerveau \emph{in utero} (organisation) et l'activent à la puberté et chez l'adulte (activation).

\begin{propriete}[Influence comportementale des stéroïdes gonadiques (admis)]
\begin{itemize}
    \item \textbf{Testostérone} : corrélée (relation complexe, non linéaire) à l'\textbf{agressivité} de statut/dominance, à la \textbf{libido}, à la recherche de partenaire (séduction), à la \textbf{compétitivité}. Chez le père, une \textbf{baisse} de testostérone est observée avec l'investissement parental (soins, portage), favorisant la tolérance et la protection.
    \item \textbf{Œstrogènes/Progestérone} : modulent l'humeur, l'anxiété, la réceptivité sexuelle (œstrus chez l'animal, libido cyclique chez la femme), les comportements maternels (préparation du « nid », protection). Fluctuations cycliques → syndrome prémenstruel, humeur postpartum.
\end{itemize}
\end{propriete}

\begin{attention}[Attention : éviter le déterminisme hormonal simpliste !]
\emph{« Il est violent parce qu'il a trop de testostérone »} est une explication \textbf{scientifiquement fausse} et socialement dangereuse.
\begin{itemize}
    \item La testostérone module la \textbf{réactivité} aux provocations, elle ne « programme » pas la violence.
    \item L'éducation, les normes sociales, le statut socio-économique, l'histoire traumatique pèsent infiniment plus lourd.
    \item Chez la femme, les œstrogènes ne « rendent » pas irrationnelle ; ils modulent la sensibilité émotionnelle.
\end{itemize}
La biologie propose des \textbf{probabilités}, la personne \textbf{choisit} (responsabilité).
\end{attention}

\courssub{L'hypothalamus : chef d'orchestre de la neuroendocrinien}

\begin{propriete}[L'hypothalamus relie les systèmes nerveux et hormonal (admis)]
L'\cle{hypothalamus} est la \textbf{clé de voûte} de l'intégration neuroendocrinienne :
\begin{enumerate}
    \item \textbf{Zone nerveuse :} reçoit des informations du système limbique (émotions), du cortex (pensée), des viscères (intéroception), de la rétine (rythme circadien).
    \item \textbf{Zone endocrine :} synthétise les \cle{neurohormones} libératrices (GnRH, CRH, TRH, GHRH, somatostatine, dopamine/PIH) → porte hypophysaire → hypophyse antérieure.
    \item \textbf{Neurosecretion :} neurones neurosecréteurs (noyaux paraventriculaire/supraoptique) → axones → neurohypophyse → libération ocytocine/vasopressine (ADH) dans le sang.
    \item \textbf{Contrôle autonome :} centre supérieur du système nerveux végétatif (sympathique/parasympathique) → contrôle direct surrénales (médullosurrénale).
\end{enumerate}
Il coordonne ainsi la réponse \textbf{rapide} (nerf vague, sympathique, adrénaline) et la réponse \textbf{durable} (axes corticotrope, gonadotrope, thyrotrope, somatotrope) aux stimuli sociaux.
\end{propriete}

\begin{experience}[Démarche d'investigation : l'hypothalamus centre de l'agressivité (expérience classique chez le chat)]
\textbf{Observation :} Stimulation électrique de zones précises de l'hypothalamus → comportements stéréotypés.
\textbf{Expérience (Delgado, Hess, années 1950-60) :}
\begin{itemize}
    \item Stimulation \textbf{hypothalamus latéral} → \textbf{agressivité prédatrice} (chasse silencieuse, pupilles dilatées, poils dressés, attaque proie). Comportement « à froid », orienté vers un but.
    \item Stimulation \textbf{hypothalamus médian (aire septale)} → \textbf{agressivité affective} (rugissements, grincement dents, piloérection, attaque congénaire). Comportement « à chaud », émotionnel, défensif.
\end{itemize}
\textbf{Interprétation :} L'hypothalamus contient des \textbf{centres moteurs organisés} pour des comportements sociaux complexes. Il intègre les signaux hormonaux (testostérone sensibilise l'hypothalamus latéral) et nerveux pour déclencher le pattern moteur adéquat.
\textbf{Conclusion :} L'hypothalamus est l'interface où l'état hormonal (testostérone, cortisol) modifie le seuil de déclenchement des programmes moteurs innés de comportement social.
\end{experience}

\courssub{Comparaison : message nerveux vs message hormonal}

Le tableau suivant synthétise les différences fondamentales entre les deux grands systèmes de communication interne, dont la synergie assure l'adaptation comportementale.

\begin{popfigure}
\begin{tikzpicture}[
    node distance=2mm,
    header/.style={fill=popblue, text=white, font=\bfseries\small, align=center, minimum height=8mm},
    cell/.style={font=\small, align=center, minimum height=7mm, inner xsep=2pt},
    oddrow/.style={fill=popblueL!30},
    evenrow/.style={fill=white},
    col1/.style={fill=popdark!10, font=\bfseries\small, align=left, text width=3.2cm},
    col2/.style={text width=4.5cm},
    col3/.style={text width=4.5cm}
]
% En-têtes
\node[header, minimum width=3.2cm] (h1) at (0,0) {Critère de comparaison};
\node[header, minimum width=4.5cm] (h2) at (3.2,0) {Message nerveux (influx + neurotransmetteur)};
\node[header, minimum width=4.5cm] (h3) at (7.7,0) {Message hormonal (hormone dans le sang)};

% Lignes
\node[cell, oddrow, col1, anchor=north west] at (0,-7.5) {Vitesse de transmission};
\node[cell, oddrow, col2, anchor=north west] at (3.2,-7.5) {Très rapide (ms à s) : influx \SI{1}{m.s^{-1}} à \SI{120}{m.s^{-1}} + synapse \SI{0,5}{ms}};
\node[cell, oddrow, col3, anchor=north west] at (7.7,-7.5) {Lente (min à h) : diffusion sanguine \SI{5}{L.min^{-1}} + fixation récepteur};

\node[cell, evenrow, col1, anchor=north west] at (0,-15) {Durée de l'action};
\node[cell, evenrow, col2, anchor=north west] at (3.2,-15) {Brève (ms à s) : inactivation rapide neurotransmetteur (recapture, enzymes)};
\node[cell, evenrow, col3, anchor=north west] at (7.7,-15) {Longue (min à jours) : demi-vie hormonale, synthèse protéines nouvelles};

\node[cell, oddrow, col1, anchor=north west] at (0,-22.5) {Ciblage et spécificité};
\node[cell, oddrow, col2, anchor=north west] at (3.2,-22.5) {Extrêmement précis : une fibre innerve quelques cellules cibles (unité motrice)};
\node[cell, oddrow, col3, anchor=north west] at (7.7,-22.5) {Diffus : hormone atteint toutes cellules vascularisées ; spécificité par présence récepteur};

\node[cell, evenrow, col1, anchor=north west] at (0,-30) {Support physique};
\node[cell, evenrow, col2, anchor=north west] at (3.2,-30) {Électrochimique : variation potentiel membranaire + diffusion chimique fente synaptique};
\node[cell, evenrow, col3, anchor=north west] at (7.7,-30) {Chimique pur : molécule transportée par le plasma sanguin};

\node[cell, oddrow, col1, anchor=north west] at (0,-37.5) {Type de réponse};
\node[cell, oddrow, col2, anchor=north west] at (3.2,-37.5) {Contraction musculaire, sécrétion exocrine, potentiel d'action suivant};
\node[cell, oddrow, col3, anchor=north west] at (7.7,-37.5) {Modification métabolisme, synthèse protéines, division cellulaire, modulation excitabilité neuronale};

\node[cell, evenrow, col1, anchor=north west] at (0,-45) {Exemple comportemental};
\node[cell, evenrow, col2, anchor=north west] at (3.2,-45) {Retrait main feu (réflexe), clignement, parole, geste volontaire};
\node[cell, evenrow, col3, anchor=north west] at (7.7,-45) {Réaction stress (cortisol), cycle menstruel, croissance, lien d'attachement (ocytocine), métamorphose};
% Lignes de grille
\draw[popline] (0,0) -- (12.2,0);
\draw[popline] (0,-52.5) -- (12.2,-52.5);
\foreach \i in {0,...,6} {
    \draw[popline] (0,-\i*7.5) -- (12.2,-\i*7.5);
}
\draw[popline] (0,0) -- (0,-52.5);
\draw[popline] (3.2,0) -- (3.2,-52.5);
\draw[popline] (7.7,0) -- (7.7,-52.5);
\draw[popline] (12.2,0) -- (12.2,-52.5);
\end{tikzpicture}
\legende{Tableau comparatif : message nerveux vs message hormonal. Le système nerveux assure la réactivité immédiate et localisée ; le système hormonal assure la coordination globale, durable et l'adaptation de fond. Les deux sont couplés au niveau de l'hypothalamus.}
\end{popfigure}

\courssub{Schéma intégré : de la perception sociale à la réponse hormonale}

Le schéma ci-dessous résume la cascade hypothalamo-hypophysaire reliant un stimulus social (ex. : menace, étreinte, séduction) à la sécrétion hormonale et à la modulation comportementale.

\begin{popfigure}
\begin{tikzpicture}[
    scale=0.95,
    every node/.style={font=\small},
    box/.style={rectangle, rounded corners=3pt, draw=popdark, thick, fill=popcream, minimum height=11mm, minimum width=28mm, align=center, inner sep=2pt},
    gland/.style={ellipse, draw=popblue, thick, fill=popblueL!40, minimum height=10mm, minimum width=24mm, align=center},
    hormone/.style={rectangle, rounded corners=2pt, draw=poporange, thick, fill=poporangeL!40, minimum height=8mm, minimum width=22mm, align=center, font=\scriptsize},
    behavior/.style={rectangle, rounded corners=2pt, draw=popgreen, thick, fill=popgreenL!40, minimum height=8mm, minimum width=26mm, align=center, font=\scriptsize},
    arrow/.style={-Stealth, thick, popdark},
    labelarrow/.style={midway, fill=white, inner sep=1pt, font=\scriptsize, text=popdark}
]

% Nœuds principaux
\node[gland, align=center] (hypo) at (0,0){Hypothalamus\\(intégration limbique,\\cortex, intéroception)};
\node[gland, align=center] (ant) at (5.5,1.8){Hypophyse\\antérieure\\(adénohypophyse)};
\node[gland, align=center] (post) at (5.5,-1.8){Neurohypophyse\\(stockage/libération)};
\node[gland, align=center] (sur) at (11,2.5){Surrénales\\(Cortisol, Adrénaline)};
\node[gland, align=center] (gon) at (11,-2.5){Gonades\\(Testostérone, Œstrogènes)};

% Hormones antérieures
\node[hormone] (acth) at (8.5,2.5) {ACTH};
\node[hormone] (lhfh) at (8.5,1.0) {LH / FSH};
\node[hormone] (tsh) at (8.5,-0.5) {TSH / GH};

% Hormones postérieures
\node[hormone] (ocyt) at (8.5,-1.8) {Ocytocine / ADH};

% Comportements
\node[behavior, align=center] (comp1) at (14,3.2){Stress aigu :\\Fuite / Agressivité\\(Adrénaline)};
\node[behavior, align=center] (comp2) at (14,1.8){Stress chronique :\\Irritabilité / Repli\\(Cortisol)};
\node[behavior, align=center] (comp3) at (14,0.2){Séduction / Libido /\\(Testostérone, Œstrogènes)};
\node[behavior, align=center] (comp4) at (14,-1.8){Attachement / Confiance /\\(Ocytocine)};
\node[behavior, align=center] (comp5) at (14,-3.2){Comportement parental /\\(Testostérone↓, Ocytocine↑)};

% Flèches principales
\draw[arrow] (hypo.east) -- node[labelarrow, align=center]{Neurohormones\\libératrices (GnRH, CRH,\\TRH, GHRH...)} (ant.west);
\draw[arrow] (hypo.east) -- node[labelarrow, below] {Axones neurosecréteurs} (post.west);

% Flèches antérieures
\draw[arrow] (ant.east) -- node[labelarrow, above] {Sang} (acth.west);
\draw[arrow] (ant.east) -- node[labelarrow] {Sang} (lhfh.west);
\draw[arrow] (acth.east) -- node[labelarrow] {Sang} (sur.west);
\draw[arrow] (lhfh.east) -- node[labelarrow] {Sang} (gon.west);

% Flèches postérieures
\draw[arrow] (post.east) -- node[labelarrow] {Sang} (ocyt.west);
\draw[arrow] (ocyt.east) -- node[labelarrow, above] {Sang + Diffusion cérébrale} (comp4.west);

% Flèches glandes cibles -> comportements
\draw[arrow] (sur.east) -- node[labelarrow, above] {Sang} (comp1.west);
\draw[arrow] (sur.east) -- node[labelarrow] {Sang} (comp2.west);
\draw[arrow] (gon.east) -- node[labelarrow] {Sang} (comp3.west);
\draw[arrow] (gon.east) -- node[labelarrow, below] {Sang} (comp5.west);

% Retour feedback (simplifié)
\draw[arrow, poppurple, dashed, bend right=20] (sur.south) to node[labelarrow, poppurple, align=center]{Feedback -\\(Cortisol)} (hypo.south east);
\draw[arrow, poppurple, dashed, bend left=20] (gon.north) to node[labelarrow, poppurple, align=center]{Feedback -\\(Stéroïdes)} (hypo.north east);

% Système nerveux direct (sympathique -> surrénale)
\draw[arrow, poppink, densely dashed] (hypo) -- ++(0,-1.2) -| node[labelarrow, poppink, near start, align=center]{Système nerveux\\sympathique (rapide)} (sur.south west);

% Légendes couleurs
\node[anchor=west, font=\scriptsize] at (0,-4.0) {\textcolor{popdark}{\textbf{---}} Voie hormonale classique (lente, durable)};
\node[anchor=west, font=\scriptsize] at (0,-4.5) {\textcolor{poppink}{\textbf{--- ---}} Voie nerveuse directe sympathique (rapide, adrénaline)};
\node[anchor=west, font=\scriptsize] at (0,-5.0) {\textcolor{poppurple}{\textbf{--- ---}} Rétrocontrôle négatif (homéostasie)};
\end{tikzpicture}
\legende{Schéma de la régulation neuroendocrinienne des comportements sociaux. L'hypothalamus intègre les signaux sociaux (perception, émotion) et orchestre la réponse via l'hypophyse (antérieure : axes corticotrope/gonadotrope ; postérieure : ocytocine) et le système nerveux sympathique (adrénaline). Les hormones ciblent le cerveau (modulation seuil réactivité) et la périphérie (métabolisme, muscle).}
\end{popfigure}

\courssub{Synthèse : une symphonie à deux chefs d'orchestre}

\begin{aretenir}[Points clés à retenir pour le Probatoire]
\begin{enumerate}
    \item Le comportement social résulte de l'interaction dynamique \textbf{Système Nerveux (rapide, ponctuel) + Système Hormonal (lent, diffus, durable)}.
    \item L'\textbf{hypothalamus} est l'interface centrale : il traduit l'état émotionnel/social en signaux hormonaux (axes corticotrope, gonadotrope, thyrotrope, somatotrope) et nerveux (sympathique/parasympathique).
    \item \textbf{Adrénaline/Noradrénaline} : urgence, lutte/fuite (secondes). \textbf{Cortisol} : stress prolongé, vigilance, coût métabolique/psychique si chronique (heures/jours).
    \item \textbf{Ocytocine} : prosocialité, attachement, confiance, lien parental (minutes/heures).
    \item \textbf{Testostérone / Œstrogènes} : modulation de l'agressivité, sexualité, parentalité, compétitivité (jours/mois/années).
    \item \textbf{Jamais de déterminisme hormonal strict :} les hormones modulent des \textbf{probabilités}, l'individu garde sa responsabilité (savoir-être : respect de la diversité, non-stigmatisation).
\end{enumerate}
\end{aretenir}

\begin{exoresolu}[Analyse d'une situation : « Le stress de l'examen du Probatoire »]
\textbf{Énoncé :} Paul, élève de Première A, prépare le Probatoire. Depuis trois semaines, il dort mal (insomnie à 3h du matin), est irritable avec sa famille, a des maux de ventre, et révise moins efficacement malgré des heures passées sur ses livres. Expliquez cette situation à la lumière de la régulation neuroendocrinienne.

\tcblower

\textbf{Solution détaillée :}
\begin{enumerate}
    \item \textbf{Stimulus stressogène chronique :} L'échéance de l'examen (représentation mentale corticale) + charge de travail + enjeu social (réussite, fierté familiale) = activation persistante du système limbique (amygdale, hippocampe) vers l'hypothalamus.
    \item \textbf{Activation de l'axe HHS (Cortisol) :} Hypothalamus $\xrightarrow{\text{CRH}}$ Hypophyse antérieure $\xrightarrow{\text{ACTH}}$ Surrénales (zone fasciculée) $\xrightarrow{\text{Cortisol}}$ Sang.
    \item \textbf{Effets du cortisol chronique (jours/semaines) :}
    \begin{itemize}
        \item \textbf{Métaboliques :} Néo-glucogenèse hépatique → hyperglycémie matinale (réveil 3h, élévation nocturne anormale du cortisol liée au stress + pic circadien anticipé) → faim, prise de poids abdominale.
        \item \textbf{Immunitaires :} Immunosuppression → vulnérabilité infections (rhumes, maux de ventre fonctionnels).
        \item \textbf{Cérébraux :} Hippocampe (récepteurs glucocorticoïdes saturés) → troubles de la mémoire déclarative (révisions inefficaces), de la concentration. Amygdale sensibilisée → \textbf{irritabilité}, anxiété, repli social (conflits familiaux).
        \item \textbf{Sommeil :} Perturbation rythme nycthéméral cortisol (perte de l'amplitude, élévation nocturne) → insomnie de maintien.
    \end{itemize}
    \item \textbf{Pics aigus (Adrénaline) :} À chaque pensée angoissante (« si je rate ? »), système sympathique → médullosurrénale → adrénaline → tachycardie, sueurs, blocage cognitif momentané.
    \item \textbf{Régulation défaillante :} Rétrocontrôle négatif cortisol sur hypothalamus/hypophyse moins efficace (résistance aux glucocorticoïdes) → emballement du cercle vicieux.
    \item \textbf{Pistes de régulation (savoir-être/savoir-faire) :} Activité physique (métabolise cortisol, libère endorphines), cohérence cardiaque / respiration (active parasympathique, freine sympathique), sommeil protégé (hygiène), soutien social ocytocinergique (parler à un proche, câlin), planification réaliste (réduit incertitude = réduit activation amygdale).
\end{enumerate}
\end{exoresolu}

\begin{exercice}[Entraînement : Le père « kangourou »]
Marc, jeune père à Bafoussam, pratique le portage de son nourrisson de 2 mois en écharpe (peau-à-peau) plusieurs heures par jour. On observe chez lui une grande patience, une sensibilité accrue aux pleurs de l'enfant, et une baisse de son intérêt pour les sorties entre amis.
\begin{enumerate}
    \item Identifiez l'hormone principale impliquée dans ce changement comportemental et son site de synthèse/libération.
    \item Expliquez le mécanisme neuroendocrinien reliant le stimulus sensoriel (contact peau-à-peau) à la libération de cette hormone.
    \item Citez une autre hormone dont le taux diminue généralement chez le père investi, et reliez cette baisse au comportement observé (moins de sorties, plus de soin).
    \item En vous appuyant sur la notion de \emph{modulation probabiliste}, expliquez pourquoi Marc pourrait néanmoins, un soir de grande fatigue, crier sur son bébé malgré son taux élevé d'ocytocine.
\end{enumerate}
\end{exercice}

\begin{corrige}[Exercice n°1 : Le père « kangourou »]
\begin{enumerate}
    \item \textbf{Ocytocine.} Synthétisée dans les noyaux paraventriculaire et supraoptique de l'\textbf{hypothalamus}, transportée le long des axones jusqu'à la \textbf{neurohypophyse} (hypophyse postérieure) d'où elle est libérée dans le sang (et dans le cerveau par projections collatérales).
    \item \textbf{Mécanisme :} Stimulation des \textbf{mécanorécepteurs} cutanés (pression, chaleur, caresse) lors du contact peau-à-peau $\rightarrow$ influx nerveux afférents (voies spinothalamiques/lemniscales) $\rightarrow$ hypothalamus (noyaux PVN/SON) $\rightarrow$ décharge des neurones neurosecréteurs $\rightarrow$ libération ocytocine dans la circulation sanguine et dans les aires limbiques (amygdale, noyau accumbens, cortex préfrontal).
    \item \textbf{Testostérone.} Des études longitudinales montrent une \textbf{baisse de la testostérone} chez les pères très investis dans les soins directs (portage, change, nuit). Cette baisse est corrélée à une \textbf{diminution des comportements de recherche de partenaire} (sorties, compétition, séduction) et une \textbf{augmentation de la tolérance, de la vigilance protectrice et de l'empathie} parentale. La testostérone favorise l'effort reproductif (mating effort) ; sa baisse favorise l'effort parental (parenting effort).
    \item \textbf{Modulation probabiliste :} L'ocytocine \textbf{abaisse le seuil} de déclenchement des comportements de soin et \textbf{augmente la récompense} liée à l'interaction avec le bébé. Elle ne \emph{supprime pas} les circuits de la colère/frustration (amygdale, cortex préfrontal). En cas de \textbf{fatigue extrême} (manque de sommeil $\rightarrow$ dysfonction cortex préfrontal = perte du contrôle inhibiteur) + \textbf{stress aigu} (cri strident bébé $\rightarrow$ activation amygdale $\rightarrow$ noradrénaline/adrénaline), la probabilité d'une réaction aggressive (crier) redevient non-nulle, malgré l'ocytocine. L'hormone module, elle n'efface pas la liberté/la responsabilité.
\end{enumerate}
\end{corrige}

\coursec{Diversité des comportements humains et respect d'autrui}

Après avoir étudié comment les hormones modulent nos comportements sociaux (ocytocine, testostérone, cortisol), nous comprenons que la biologie fournit un \emph{support}, mais ne détermine pas tout. Pourquoi, face à la même situation — une remarque d'un enseignant, une provocation dans la cour, un échec scolaire —, deux élèves réagissent-ils si différemment ? L'un sourit, l'autre explose de colère ; l'un cherche le dialogue, l'autre s'isole. Cette diversité n'est pas du hasard : elle naît de l'interaction constante entre notre biologie singulière et notre histoire psychosociale. Comprendre cette interaction, c'est acquérir les clés pour \textbf{maîtriser ses propres comportements} et \textbf{respecter ceux d'autrui}, fondement du vivre-ensemble.

\courssub{Facteurs biologiques et psychosociaux de la diversité}

\textbf{Observation initiale.}
Dans une classe de Première A au Lycée de Nkongsamba, M.~Kamga fait une remarque sur un devoir mal rendu : « Ce travail est en dessous de vos capacités. » Élève A (16 ans, garçon) rougeoit, claque la porte et crie : « C'est n'importe quoi ! » Élève B (17 ans, fille) baisse les yeux, murmure « Je ferai mieux » et rentre chez elle en pleurant. Même stimulus, deux réponses radicalement opposées.

\textbf{Hypothèse.}
Le comportement observable résulte de la confluence de déterminants \textbf{biologiques} (génétique, maturation cérébrale, état hormonal) et de déterminants \textbf{psychosociaux} (éducation, culture, expériences vécues, contexte immédiat).

\textbf{Analyse des facteurs biologiques.}
Chaque individu possède un \cle{patrimoine génétique} unique (sauf jumeaux monozygotes) qui influence son tempérament de base (timidité, impulsivité, seuil de réactivité émotionnelle). Le \cle{sexe} (chromosomes, hormones gonadiques) module certaines réponses : la testostérone favorise l'approche et la compétition, les œstrogènes la sensibilité sociale. L'\cle{âge} et la \cle{maturation cérébrale} sont cruciaux : chez l'adolescent, le \cle{cortex préfrontal} (siège du contrôle inhibiteur, de la planification, de l'empathie cognitive) n'est pas encore mature, tandis que l'\cle{amygdale} (détection de la menace, peur, colère) est hyper-réactive. L'\cle{état hormonal} du moment (cortisol élevé = stress, fatigue) abaisse le seuil de déclenchement des réactions émotionnelles. Enfin, certaines \cle{pathologies} (troubles du spectre autistique, TDAH, dépression, lésions cérébrales) altèrent spécifiquement le traitement de l'information sociale.

\begin{propriete}[Maturation du cortex préfrontal et impulsivité adolescente -- Admise]
Les études de neuro-imagerie (IRMf longitudinales) montrent que le cortex préfrontal, zone clé de la régulation émotionnelle, de la prise de décision et de l'inhibition des comportements inadaptés, n'achève sa myélinisation et son élagage synaptique qu'autour de \textbf{20--25 ans}. Chez l'adolescent, le décalage entre une amygdale mature et réactive et un cortex préfrontal encore immature explique en partie l'impulsivité, la prise de risque et la labilité émotionnelle observées à cet âge.
\end{propriete}

\textbf{Analyse des facteurs psychosociaux.}
L'\cle{éducation} reçue (style parental : autoritaire, démocratique, permissif, négligent) façonne les stratégies de coping. La \cle{culture} et la \cle{religion} définissent les normes d'expression des émotions (au Cameroun, dans certaines cultures, un garçon ne doit pas pleurer en public ; dans d'autres, la colère s'exprime fort). Les \cle{expériences vécues} (traumatismes, réussites, harcèlement antérieur) créent des « filtres » cognitifs : un élève moqué souvent interprétera une critique neutre comme une attaque. L'\cle{environnement familial} (climat, communication, violence conjugale) et le \cle{groupe de pairs} (normes du groupe, pression de conformité, recherche de statut) modèlent les réponses comportementales.

\textbf{Interaction inné / acquis.}
Il n'y a pas de « part d'inné » et de « part d'acquis » qui s'additionnent ; le comportement \emph{émerge} de leur interaction dynamique. Un génotype « à risque » (ex. allèle court du transporteur de la sérotonine \emph{5-HTTLPR}) ne conduit à la dépression ou à l'agressivité \emph{que} s'il est couplé à un environnement adverse (maltraitance, stress chronique) : c'est la \cle{diathèse-stress} ou la \cle{sensibilité différentielle} (certains génotypes confèrent une plus grande plasticité, pour le meilleur comme pour le pire). L'épigénétique (méthylation de l'ADN, modifications des histones) est le mécanisme moléculaire par lequel l'environnement « parle » au génome et module l'expression des gènes impliqués dans la réactivité au stress.

\begin{popfigure}
\begin{tikzpicture}[scale=0.9, every node/.style={font=\small}]
  % Styles
  \tikzset{
    bio/.style={ellipse, draw=popblue, thick, fill=popblue!10, minimum width=4.5cm, minimum height=3.5cm},
    psycho/.style={ellipse, draw=poporange, thick, fill=poporange!10, minimum width=4.5cm, minimum height=3.5cm},
    inter/.style={ellipse, draw=popgreen, thick, fill=popgreen!20, minimum width=3.5cm, minimum height=2.5cm},
    label/.style={font=\bfseries\small, text=popdark},
    item/.style={font=\small, text=popdark, align=left}
  }
  % Cercles
  \node[bio] (B) at (-2.5,0) {};
  \node[psycho] (P) at (2.5,0) {};
  \node[inter] (I) at (0,0) {};
  % Étiquettes des ensembles
  \node[label] at (-2.5, 2.0) {Facteurs \textbf{Biologiques}};
  \node[label] at (2.5, 2.0) {Facteurs \textbf{Psychosociaux}};
  \node[label, text=popgreen!70!black] at (0, -2.2) {\textbf{COMPORTEMENT OBSERVÉ}};
  % Contenu Biologique
  \node[item, anchor=north west, align=center] at (-4.5, 1.5){
    \textbullet\ Patrimoine génétique (tempérament)\\
    \textbullet\ Sexe / Hormones gonadiques\\
    \textbullet\ Âge / Maturation cortex préfrontal\\
    \textbullet\ État hormonal instantané (cortisol)\\
    \textbullet\ Pathologies neurologiques / psychiatriques
  };
  % Contenu Psychosocial
  \node[item, anchor=north east, align=center] at (4.5, 1.5){
    \textbullet\ Éducation / Style parental\\
    \textbullet\ Culture / Religion / Normes sociales\\
    \textbullet\ Expériences vécues (trauma, succès)\\
    \textbullet\ Environnement familial\\
    \textbullet\ Groupe de pairs / Pression sociale
  };
  % Contenu Intersection
  \node[item, align=center] at (0,0) {
    \textbf{Interaction dynamique}\\
    (Épigénétique, Plasticité cérébrale,\\
    Diathèse-stress, Apprentissage)
  };
  % Flèches d'interaction
  \draw[->, thick, popblue] (-1.0, -1.2) -- (-0.5, -0.8) node[midway, above, sloped, font=\tiny] {Module};
  \draw[->, thick, poporange] (1.0, -1.2) -- (0.5, -0.8) node[midway, above, sloped, font=\tiny] {Façonne};
\end{tikzpicture}
\legende{Diagramme de Venn : le comportement observable résulte de l'intersection et de l'interaction dynamique entre les facteurs biologiques (bleu) et les facteurs psychosociaux (orange). L'intersection (vert) représente les mécanismes d'intégration : épigénétique, plasticité neuronale, apprentissage.}
\end{popfigure}

\begin{exemplebox}[Deux élèves, une remarque, deux réactions : analyse factorielle]
\textbf{Situation :} M.~Kamga dit : « Ce travail est en dessous de vos capacités. »

\textbf{Élève A (Réaction explosive / Fuite) :}
\begin{itemize}
  \item \textbf{Biologique :} Garçon, 16 ans, taux élevé de testostérone ; cortex préfrontal immature (faible inhibition) ; amygdale hyper-réactive ; peut-être variant génétique MAOA « low activity » (associé à réactivité accrue à la provocation si maltraitance antérieure) ; cortisol élevé (nuit blanche révision).
  \item \textbf{Psychosocial :} Père autoritaire, critiques humiliantes fréquentes à la maison → hypersensibilité à l'autorité perçue comme injuste ; pairs valorisant la « dureté » et le refus de se laisser faire ; culture locale : « Un homme ne s'excuse pas. »
\end{itemize}
\textbf{Élève B (Réaction inhibition / Repli) :}
\begin{itemize}
  \item \textbf{Biologique :} Fille, 17 ans, phase prémenstruelle (variations œstrogènes/progestérone, sensibilité accrue au stress) ; cortex préfrontal légèrement plus mature (filles en avance ~1-2 ans) ; tempérament anxieux (génétique) ; fatigue.
  \item \textbf{Psychosocial :} Famille valorisant l'effort et le respect du maître ; intériorisation de l'échec comme faute personnelle ; pairs soutenant mais elle ne veut pas les inquiéter ; religion prônant l'humilité et la patience.
\end{itemize}
\textbf{Conclusion :} Le stimulus identique active des réseaux neuronaux différents (amygdale + faible contrôle préfrontal chez A ; amygdale + contrôle préfrontal excessif menant à l'inhibition chez B) \emph{parce que} l'histoire de vie a câblé différemment l'évaluation de la menace (critique = agression vs critique = vérité).
\end{exemplebox}

\courssub{Comportements déviants ou à risque : violence et addictions}

Lorsque l'équilibre bascule vers une réactivité émotionnelle extrême non régulée, ou vers une recherche compulsive de récompense, on bascule dans les \cle{comportements déviants} (transgression des normes sociales) ou \cle{à risque} (mise en danger de soi ou d'autrui).

\textbf{La violence : le court-circuit amygdale -- cortex préfrontal.}
Face à une provocation (insulte, bousculade), l'information visuelle/auditive prend deux voies :
\begin{enumerate}
  \item \textbf{Voie basse (rapide, sous-corticale) :} Thalamus $\rightarrow$ Amygdale $\rightarrow$ Réaction motrice (frapper, fuir) en $\sim$\SI{12}{ms}. C'est la \cle{réaction impulsive}, adaptée pour la survie immédiate (prédateur), inadaptée en société.
  \item \textbf{Voie haute (lente, corticale) :} Thalamus $\rightarrow$ Cortex sensoriel $\rightarrow$ Cortex préfrontal (évaluation : « Est-ce une vraie menace ? Quelle conséquence ? ») $\rightarrow$ Amygdale (modulation) $\rightarrow$ \cle{Réponse maîtrisée} en $\sim$\SI{300}{ms}--\SI{500}{ms}.
\end{enumerate}
Chez l'adolescent violent (ou l'adulte à cortex préfrontal lésé/dysfonctionnel), la voie haute est trop lente ou trop faible : l'amygdale « tire » avant que le cortex ne dise « stop ». L'alcool et les drogues potentialisent ce court-circuit en inhibant le cortex préfrontal (désinhibition).

\textbf{Les addictions : détournement du circuit de la récompense.}
Les drogues (cannabis, tramadol — fléau au Cameroun —, alcool, cocaïne, nicotine) et les addictions comportementales (jeux vidéo, réseaux sociaux) inondent le \cle{noyau accumbens} (centre de la récompense) de \cle{dopamine}, bien au-delà des récompenses naturelles (manger, lien social, réussite). Cela produit :
\begin{itemize}
  \item \textbf{Tolérance :} Récepteurs D2 down-régulés $\rightarrow$ besoin de doses croissantes.
  \item \textbf{Sevrage :} Chute dopaminergique $\rightarrow$ dysphorie, anxiété, \emph{craving} (envie irrésistible).
  \item \textbf{Neuroplasticité maladaptative :} Renforcement des circuits « drogue $\rightarrow$ soulagement » au détriment des circuits « effort $\rightarrow$ récompense différée » (cortex préfrontal atrophié).
\end{itemize}
Conséquences sur les relations : mensonge, vol, rupture familiale, isolement, violence pour se procurer le produit.

\begin{popfigure}
\begin{tikzpicture}[scale=0.85, every node/.style={font=\small},
  decision/.style={diamond, draw=popdark, thick, fill=poporange!20, aspect=2, inner sep=4pt},
  action/.style={rectangle, rounded corners, draw=popblue, thick, fill=popblue!10, minimum width=3cm, minimum height=1cm},
  outcome/.style={rectangle, rounded corners, draw=popgreen, thick, fill=popgreen!10, minimum width=3.5cm, minimum height=1cm},
  arrow/.style={->, thick, popdark},
  label/.style={font=\scriptsize, text=popdark, midway, fill=white, inner sep=1pt}
]
  % Nœuds
  \node[action, align=center] (stim) at (0, 3.5){Stimulus : Provocation\\ (insulte, bousculade, injustice)};
  \node[decision, align=center] (eval) at (0, 1.8){Évaluation rapide\\ (Amygdale : Menace ?)};
  % Voie basse
  \node[action, align=center] (impulsive) at (-4, 0){Réaction IMPULSIVE\\ (Voie basse : Amygdale $\rightarrow$ Corps strié)};
  \node[outcome, align=center] (cons1) at (-4, -1.8){Conséquences : Violence,\\ regret, sanction, escalade};
  % Voie haute
  \node[action, align=center] (prefrontal) at (4, 0){Analyse \& Inhibition\\ (Voie haute : Cortex Préfrontal)};
  \node[outcome, align=center] (cons2) at (4, -1.8){Réponse MAÎTRISÉE :\\ Dialogue, recul, aide, CNV};
  % Flèches
  \draw[arrow] (stim) -- (eval);
  \draw[arrow] (eval) -- node[label, left, align=center]{NON (pas de menace\\ ou contrôle OK)} (prefrontal);
  \draw[arrow] (eval) -- node[label, right, align=center]{OUI (Menace forte\\ + Contrôle faible)} (impulsive);
  \draw[arrow] (impulsive) -- (cons1);
  \draw[arrow] (prefrontal) -- (cons2);
  % Annotation cortex préfrontal
  \node[draw=poppurple, thick, rounded corners, fill=poppurple!10, fit=(prefrontal)(cons2), inner sep=5pt, label={[font=\tiny\bfseries, poppurple]above:Cortex Préfrontal (Frein)}] {};
  % Annotation amygdale
  \node[draw=poppink, thick, rounded corners, fill=poppink!10, fit=(impulsive)(cons1), inner sep=5pt, label={[font=\tiny\bfseries, poppink]above:Amygdale (Accélérateur)}] {};
\end{tikzpicture}
\legende{Arbre de décision face à une provocation. La voie basse (amygdale) mène à l'impulsivité ; la voie haute (cortex préfrontal) permet la maîtrise. L'entraînement renforce la voie haute.}
\end{popfigure}

\begin{attention}[La diversité n'est pas une anomalie]
Ne confondez jamais \textbf{diversité des comportements} et \textbf{pathologie}. Réagir différemment face au même stimulus est la norme biologique (variabilité génétique, épigénétique, historique). Un élève « trop timide », « trop colérique », « trop distrait » n'est pas « anormal » : son comportement a des causes identifiables (biologiques et/ou psychosociales). Le piège est le \textbf{jugement moralisateur} (« il est méchant », « elle est faible ») qui empêche la compréhension et l'aide. La rigueur scientifique appelle à l'\textbf{analyse causale}, non au rejet. Au Probatoire, on attend une explication physiologique et psychosociale, pas un avis personnel.
\end{attention}

\courssub{Maîtrise des comportements : stratégies de régulation}

La maîtrise n'est pas la suppression de l'émotion (impossible et inadapté), mais la \textbf{régulation} : transformer la réaction automatique en réponse choisie. Cela s'apprend, comme un muscle, grâce à la \cle{plasticité cérébrale}.

\begin{methode}[Analyser une situation de vie courante en 4 questions (Savoir-faire)]
Pour expliquer un comportement interindividuel ou proposer une réponse maîtrisée, suivez cette grille :
\begin{enumerate}
  \item \textbf{Quel est le stimulus déclencheur ?} (Description factuelle : qui, quoi, où, quand, ton utilisé).
  \item \textbf{Quel est le support physiologique activé ?} (Amygdale ? Cortex préfrontal ? Axis HPA/cortisol ? Système dopaminergique ? Ocytocine ? Testostérone ?). Identifier la voie (basse/haute).
  \item \textbf{Quelle émotion primaire émerge ?} (Nommer : colère, peur, tristesse, honte, dégoût, surprise, joie). \emph{Nommer apaise l'amygdale (effet « name it to tame it »)}.
  \item \textbf{Quelle réponse maîtrisée est possible ?} (Stratégie : respiration diaphragmatique 4-7-8 pour baisser le cortisol ; reformulation / Communication Non Violente (CNV) : « Quand j'entends X, je ressens Y, car j'ai besoin de Z, je demande... » ; recul cognitif : « Qu'est-ce que cette situation active en moi ? » ; recherche d'aide : pair, adulte de confiance, professionnel).
\end{enumerate}
\end{methode}

\textbf{Stratégies concrètes de renforcement du contrôle préfrontal.}
\begin{itemize}
  \item \textbf{Respiration cohérente / cohérence cardiaque :} \SI{6}{cycles/min} (inspirer \SI{5}{s}, expirer \SI{5}{s}) $\rightarrow$ stimule le nerf vague $\rightarrow$ frein parasympathique sur le cœur $\rightarrow$ signal de sécurité au cerveau $\rightarrow$ cortex préfrontal réactivé.
  \item \textbf{Communication Non Violente (CNV - Rosenberg) :} Observation factuelle $\neq$ Jugement ; Sentiment $\neq$ Pensée ; Besoin $\neq$ Stratégie ; Demande claire $\neq$ Exigence. Exemple : « Tu ne m'écoutes jamais ! » (Jugement/Exigence) $\rightarrow$ « Quand tu regardes ton téléphone pendant que je te parle (Obs), je me sens ignoré (Sentiment), j'ai besoin de considération (Besoin), accepterais-tu de me regarder \SI{2}{min} ? (Demande) ».
  \item \textbf{Empathie cognitive (Théorie de l'Esprit) :} S'efforcer de modéliser l'état mental d'autrui : « Pourquoi a-t-il réagi ainsi ? Quelle peur, quel besoin non satisfait ? » Active le cortex préfrontal médian et la jonction temporo-pariétale.
  \item \textbf{Activité physique régulière :} Augmente le BDNF (facteur neurotrophique), favorise la neurogenèse hippocampique, régule le cortisol, améliore les fonctions exécutives.
  \item \textbf{Sommeil suffisant (\SI{8}{h}--\SI{9}{h} ado) :} Le sommeil nettoie les métabolites (système glymphatique) et consolide les circuits de régulation émotionnelle.
\end{itemize}

\begin{savaistu}[L'entraînement renforce le frein préfrontal]
Des études en IRMf (ex. \emph{Hölzel et al., 2011}) montrent que \SI{8}{semaines} de méditation de pleine conscience (mindfulness) ou de pratique régulière de la régulation émotionnelle augmentent la densité de matière grise dans le \textbf{cortex préfrontal dorsolatéral} et l'\textbf{hippocampe}, et \textbf{diminuent le volume de l'amygdale}. Concrètement : le « muscle » du self-control grossit, le « détecteur de menace » se calme. C'est la neuroplasticité dirigée par l'intention : \emph{vous changez votre cerveau en changeant vos habitudes mentales}.
\end{savaistu}

\courssub{Vers une attitude de respect : accepter la différence et promouvoir le vivre-ensemble}

\textbf{Du savoir au savoir-être.}
Comprendre les bases physiologiques et psychosociales de la diversité transforme le regard que l'on pose sur autrui :
\begin{itemize}
  \item \textbf{Accepter la différence :} L'autre n'est pas « bizarre » ou « mauvais », il est \emph{autre} : son cerveau, son histoire, ses hormones, sa culture ont câblé des réponses différentes des miennes face au même monde.
  \item \textbf{Lutter contre les préjugés et stéréotypes :} « Les Bamileke sont... », « Les filles sont... », « Les musulmans sont... », « Les drogués sont... » sont des \cle{raisonnements heuristiques} (raccourcis cognitifs économes en énergie préfrontale) mais \emph{scientifiquement faux} (généralisation abusive, ignorance de la variance intra-groupe $>$ variance inter-groupe) et \emph{socialement destructeurs} (discrimination, exclusion, violence symbolique).
  \item \textbf{Promouvoir le vivre-ensemble (Cameroun : « Unité dans la diversité ») :} Cela passe par la \cle{justice restaurative} à l'école (médiation par les pairs au lieu de l'exclusion systématique), l'\cle{éducation à la citoyenneté} (droits de l'homme, genre, handicap), et la \cle{valorisation des compétences psychosociales} (résolution de conflits, écoute active) dans les programmes.
\end{itemize}

\begin{aretenir}[Synthèse : Diversité, Maîtrise, Respect]
\begin{itemize}
  \item Le comportement = Interaction dynamique \textbf{Biologique (gènes, hormones, cerveau immature)} $\times$ \textbf{Psychosocial (éducation, culture, pairs, trauma)}.
  \item L'adolescence est une fenêtre de \textbf{vulnérabilité} (amygdale $>$ cortex préfrontal) mais aussi de \textbf{plasticité maximale} : c'est le moment d'apprendre à réguler.
  \item Violence et addictions = \textbf{Déséquilibre} voie basse / voie haute ; circuit de la récompense détourné.
  \item \textbf{Maîtrise} = Entraînement du cortex préfrontal (respiration, CNV, empathie, sommeil, sport) $\rightarrow$ Neuroplasticité positive.
  \item \textbf{Respect} = Reconnaissance scientifique de la causalité complexe du comportement $\rightarrow$ Refus du jugement essentialiste $\rightarrow$ Engagement pour l'inclusion et la justice sociale.
\end{itemize}
\end{aretenir}

\begin{exoresolu}[Application de la méthode d'analyse à une situation de harcèlement]
\textbf{Énoncé :} Kevin (15 ans) subit des moqueries quotidiennes sur son poids de la part d'un groupe de garçons dans la cour. Un jour, l'un d'eux lui lance : « Encore un gros qui prend la place de deux ! » Kevin, d'ordinaire calme, se rue sur lui et le frappe au visage. Il est convoqué chez le proviseur. Analysez la situation avec la méthode en 4 questions pour Kevin, puis proposez une issue éducative (pas seulement punitive).

\textbf{Solution détaillée :}
\begin{enumerate}
  \item \textbf{Stimulus :} Moqueries répétées (harcèlement) + insulte publique ciblant l'image corporelle (« gros », « place de deux ») devant témoins (pairs, filles). Atteinte à l'intégrité physique et à la dignité sociale.
  \item \textbf{Support physiologique :}
      \begin{itemize}
        \item \textbf{Amygdale :} Hyper-activation par la menace sociale répétée (système de détection de l'exclusion = douleur sociale, cingulaire antérieur).
        \item \textbf{Cortex préfrontal (immature à 15 ans) :} Submergé par l'afflux noradrénergique (locus coeruleus) et le cortisol chronique (stress post-traumatique latent). Perte de l'inhibition top-down.
        \item \textbf{Axis HPA :} Cortisol basalement élevé (stress chronique) $\rightarrow$ seuil de réactivité abaissé.
        \item \textbf{Via basse exclusive :} Thalamus $\rightarrow$ Amygdale $\rightarrow$ Circuit moteur (frapper) sans passage par l'évaluation préfrontale.
      \end{itemize}
  \item \textbf{Émotion primaire :} \textbf{Colère} (réaction à l'injustice, à l'humiliation) + \textbf{Honte} (exposition du corps stigmatisé) + \textbf{Peur} (impuissance face au groupe, anticipation de la récidive). Le mélange colère/honte est explosif.
  \item \textbf{Réponse maîtrisée possible (à construire a posteriori / préventivement) :}
      \begin{itemize}
        \item \textbf{Immédiat (si possible) :} Respiration 4-7-8 (frein vague) ; phrase CNV apprise : « Arrête. Ça me blesse. J'ai besoin de respect. » ; S'éloigner physiquement vers un adulte référent (surveillant, CPE).
        \item \textbf{Institutionnel (Proviseur/Vie scolaire) :} Ne pas sanctionner \emph{seulement} Kevin (victime qui a craqué). Appliquer la \textbf{méthode de la préoccupation partagée (Pikas)} ou \textbf{justice restaurative} : rencontre médiatisée Kevin / harceleurs pour nommer le tort, réparer (excuses publiques, engagement écrit), suivre les harceleurs (souvent eux-mêmes en souffrance). Sanction éducative pour Kevin (réparation, atelier gestion colère) pas exclusion (qui renforce l'isolement).
        \item \textbf{Structurel :} Programme école « Bienveillance / Lutte anti-harcèlement » (ex. programme \emph{pHARe} adapté au Cameroun) : formation des élèves sentinelles, climat de classe, travail sur l'image corporelle, stigmatisation du poids.
      \end{itemize}
\end{enumerate}
\textbf{Conclusion :} L'analyse physiologique (cortex préfrontal immature + stress chronique) \textbf{exonère de la responsabilité morale totale} (Kevin n'est pas un « voyou ») mais \textbf{n'exonère pas de la responsabilité légale/réparatrice}. Elle oriente vers une réponse \textbf{éducative et thérapeutique} (renforcer le frein préfrontal, soigner la honte, sanctionner le harcèlement) plutôt que purement répressive.
\end{exoresolu}

\begin{exercice}[Entraînement : Analysez votre propre réaction]
Rappelez-vous une situation récente (cette semaine) où vous avez réagi « trop fort » (colère, fuite, paralysie) face à un pair, un parent, un professeur.
\begin{enumerate}
  \item Décrivez le stimulus factuel.
  \item Nommez l'émotion(s) ressentie(s) sur le coup.
  \item Identifiez 2 facteurs biologiques (fatigue ? cycle ? âge ? tempérament ?) et 2 facteurs psychosociaux (stress examen ? conflit familial ? norme de groupe ?) qui ont baissé votre seuil de tolérance.
  \item Écrivez la réponse que vous \emph{auriez voulu} donner (style CNV) et l'action concrète que vous allez mettre en place cette semaine pour renforcer votre cortex préfrontal (ex: \SI{5}{min} respiration quotidienne, sport, journaling émotions).
\end{enumerate}
\end{exercice}

\begin{corrige}[Exercice n°1 -- Exemple de réponse type]
\textit{(Réponse personnelle, pas de solution unique. Voici un modèle de rédaction attendue.)}
\begin{enumerate}
  \item \textbf{Stimulus :} Ma mère crie depuis la cuisine : « Tu n'as encore pas fait la vaisselle ! Tu es paresseux ! » (alors que je venais de finir mes devoirs de maths à \SI{22}{h}).
  \item \textbf{Émotions :} Colère vive (injustice) + Tristesse (non-reconnaissance effort) + Honte (étiquette « paresseux »).
  \item \textbf{Facteurs biologiques :} Fatigue aiguë (\SI{22}{h}, \SI{5}{h} sommeil nuit précédente) $\rightarrow$ cortisol élevé, glucose cérébral bas $\rightarrow$ cortex préfrontal « offline » ; Tempérament sensible (génétique) $\rightarrow$ réactivité amygdalienne haute. \textbf{Facteurs psychosociaux :} Stress Probatoire (maths difficile) ; Dynamique familiale : critique globale (« tu es... ») au lieu de comportement (« tu n'as pas... ») ; Besoin non satisfait de reconnaissance/autonomie.
  \item \textbf{Réponse maîtrisée (CNV) :} « Maman, quand j'entends "tu es paresseux" alors que je viens de bosser \SI{3}{h} mes maths, je me sens blessé et découragé, car j'ai besoin que mon effort soit vu. Est-ce qu'on peut convenir que je fais la vaisselle demain matin avant l'école, et qu'on se parle sans étiquettes ? » \textbf{Action semaine :} \SI{10}{min} cohérence cardiaque chaque soir \SI{22}{h}30 + 2 séances footing \SI{30}{min}.
\end{enumerate}
\end{corrige}

\newpage

\coursec{Activité d'intégration}

\coursec{Maîtrise des comportements interindividuels}

\courssub{Objectifs de l'activité}
Cette activité d'intégration vise à mobiliser l'ensemble des savoirs et savoir-faire de la leçon \og Maîtrise des comportements interindividuels \fg à travers l'analyse d'une situation de conflit réel. Elle permet à l'élève de :
\begin{itemize}
    \item Identifier et distinguer les trois composantes fondamentales de la relation interindividuelle : la \cle{perception}, l'\cle{émotion} et la \cle{communication} (verbale et non verbale).
    \item Reconstituer le \cle{trajet nerveux} (voie rapide sous-corticale et voie lente corticale) et le \cle{trajet hormonal} (axe médullo-surrénalien, adrénaline) sous-tendant une réaction émotionnelle intense (colère/peur).
    \item Représenter graphiquement l'évolution d'un paramètre physiologique (fréquence cardiaque) corrélée aux phases de l'émotion et de sa régulation.
    \item Proposer et justifier physiologiquement des \cle{stratégies de maîtrise de soi} (coping) faisant appel au cortex préfrontal et au système parasympathique.
    \item Élaborer une \cle{règle de vie} favorisant le \cle{vivre-ensemble} et le respect de la diversité des comportements.
\end{itemize}

\courssub{Situation : La bagarre au Marché Mokolo}
\emph{Contexte :} Nous sommes un samedi matin, vers 10 h 30, au Marché Mokolo à Yaoundé. L'affluence est dense, l'air chargé d'odeurs d'épices, de poisson fumé et d'essence. Le bruit est continu : cris des vendeurs, klaxons des motos-taxis, discussions animées.

\emph{Les acteurs :}
\begin{itemize}
    \item \textbf{Mme Ngo}, 45 ans, commerçante de pagnes depuis 20 ans, connue pour sa vigilance.
    \item \textbf{M. Essomba}, 32 ans, fonctionnaire, pressé d'acheter un tissu pour le mariage de sa sœur.
    \item \textbf{M. Atangana}, 50 ans, ancien enseignant à la retraite, passant par là.
\end{itemize}

\emph{Déroulement de la scène :}
\begin{enumerate}
    \item \textbf{Étape 1 (Perception \& Déclencheur) :} M. Essomba choisit un pagne à 12\,000 FCFA. Il tend un billet de 10\,000 FCFA et un billet de 2\,000 FCFA à Mme Ngo. Celle-ci prend le billet de 10\,000 FCFA, le frotte entre ses doigts, le regarde en transparence face au soleil, puis secoue la tête énergiquement. D'une voix forte qui domine le brouhaha, elle lance : \og Eh jeune homme, tu veux m'arnaquer avec ce faux billet ? Tu te prends pour qui ? Tout le monde, regardez ce voleur ! \fg
    \item \textbf{Étape 2 (Réaction émotionnelle immédiate) :} M. Essomba fige. Son visage pâlit puis rougit brutalement. Il sent son \textbf{cœur battre violemment contre sa cage thoracique} (estimation : passage de ~75 à ~130 battements/min en quelques secondes), ses mains deviennent moites et tremblantes, sa respiration devient courte et saccadée. Une chaleur intense lui monte au visage. Il a l'impression que \og le sang lui bout \fg. Sans réfléchir, il crie : \og C'est faux ! Tu m'insultes devant tout le monde ! \fg et lève la main droite, poing serré, en direction de l'épaule de Mme Ngo.
    \item \textbf{Étape 3 (Intervention \& Régulation) :} M. Atangana, qui observait la scène à trois mètres, s'avance rapidement mais calmement. Il pose fermement sa main sur le bras levé de M. Essomba et dit d'une voix grave et posée : \og Mon fils, abaisse la main. Respire. Inspire par le nez... expire par la bouche. Personne ne t'attaque. Mme Ngo, rangez votre billet, parlez calmement. \fg M. Essomba cligne des yeux, comme sorti d'un rêve. Il laisse retomber son bras. Ses épaules s'affaissent. Il prend une grande inspiration tremblante, puis une seconde, plus ample. Son visage dégonfle, la rougeur s'estompe.
    \item \textbf{Étape 4 (Résolution) :} M. Essomba, voix redevenue contrôlée : \og Maman, je vous prie de vérifier avec votre lampe UV ou le stylo détecteur. Ce billet vient du distributeur de la banque tout à l'heure. \fg Mme Ngo, un peu déstabilisée par le calme revenu, sort son stylo détecteur : le trait reste clair. \og Ah... c'est bon, mon fils. Pardon, l'habitude... \fg M. Essomba récupère son billet, paie, prend son pagne et part sans un mot de plus, le pas pressé mais régulier.
\end{enumerate}

\courssub{Consignes}
Travailler par groupes de 4 à 5 élèves. Chaque groupe rend une fiche unique structurée comme suit :

\begin{enumerate}
    \item \textbf{Analyse des composantes :} Pour chacune des 4 étapes du récit, identifier et citer explicitement les éléments relevant de la \textbf{perception} (sensorielle, traitement), de l'\textbf{émotion} (composante subjective, physiologique, comportementale) et de la \textbf{communication} (verbale, paraverbale, non verbale). Compléter le tableau récapitulatif fourni en annexe.
    \item \textbf{Schéma physiologique de la réaction de M. Essomba (Étape 2) :} Sur une feuille A4 paysage, tracer le schéma fonctionnel du trajet de l'information menant à la réaction de colère/peur. Faire apparaître obligatoirement : \textbf{Organes des sens} $\rightarrow$ \textbf{Thalamus} $\rightarrow$ \textbf{Amygdale} (voie rapide) $\rightarrow$ \textbf{Hypothalamus} $\rightarrow$ \textbf{Système nerveux sympathique} / \textbf{Médullo-surrénale (Adrénaline)} $\rightarrow$ \textbf{Effeteurs (Cœur, vaisseaux, muscles)}. Indiquer \textbf{en parallèle} la voie lente : \textbf{Thalamus} $\rightarrow$ \textbf{Cortex sensoriel} $\rightarrow$ \textbf{Cortex préfrontal} (évaluation, inhibition). Légender les flèches \og Voie rapide (inconsciente, automatique) \fg et \og Voie lente (consciente, contrôlée) \fg.
    \item \textbf{Graphique de l'évolution de la Fréquence Cardiaque (FC) :} Sur un repère orthonormé (axe des abscisses : Temps en minutes, de 0 à 5 min ; axe des ordonnées : FC en bpm, de 50 à 150 bpm), tracer la courbe qualitative de la FC de M. Essomba couvrant les 4 étapes. Annoter 4 points caractéristiques : $T_0$ (base), $T_1$ (pic émotionnel), $T_2$ (début régulation), $T_3$ (retour au calme). Justifier la forme de la courbe par les systèmes nerveux impliqués.
    \item \textbf{Réécriture et justification physiologique :} Réécrire la fin de la scène (Étape 4) en imaginant que M. Essomba \textbf{n'a pas été interrompu} par M. Atangana, mais qu'il a \textbf{lui-même} mis en œuvre une stratégie de maîtrise de soi (ex: respiration diaphragmatique, comptage, retrait physique, reformulation interne). Rédiger le dialogue intérieur et l'action. Justifier point par point l'efficacité physiologique de cette stratégie (structures cérébrales, systèmes nerveux, hormones).
    \item \textbf{Règle de vie :} Formuler collectivement une \og Règle d'or du marché \fg (ou de tout espace public) sous forme d'une phrase courte, positive, applicable, qui synthétise la leçon tirée de cette analyse sur le respect d'autrui et la gestion des émotions.
\end{enumerate}

\courssub{Ressources à mobiliser}
\begin{itemize}
    \item \textbf{Perception :} Transduction sensorielle $\rightarrow$ Voies afférentes $\rightarrow$ Thalamus (relais) $\rightarrow$ Aires corticales primaires $\rightarrow$ Intégration multisensorielle.
    \item \textbf{Émotion (Modèle composantaire) :} Composante \emph{subjective} (ressenti), \emph{neurophysiologique} (SNA, hormones), \emph{comportementale} (expression faciale, fuite/attaque), \emph{cognitive} (évaluation).
    \item \textbf{Voies de la peur/colère (LeDoux) :}
        \begin{itemize}
            \item \cle{Voie rapide (basse)} : Thalamus $\rightarrow$ Amygdale $\rightarrow$ Hypothalamus/Tronc cérébral $\rightarrow$ Réaction immédiate (Sympathique + Adrénaline). \textbf{Temps :} $\approx$ 12-20 ms. Pas de conscience.
            \item \cle{Voie lente (haute)} : Thalamus $\rightarrow$ Cortex sensoriel $\rightarrow$ Hippocampe (contexte) $\rightarrow$ Cortex préfrontal (évaluation, décision, inhibition) $\rightarrow$ Amygdale (modulation). \textbf{Temps :} $\approx$ 300-500 ms. Conscience et contrôle.
        \end{itemize}
    \item \textbf{Système Nerveux Autonome (SNA) :}
        \begin{center}
            \begin{tabular}{|l|c|c|}\hline
            \textbf{Organe/Effeteur} & \textbf{Sympathique (Orthosympathique)} & \textbf{Parasympathique (Vague)} \\ \hline
            Cœur (Fréquence/Force) & $\uparrow\uparrow$ (Noradrénaline/Adrénaline sur $\beta_1$) & $\downarrow$ (Acétylcholine sur Muscariniques $M_2$) \\ \hline
            Vaisseaux (Muscle lisse) & Vasoconstriction ($\alpha_1$) / Vasodilatation muscle ($\beta_2$) & Quasiment nul \\ \hline
            Bronches & Bronchodilatation ($\beta_2$) & Bronchoconstriction ($M_3$) \\ \hline
            Transpiration & $\uparrow\uparrow$ (Cholinergique sympathique) & -- \\ \hline
            \end{tabular}
        \end{center}
    \item \textbf{Communication :} \cle{Verbale} (mots), \cle{Paraverbale} (ton, débit, volume, timbre), \cle{Non verbale} (regard, posture, gestuelle, distance interpersonnelle/proxémie). \cle{Congruence} vs \cle{Incongruence}.
    \item \textbf{Régulation émotionnelle (Stratégies) :} \cle{Respiration lente/abdominale} $\rightarrow$ Stimulation barorécepteurs / Nerf vague $\rightarrow$ Activation Parasympathique $\rightarrow$ Inhibition Amygdale via Cortex Préfrontal. \cle{Recadrage cognitif (Réévaluation)} $\rightarrow$ Cortex Préfrontal Dorsolatéral/Ventromédian $\rightarrow$ Modulation Amygdale.
\end{itemize}

\courssub{Démarche et corrigé commenté}
\begin{exoresolu}[Corrigé détaillé de l'activité d'intégration]
\textbf{Consigne 1 : Analyse des composantes (Tableau récapitulatif)}

\begin{center}
\begin{tabularx}{\linewidth}{|l|X|X|X|}\hline
\textbf{Étape} & \textbf{Perception} & \textbf{Émotion} & \textbf{Communication} \\ \hline
\textbf{1. Accusation} &
\begin{itemize}
\item \emph{Visuelle :} Mme Ngo voit le billet, transparence, texture.
\item \emph{Tactile :} Frottement papier.
\item \emph{Auditive :} Brouhaha marché (bruit de fond).
\item \emph{Traitement :} Reconnaissance motif \og faux billet \fg (mémoire sémantique) $\rightarrow$ Décision \og danger/arnaque \fg.
\end{itemize} &
\begin{itemize}
\item \emph{Sujet (Mme Ngo) :} Méfiance, colère froide, sentiment de justice/légitimité.
\item \emph{Physio (Mme Ngo) :} Activation modérée SNA (vigilance).
\item \emph{Comportement :} Accusation publique, désignation du doigt.
\end{itemize}
&
\begin{itemize}
\item \emph{Verbale :} \og Faux billet \fg, \og Voleur \fg (mots blessants, accusatoires).
\item \emph{Paraverbale :} Voix forte, ton aigu, débit rapide (perçage bruit).
\item \emph{Non verbale :} Doigt tendu, regard dur, posture droite/dominante, proxémie proche.
\end{itemize} \\ \hline
\textbf{2. Réaction M. Essomba} &
\begin{itemize}
\item \emph{Auditive :} Mots \og voleur \fg, \og arnaquer \fg perçus comme menace sociale/physique.
\item \emph{Visuelle :} Doigt pointé, foule qui se tourne (regards juges).
\item \emph{Interception :} Tachycardie, sueurs, chaleur (signaux corporels perçus).
\item \emph{Traitement :} \textbf{Via voie rapide} : Amygdale $\rightarrow$ \og Menace vitale \fg. Pas de temps pour analyse cortex.
\end{itemize} &
\begin{itemize}
\item \emph{Subjective :} Sidération $\rightarrow$ Honte intense $\rightarrow$ Colère explosive (défense honneur).
\item \emph{Physio :} \textbf{Explosion sympathique + Adrénaline}. FC $\uparrow\uparrow$, TA $\uparrow$, Tremblements (adrénaline $\beta_2$), Pâleur $\rightarrow$ Rougeur (vasoconstriction puis dilatation), Apnée puis hyperventilation.
\item \emph{Comportementale :} Cri, levée du poing (attaque), figement initial.
\end{itemize}
&
\begin{itemize}
\item \emph{Verbale :} \og C'est faux ! \fg (Négation), \og Tu m'insultes \fg (Plainte).
\item \emph{Paraverbale :} Cri strident, voix cassée par l'émotion, débit haché.
\item \emph{Non verbale :} Poing levé (menace), corps en avant (anticipation attaque), regard fixe/évitant (honte/colère), distance réduite.
\end{itemize} \\ \hline
\textbf{3. Intervention M. Atangana} &
\begin{itemize}
\item \emph{Visuelle :} Bras levé, visage rouge, femme effrayée.
\item \emph{Auditive :} Cris, bruit foule.
\item \emph{Tactile :} Main sur bras (contact physique apaisant, stimulation cutanée C-tactiles).
\item \emph{Traitement :} \textbf{Voie lente activée}. Perception \og Aide/Autorité bienveillante \fg $\rightarrow$ Cortex Préfrontal $\rightarrow$ Inhibition Amygdale.
\end{itemize} &
\begin{itemize}
\item \emph{Sujet (M. Atangana) :} Empathie, calme, autorité sereine.
\item \emph{Sujet (M. Essomba) :} \textbf{Bascule} : Colère $\rightarrow$ Surprise $\rightarrow$ Soulagement/Honte résiduelle. Physio : Début activation Parasympathique (Freinage).
\item \emph{Comportement :} Bras retombe, épaules basses, respiration reprise.
\end{itemize}
&
\begin{itemize}
\item \emph{Verbale :} \og Abaisse la main \fg (Ordre court), \og Respire \fg (Instruction corporelle), \og Personne ne t'attaque \fg (Recadrage cognitif/Réassurance).
\item \emph{Paraverbale :} Voix grave, lente, stable, volume modéré (contraste apaisant).
\item \emph{Non verbale :} Contact main-bras (ancrage), posture ouverte, regard dans les yeux (connexion), distance respectueuse.
\end{itemize} \\ \hline
\textbf{4. Résolution} &
\begin{itemize}
\item \emph{Visuelle :} Stylo détecteur, trait clair.
\item \emph{Auditive :} Silence relatif, mots d'excuses.
\item \emph{Traitement :} Confirmation \og Billet vrai \fg $\rightarrow$ Réévaluation : \og Danger écarté \fg $\rightarrow$ Inhibition totale voie rapide.
\end{itemize} &
\begin{itemize}
\item \emph{M. Essomba :} Soulagement, fatigue post-adrénalinique, résidu de honte/injustice (géré par cortex).
\item \emph{Mme Ngo :} Gêne, empathie tardive.
\item \emph{Physio :} Retour base (Homeostasie). Parasympathique dominant.
\end{itemize}
&
\begin{itemize}
\item \emph{Verbale :} \og Vérifiez \fg (Demande rationnelle), \og Pardon \fg (Réparation sociale), \og C'est bon \fg (Validation).
\item \emph{Paraverbale :} Calme revenu, débit normal.
\item \emph{Non verbale :} Prise pagne, départ rapide (évitement lieu stress), pas de regard soutenu (protection visage).
\end{itemize} \\ \hline
\end{tabularx}
\end{center}

\textbf{Consigne 2 : Schéma physiologique de la réaction (Étape 2)}

\begin{popfigure}
\begin{tikzpicture}[
    node distance=1.2cm and 1.5cm,
    box/.style={rectangle, draw=popdark, thick, fill=popblueL, rounded corners, minimum height=1.1cm, text width=2.8cm, align=center, font=\small},
    arrow/.style={-Stealth, thick, popdark},
    fast/.style={arrow, poporange, line width=1.5pt},
    slow/.style={arrow, popgreen, line width=1.5pt, dashed},
    hormone/.style={arrow, poppurple, line width=1.2pt, densely dotted}
]
% Nœuds
\node[box, align=center] (sense){Organes des sens\\ (Yeux, Oreilles, Peau)};
\node[box, right=of sense, align=center] (thal){Thalamus\\(Relais sensoriel)};
\node[box, below right=of thal, align=center] (amyg){Amigdale\\(Centre peur/colère,\\ détection menace)};
\node[box, above right=of thal, align=center] (ctxsens){Cortex sensoriel\\(Analyse fine stimulus)};
\node[box, right=of ctxsens, align=center] (hippo){Hippocampe\\(Contexte, Mémoire)};
\node[box, right=of hippo, align=center] (pfc){Cortex Préfrontal\\(Évaluation, Décision,\\ Inhibition contrôle)};
\node[box, below=of amyg, align=center] (hypo){Hypothalamus\\(Centre commande SNA)};
\node[box, right=of hypo, align=center] (symp){Système Nerveux\\ Sympathique\\ (Ganglions paravertébraux)};
\node[box, below=of symp, align=center] (medulla){Médullo-surrénale\\ (Libération Adrénaline\\ dans le sang)};
\node[box, right=of symp, align=center] (effect){Effecteurs :\textbf{Cœur} (FC $\uparrow$),\\ \textbf{Vaisseaux}, \textbf{Muscles},\\ \textbf{Glandes sudoripares}};

% Flèches Voie Rapide (Orange)
\draw[fast] (sense) -- node[above, font=\tiny] {Influx nerveux} (thal);
\draw[fast] (thal) -- node[above left, font=\tiny] {\textbf{Voie Rapide (Basse)}} (amyg);
\draw[fast] (amyg) -- node[left, font=\tiny] {Projections directes} (hypo);
\draw[fast] (hypo) -- (symp);
\draw[fast] (symp) -- (effect);
\draw[hormone] (symp) -- node[left, font=\tiny] {Innervation cholinergique} (medulla);
\draw[hormone] (medulla) -- node[below, font=\tiny, align=center] {Adrénaline $\rightarrow$ Sang\\ $\rightarrow$ Cible lointaine} (effect);

% Flèches Voie Lente (Vert pointillé)
\draw[slow] (thal) -- node[above, font=\tiny] {\textbf{Voie Lente (Haute)}} (ctxsens);
\draw[slow] (ctxsens) -- (hippo);
\draw[slow] (hippo) -- (pfc);
\draw[slow, bend left=15] (pfc) to node[above, font=\tiny, align=center]{Modulation / Inhibition\\ (GABA, Glutamate)} (amyg);
\draw[slow, bend right=15] (pfc) to node[below, font=\tiny, align=center]{Contrôle moteur\\ (Inhibition geste)} (symp);

% Légendes boîtes
\node[draw=poporange, thick, rounded corners, fill=poporange!10, anchor=north west, font=\small, align=left] at (amyg.north east) {\textbf{Légende :}\\\textcolor{poporange}{---} Voie rapide (Sous-corticale, < 50 ms, automatique, inconsciente)\\\textcolor{popgreen}{- - -} Voie lente (Corticale, > 300 ms, consciente, contrôlée)\\\textcolor{poppurple}{\ldots\ldots} Voie hormonale (Adrénaline, secondes-minutes, diffuse)};
\end{tikzpicture}
\legende{Schéma fonctionnel des deux voies de traitement de l'information émotionnelle (modèle de LeDoux) chez M. Essomba à l'étape 2. La voie rapide (orange) explique la réaction explosive immédiate (levée du poing) avant toute analyse consciente. La voie lente (verte) arrive trop tard pour empêcher l'acte, mais est la cible de la régulation (M. Atangana / Auto-contrôle). La voie hormonale (violette) prolonge et amplifie l'état physiologique.}
\end{popfigure}

\textbf{Consigne 3 : Graphique de l'évolution de la Fréquence Cardiaque (FC)}

\begin{popfigure}
\begin{tikzpicture}
\begin{axis}[
    width=\linewidth, height=0.55\linewidth,
    xlabel={Temps (min)}, ylabel={Fréquence Cardiaque (bpm)},
    xmin=0, xmax=5, ymin=50, ymax=150,
    xtick={0,1,2,3,4,5}, xticklabels={$T_0$ (Base), $T_1$ (Pic), $T_2$ (Régul.), $T_3$ (Calme), $T_4$, $T_5$},
    ytick={60,75,90,110,130,150},
    grid=major, grid style={popline},
    axis lines=middle, axis line style={thick, popdark},
    tick label style={font=\small}, label style={font=\small},
    clip=false
]
% Courbe qualitative tracée à la main via coordinates
\addplot[popcurve, very thick, smooth, tension=0.7] coordinates {
    (0, 72)   % T0 Base
    (0.5, 75)
    (1, 132)  % T1 Pic Émotionnel (Étape 2)
    (1.5, 125)
    (2, 110)  % T2 Début Régulation (Intervention M. Atangana)
    (2.5, 95)
    (3, 85)
    (3.5, 80)
    (4, 76)   % T3 Retour calme (Étape 4)
    (4.5, 74)
    (5, 72)
};
% Annotations
\node[anchor=south, font=\small, text=popdark] at (axis cs:0,72) {$T_0 \approx 72$ bpm};
\node[anchor=south, font=\small, text=poporange, align=center] at (axis cs:1,132){$T_1 \approx 132$ bpm \\ \textbf{Pic Sympathique + Adrénaline}};
\node[anchor=north east, font=\small, text=popgreen, align=center] at (axis cs:2,110){$T_2 \approx 110$ bpm \\ \textbf{Activation Parasympathique (Vague)}};
\node[anchor=north, font=\small, text=popblue, align=center] at (axis cs:4,76){$T_3 \approx 76$ bpm \\ \textbf{Retour Homeostasie}};
% Étiquettes de phases (remplacement des zones colorées fillbetween)
\node[rotate=90, font=\small, poporange, anchor=south] at (axis cs:1,140) {\textbf{Phase Aiguë (Sympathique)}};
\node[rotate=90, font=\small, popgreen, anchor=south] at (axis cs:2.8,140) {\textbf{Phase Récupération (Parasympathique)}};
\end{axis}
\end{tikzpicture}
\legende{Évolution qualitative de la fréquence cardiaque de M. Essomba. $T_0$ : État basal (équilibre SNA). $T_1$ : Pic brutal lors de l'accusation et de la levée du poing (activation massive orthosympathique + adrénaline circulante). $T_2$ : Infléchissement net dès l'intervention de M. Atangana (stimulation du nerf vague par le contact, la voix, la respiration guidée). $T_3$ : Retour progressif à la normale après vérification du billet (levée de l'incertitude cognitive, inhibition corticale de l'amygdale).}
\end{popfigure}

\textbf{Justification physiologique de la courbe :}
\begin{itemize}
    \item \textbf{Segment $T_0 \rightarrow T_1$ (Pente raide $\nearrow$) :} Activation de la \textbf{voie rapide} (Amigdale $\rightarrow$ Hypothalamus $\rightarrow$ Centre cardiaque bulbaire $\rightarrow$ Nerfs sympathiques cardiaques). Libération de \textbf{noradrénaline} sur récepteurs $\beta_1$ (augmentation FC et force de contraction) + arrivée de l'\textbf{adrénaline} sanguine (effet plus lent, plus durable). Inhibition simultanée du tonus vagal (parasympathique).
    \item \textbf{Segment $T_1 \rightarrow T_2$ (Palier/Légère descente) :} Persistance de l'adrénaline circulante (demi-vie $\approx$ 2-3 min) maintient FC haute malgré début d'inhibition centrale.
    \item \textbf{Segment $T_2 \rightarrow T_3$ (Pente descendante $\searrow$ prononcée) :} \textbf{Réactivation massive du parasympathique (Nerf Vague)}. Mécanismes : 1) \emph{Baroréflexe} : Hypertension artérielle $\rightarrow$ Barorécepteurs carotidiens/aortiques $\rightarrow$ Nucleus Tractus Solitarius (NTS) $\rightarrow$ Noyau Ambigu $\rightarrow$ Vague $\rightarrow$ Cœur (Acétylcholine sur $M_2$). 2) \emph{Commande volontaire (Cortex Préfrontal $\rightarrow$ NTS/Noyau Ambigu)} via respiration lente et contact tactile (fibres C-tactiles $\rightarrow$ Insula $\rightarrow$ NTS). 3) \emph{Inhibition amygdale} par Cortex Préfrontal $\rightarrow$ Arrêt drive sympathique.
    \item \textbf{Segment $T_3 \rightarrow T_5$ (Plateau) :} Retour à l'équilibre homéostatique (Tonus vagal dominant au repos).
\end{itemize}

\textbf{Consigne 4 : Réécriture de la fin (Auto-régulation) et Justification}

\emph{Scénario alternatif (M. Essomba seul) :}
\og \textit{M. Essomba sent son poing trembler à deux centimètres de l'épaule de Mme Ngo. La foule crie. Soudain, une image lui traverse l'esprit : le visage de sa petite fille qui l'attend à la maison. \textbf{Il inspire brutalement par le nez en gonflant le ventre (5 secondes)}, bloque sa respiration (2 secondes), \textbf{expire longuement par la bouche entrouverte comme s'il soufflait dans une paille (8 secondes)}. Il répète : \og Un... Deux... Trois... \fg (comptage interne). Il sent ses épaules descendre. Il \textbf{recule d'un pas} (augmentation distance proxémique). Il se dit intérieurement : \og Elle a peur pour son commerce, elle ne me connaît pas. Ce n'est pas une attaque contre moi. \fg (Réévaluation cognitive). Il abaisse le poing, ouvre la main, paume visible. D'une voix qu'il force à être basse : \og Maman, calmez-vous. Je ne suis pas un voleur. Vérifiez avec votre machine, s'il vous plaît. \fg} \fg

\emph{Justification physiologique point par point :}

\begin{center}
\begin{tabularx}{\linewidth}{|l|X|X|}\hline
\textbf{Action de maîtrise} & \textbf{Mécanisme Neurophysiologique Mobilisé} & \textbf{Effet sur la Réaction de Stress} \\ \hline
\textbf{1. Respiration lente, abdominale, expiration longue} & Stimulation mécanique des \textbf{barorécepteurs} (pression thoracique/abdominale) + Stimulation des \textbf{récepteurs d'étirement poumons} $\rightarrow$ Afférences Vague (X) $\rightarrow$ \textbf{NTS / Noyau Ambigu} $\rightarrow$ \textbf{Efférences Vague (Parasympathique)} $\rightarrow$ Cœur (Acétylcholine/$M_2$). & \textbf{Freinage cardiaque immédiat} (RSA - Arythmie Sinusale Respiratoire). Baisse FC et TA $\rightarrow$ Signal de sécurité envoyé au cerveau (Interoception). \\ \hline
\textbf{2. Comptage interne / Focalisation attention} & Activation du \textbf{Cortex Préfrontal Dorsolatéral (CPFdl)} (Contrôle exécutif, Mémoire de travail, Attention focalisée) $\rightarrow$ Inhibition des réseaux de mode par défaut (rumination) et de l'\textbf{Amigdale} (via projections glutamatergiques/GABAergiques vers noyaux intercalaires). & \textbf{Découplage "Stimulus $\rightarrow$ Réponse"}. Arrêt de la boucle de rétroaction émotionnelle (pensées catastrophistes). Passage du "Mode Menace" au "Mode Tâche". \\ \hline
\textbf{3. Recul physique (Pas en arrière)} & Modification de la \textbf{proxémie} (distance de sécurité). Signal visuel/proprioceptif $\rightarrow$ Cortex Pariétal / Prémoteur $\rightarrow$ \textbf{Cortex Préfrontal Ventromédian (CPFvm)} $\rightarrow$ Évaluation \og Danger physique écarté \fg $\rightarrow$ Inhibition Amygdale. & Réduction de l'urgence motrice (Préparation attaque/fuite). Signal corporel de \og non-menace \fg (Embodied Cognition). \\ \hline
\textbf{4. Réévaluation cognitive (Recadrage : "Elle a peur, ce n'est pas personnel")} & \textbf{CPFvm / Cortex Orbito-Frontal} : Attribution d'intention bienveillante/neutre (Théorie de l'Esprit) $\rightarrow$ Changement de la \textbf{valence émotionnelle} du stimulus (Menace $\rightarrow$ Problème technique). Projections inhibitrices directes sur \textbf{Amigdale} (via noyaux intercalaires GABAergiques). & \textbf{Extinction de la réponse de peur/colère} à la source. Arrêt de la commande hypothalamique (Sympathique) et surrénalienne (Adrénaline). \\ \hline
\textbf{5. Communication non verbale apaisée (Main ouverte, regard direct mais doux, voix basse)} & Boucle de rétroaction \textbf{Sensorimotrice} : Production de signaux de sécurité (visage détendu, prosodie basse) $\rightarrow$ Perception propre (interoception/exteroception) $\rightarrow$ Renforcement état calme (Théorie de James-Lange inversée / Facial Feedback). Activation système \textbf{d'engagement social} (Porges - Vague Myélinisé). & \textbf{Co-régulation} : Apaise aussi Mme Ngo (neurones miroirs, contagion émotionnelle positive). Stabilise l'état physiologique de l'émetteur. \\ \hline
\end{tabularx}
\end{center}

\textbf{Consigne 5 : Règle de vie (Exemples de productions attendues)}
\begin{itemize}
    \item \og \textbf{Avant de réagir à une offense, je respire trois fois et je cherche à comprendre l'intention de l'autre.} \fg
    \item \og \textbf{Ma dignité ne se défend pas par la violence, mais par la maîtrise de ma parole et de mes gestes.} \fg
    \item \og \textbf{Au marché comme ailleurs : "Vérifier avant d'accuser, Respirer avant de frapper, Parler pour se comprendre."} \fg
\end{itemize}
\emph{Critères de validation :} Phrase courte, impératif positif, mobilise une compétence (respiration, vérification, parole), applicable immédiatement, favorise le lien social.

\end{exoresolu}

\courssub{Bilan de l'activité}
Cette activité d'intégration a permis de mettre en évidence, de manière concrète et contextualisée, l'articulation fine entre les trois piliers de la relation interindividuelle :
\begin{enumerate}
    \item \textbf{L'unité fonctionnelle Perception-Émotion-Communication :} On ne perçoit jamais "neutrement" ; la perception est teintée d'émotion (évaluation amygdalaire) et s'exprime immédiatement par une communication multimodale (verbale, paraverbale, non verbale). L'incongruence entre le message verbal (\og C'est faux \fg) et le non-verbal (poing levé, visage rouge) chez M. Essomba est la signature de la \cle{désorganisation émotionnelle} (prise de contrôle par la voie rapide).
    \item \textbf{La temporalité double du cerveau émotionnel :} Le schéma et le graphique FC illustrent parfaitement le conflit entre la \cle{voie rapide sous-corticale} (survie, automaticité, 20 ms, source de l'agression) et la \cle{voie lente corticale} (adaptation, socialisation, 300 ms, source de la maîtrise). L'éducation et la culture visent à \textbf{raccourcir le délai} et \textbf{renforcer le poids} de la voie lente sur la voie rapide (plasticité synaptique du circuit Préfrontal-Amygdale).
    \item \textbf{Le corps comme levier de régulation :} L'analyse de la courbe FC et de la stratégie respiratoire démontre que la \cle{maîtrise de soi n'est pas une abstraction mentale} mais un acte physiologique précis : stimulation volontaire du \textbf{nerf vague} (parasympathique) pour freiner le cœur et inhiber l'amygdale. La respiration est le "frein à main" accessible à tous, partout, sans matériel.
    \item \textbf{La dimension sociale de la régulation (Co-régulation) :} L'intervention de M. Atangana montre que la régulation émotionnelle est souvent \textbf{interpersonnelle} avant d'être intrapersonnelle. Le contact tactile (fibres C-tactiles $\rightarrow$ Insula $\rightarrow$ Système de l'engagement social), la prosodie apaisante et le recadrage cognitif externe ("Personne ne t'attaque") sont des "béquilles corticales" prêtées à celui qui a perdu son cortex préfrontal.
    \item \textbf{Respect de la diversité et Vivre-ensemble :} La règle de vie finale ancre la biologie dans l'éthique. Comprendre que l'agressivité de l'autre est souvent une \textbf{réaction physiologique de défense} (peur, protection honneur, stress) et non une essence méchante, fonde l'\cle{empathie} et la \cle{tolérance}. C'est la base scientifique du \cle{vivre-ensemble} prôné par nos programmes.
\end{enumerate}
L'élève sort de cette activité avec une \textbf{grammaire interne} lui permettant de décoder ses propres crises et celles d'autrui, et une \textbf{boîte à outils physiologique} (Respiration, Recadrage, Recul, Communication non violente) validée par les neurosciences affectives.

\newpage

\coursec{Méthodes et exercices résolus}

\courssub{Méthode 1 : Expliquer un comportement interindividuel à partir de son support physiologique (système nerveux et hormonal)}
\begin{methode}[Démarche pour expliquer un comportement par ses bases physiologiques]
    \textbf{Étape 1 : Identifier le comportement observé et le stimulus déclencheur.} Préciser s'il s'agit d'une réaction de fuite, d'agression, de coopération, de séduction, etc., et quel événement l'a provoqué (menace, présence d'un congénère, signal sensoriel).
    \textbf{Étape 2 : Décrire la voie nerveuse sensorielle.} Nommer le récepteur (visuel, auditif, olfactif...), le nerf afférent et les aires corticales d'intégration primaire (cortex visuel, auditif...).
    \textbf{Étape 3 : Identifier les centres d'intégration émotionnelle et décisionnelle.} Citer le \textbf{système limbique} (amygdale pour la peur/agression, hippocampe pour la mémoire contextuelle, hypothalamus pour la réponse neurovégétative/endocrine) et le \textbf{cortex préfrontal} (contrôle inhibiteur, décision sociale).
    \textbf{Étape 4 : Décrire la voie efférente et les effecteurs.}
    \begin{itemize}
        \item \textbf{Via somatique :} Cortex moteur $\to$ moelle épinière $\to$ nerfs moteurs $\to$ muscles squelettiques (mouvements volontaires : fuite, geste, expression faciale).
        \item \textbf{Via neurovégétative (orthosympathique/parasympathique) :} Hypothalamus $\to$ moelle $\to$ ganglions $\to$ organes (cœur, poumons, tube digestif, glandes sudoripares) pour l'adaptation physiologique.
        \item \textbf{Via endocrine :} Hypothalamus $\to$ hypophyse (antérieure/postérieure) $\to$ glandes cibles (surrénales, gonades, thyroïde) $\to$ hormones dans le sang (adrénaline, cortisol, ocytocine, testostérone) pour une action diffuse et durable.
    \end{itemize}
    \textbf{Étape 5 : Rédiger l'explication synthétique.} Relier le stimulus $\to$ traitement central $\to$ réponse effectrice (motrice, végétative, hormonale) $\to$ comportement final. Utiliser les termes précis : \cle{amygdale}, \cle{hypothalamus}, \cle{système orthosympathique}, \cle{adrénaline}, \cle{ocytocine}, \cle{cortex préfrontal}.
\end{methode}

\begin{exoresolu}[Analyse physiologique de la réaction de « trac » avant un exposé oral]
Un élève de Première A, Pierre, doit passer un exposé oral devant sa classe. Quelques minutes avant son passage, il ressent une accélération de son rythme cardiaque, des mains moites, une gorge sèche et une envie pressante d'uriner. Il a du mal à se concentrer sur ses notes.
\tcblower
\textbf{Solution détaillée}

\textbf{1. Identification du comportement et du stimulus}
\begin{itemize}
    \item \textbf{Comportement :} Réaction de stress aigu (anxiété de performance / « trac »).
    \item \textbf{Stimulus :} Situation d'évaluation sociale (exposé oral face à un groupe de pairs et un enseignant), perçue comme une menace pour l'estime de soi et le statut social.
\end{itemize}

\textbf{2. Voie sensorielle et intégration corticale}
La situation est perçue par la \textbf{vision} (public, salle) et l'\textbf{ouïe} (silence, attente). L'information visuelle (nerf optique $\to$ corps genouillé latéral $\to$ cortex visuel primaire occipital) et auditive (nerf auditif $\to$ corps genouillé médial $\to$ cortex auditif primaire temporal) converge vers les \textbf{aires d'association multimodales} (cortex préfrontal dorsolatéral, cortex cingulaire antérieur) qui analysent la signification sociale : « Je suis évalué, je peux échouer ».

\textbf{3. Centres d'intégration émotionnelle (Système limbique)}
L'information parvient rapidement (voie courte, sous-corticale) à l'\textbf{amygdale} (noyau latéral $\to$ noyau central). L'amygdale est la structure clé de la détection de la menace et de la peur conditionnée. Elle active simultanément :
\begin{itemize}
    \item L'\textbf{hypothalamus} (centre de commande neurovégétatif et endocrinien).
    \item Le \textbf{locus cœruleus} (noradrénaline) pour la vigilance.
    \item Le \textbf{noyau du raphé} (sérotonine) modulant l'humeur.
\end{itemize}
Le \textbf{cortex préfrontal ventromédian} (siège du contrôle inhibiteur et de la régulation émotionnelle) tente de moduler l'activité de l'amygdale (réévaluation cognitive : « Ce n'est qu'un exercice, je suis préparé »), mais chez Pierre, l'activation amygdalienne domine (désinhibition préfrontale par le stress).

\textbf{4. Voies efférentes et effecteurs (Réponse physiologique)}

\begin{center}
\begin{tabularx}{\linewidth}{|l|X|X|}\hline
\textbf{Voie} & \textbf{Mécanisme} & \textbf{Manifestations observées chez Pierre} \\ \hline
\textbf{Neurovégétative Orthosympathique} & Hypothalamus $\to$ colonne latérale intermédiolatérale moelle (T1-L2) $\to$ ganglions paravertébraux $\to$ organes cibles. Libération de \textbf{noradrénaline} sur les effecteurs. & \textbf{Cœur :} Tachycardie (fréquence $\uparrow$, force $\uparrow$) $\to$ « cœur qui bat la chamade ». \\
& & \textbf{Vaisseaux cutanés :} Vasoconstriction $\to$ pâleur, mains froides. \\
& & \textbf{Glandes sudoripares :} Stimulation cholinergique $\to$ \textbf{sudation profuse} (mains moites). \\
& & \textbf{Tube digestif :} Inhibition motrice et sécrétoire $\to$ \textbf{bouche sèche} (xérostomie), nausées. \\
& & \textbf{Vessie :} Relâchement du détrusor + contraction sphincter interne $\to$ \textbf{envie d'uriner} (réflexe de vidange par peur). \\ \hline
\textbf{Endocrine (Axe HHS)} & Hypothalamus (CRH) $\to$ Hypophyse anté (ACTH) $\to$ Corticosurrénales (Cortisol). & Action lente (minutes/heures) : Mobilisation énergie (glycémie $\uparrow$), inhibition immunitaire, effet sur mémoire (hippocampe). \\ \hline
\textbf{Endocrine (Médullosurrénale)} & Orthosympathique $\to$ Médullosurrénale $\to$ \textbf{Adrénaline} (80\%) + Noradrénaline (20\%) dans le sang. & Amplifie et prolonge les effets orthosympathiques sur tout l'organisme (cœur, foie, muscles). \\ \hline
\textbf{Somatique (Volontaire)} & Cortex moteur $\to$ Voie pyramidale $\to$ Muscles. & Tremblements (co-contraction musculaire), voix tremblante, difficultés motrices fines (tenir feuille). \\ \hline
\end{tabularx}
\end{center}

\textbf{5. Conclusion}
Le « trac » de Pierre est une \textbf{réaction de stress aigu} déclenchée par l'évaluation sociale. L'\textbf{amygdale} détecte la menace, l'\textbf{hypothalamus} orchestre la réponse via le \textbf{système orthosympathique} (effets immédiats : tachycardie, sudation, bouche sèche, envie mictionnelle) et l'\textbf{axe corticosurrénalien} (cortisol) + \textbf{médullosurrénalien} (adrénaline) pour la maintien de l'alerte. La difficulté à se concentrer résulte de l'inhibition du \textbf{cortex préfrontal} par l'amygdale (phénomène de \emph{hijacking} amygdalien).
\end{exoresolu}

\courssub{Méthode 2 : Analyser la perception d'autrui (reconnaissance faciale et théorie de l'esprit)}
\begin{methode}[Démarche pour analyser la perception d'autrui]
    \textbf{Étape 1 : Distinguer les deux niveaux de traitement.}
    \begin{itemize}
        \item \textbf{Traitement bas-niveau (automatique, pré-attentif) :} Détection de la présence d'un visage, analyse des traits physiques (genre, âge, émotion faciale brute). Structures : \textbf{gyrus fusiforme} (FFA - Fusiform Face Area), \textbf{gyrus occipital inférieur} (OFA), \textbf{sillon temporal supérieur} (STS - regard, mouvement lèvres).
        \item \textbf{Traitement haut-niveau (contrôlé, contextuel) :} Reconnaissance de l'identité (mémoire sémantique), attribution d'intentions, d'états mentaux (croyances, désirs) : \textbf{Théorie de l'Esprit (TdE)}. Structures : \textbf{jonction temporo-pariétale} (TPJ), \textbf{cortex préfrontal médian} (mPFC), \textbf{pôle temporal antérieur}.
    \end{itemize}
    \textbf{Étape 2 : Identifier le rôle des neurones miroirs.} Situés dans le \textbf{cortex préfrontal ventral} (aire F5 chez le singe, homologue aire 44/45 de Broca chez l'homme) et \textbf{lobe pariétal inférieur}. Ils s'activent \textbf{à l'exécution} ET \textbf{à l'observation} d'une action/émotion $\to$ base neuronale de l'\cle{empathie} (résonance motrice/émotionnelle) et de l'imitation.
    \textbf{Étape 3 : Intégrer le contexte culturel et expérientiel.} La perception n'est pas une copie du réel : elle est une \textbf{inférence bayésienne} (prédiction based on \cle{a priori} culturels, stéréotypes, expériences passées).
    \textbf{Étape 4 : Formaliser l'analyse.} Stimulus (visage/voix/geste) $\to$ Voie ventrale (Quoi/Qui : FFA, Temporal antérieur) + Voie dorsale (Où/Comment : STS, Pariétal) $\to$ Réseau TdE (TPJ, mPFC) $\to$ Attribution d'état mental $\to$ Comportement social adapté.
\end{methode}

\begin{exoresolu}[Cas de méprise identitaire et attribution d'intention]
Lors d'une fête de fin d'année au lycée, Amina aperçoit son ami Koffi de dos. Elle lui tape joyeusement sur l'épaule en criant « Koffi, tu danses bien ! ». L'homme se retourne : c'est un inconnu qui ressemble à Koffi (même taille, même coupe de cheveux, même chemise). L'inconnu grimace, se frotte l'épaule et regarde Amina avec un air mêlant surprise et irritation.
\tcblower
\textbf{Solution détaillée}

\textbf{1. Analyse de l'erreur de perception (Traitement bas-niveau $\to$ Reconnaissance d'identité)}
\begin{itemize}
    \item \textbf{Entrée sensorielle :} Vue de dos (silhouette, vêtements, chevelure). Information partielle (pas de visage).
    \item \textbf{Traitement visuel dorsal (Où/Comment) :} Cortex pariétal postérieur $\to$ analyse de la posture, de la démarche, de la position dans l'espace.
    \item \textbf{Traitement visuel ventral (Quoi/Qui) :} L'information remonte vers le \textbf{gyrus fusiforme (FFA)} et le \textbf{pôle temporal antérieur} (stockage des représentations visuelles des visages/identités connues).
    \item \textbf{Mécanisme d'erreur :} Le cerveau utilise une \textbf{inférence prédictive} (prédiction bayésienne). L'\emph{a priori} (probabilité a priori) : « Koffi est à cette fête, il porte cette chemise ». L'entrée sensorielle partielle (dos, chemise, cheveux) est \textbf{ambigue} mais compatible avec la prédiction. Le système « complète » l'information manquante (le visage) par la représentation stockée de Koffi $\to$ \textbf{Illusion perceptive / Reconnaissance fausse positive}. Le seuil de décision pour « C'est Koffi » a été abaissé par l'attente contextuelle.
\end{itemize}

\textbf{2. Analyse de la réaction de l'inconnu (Perception de l'action d'autrui $\to$ Attribution d'intention)}
\begin{itemize}
    \item \textbf{Observation de l'action :} L'inconnu perçoit un contact physique soudain (tapement épaule) + stimulus auditif fort (cri) + visage d'Amina (sourire, joie).
    \item \textbf{Activation du système miroir :} Ses \textbf{neurones miroirs} (aire 44/45, pariétal inf) s'activent en résonance avec l'action motrice (tapement) et l'expression émotionnelle (joie) d'Amina $\to$ simulation interne de l'état d'Amina (base de l'empathie affective).
    \item \textbf{Attribution d'intention (Théorie de l'Esprit - Niveau haut) :} Le \textbf{TPJ (Jonction Temporo-Pariétale)} et le \textbf{mPFC} traitent la discordance :
    \begin{itemize}
        \item \textit{Contexte :} Je ne connais pas cette fille.
        \item \textit{Action :} Elle me touche familièrement et crie mon nom (ou un nom).
        \item \textit{Inférence :} « Elle me prend pour quelqu'un d'autre » (Fausse croyance d'Amina) OU « Elle est bizarre/agressive ».
    \end{itemize}
    \item \textbf{Résultat comportemental :} L'inconnu attribue à Amina une \textbf{fausse croyance} (elle croit que je suis Koffi). Son irritation (grimace, frottement épaule) traduit la violation de son \textbf{espace personnel} (zone intime < 45 cm) par un inconnu, amplifiée par la surprise. La douleur physique (tapement fort) active la \textbf{voie nociceptive} $\to$ insula antérieure $\to$ composante affective de la douleur.
\end{itemize}

\textbf{3. Résolution de la situation par Amina (Réévaluation / Correction)}
Quand Amina voit le visage (information visuelle complète $\to$ FFA : « Inconnu » $\neq$ « Koffi »), le \textbf{cortex préfrontal latéral} détecte l'erreur (signal d'erreur de prédiction, \emph{Prediction Error}) $\to$ inhibition de la réponse motrice en cours $\to$ activation du \textbf{cortex cingulaire antérieur} (conflit, émotion de gêne) $\to$ excuses, rougeur (activation parasympathique soudaine / withdrawal sympathique).

\textbf{Conclusion}
Cet exemple illustre la nature \textbf{constructive et prédictive} de la perception sociale : Amina a perçu ce qu'elle \emph{s'attendait} à voir (Koffi) et non ce qui était là (un inconnu), du fait d'une entrée sensorielle partielle et d'un \emph{a priori} fort. L'inconnu, lui, a utilisé son \textbf{système miroir} (résonance) et son \textbf{réseau de Théorie de l'Esprit (TPJ/mPFC)} pour inférer l'intention non hostile mais erronée d'Amina, générant une réaction de surprise/défense de son intégrité corporelle.
\end{exoresolu}

\courssub{Méthode 3 : Analyser une situation de communication interindividuelle (verbal, non-verbal, paraverbal)}
\begin{methode}[Démarche pour analyser une situation de communication]
    \textbf{Étape 1 : Identifier les trois canaux de communication.}
    \begin{itemize}
        \item \textbf{Verbal (7\% selon Mehrabian) :} Contenu sémantique, syntaxe, vocabulaire. Aire de \textbf{Broca} (production, syntaxe), Aire de \textbf{Wernicke} (compréhension, sémantique), \textbf{Faisceau arqué} (connexion).
        \item \textbf{Paraverbal (38\%) :} Prosodie (intensité, timbre, débit, pauses, intonation). Hémisphère droit (prosodie émotionnelle), aire de Broca (prosodie linguistique).
        \item \textbf{Non-verbal (55\%) :} Expression faciale (muscles faciaux, nerf facial VII), regard, posture, gestes, distance interpersonnelle (proxémie), contact tactile. \textbf{Hémisphère droit} dominant pour le décodage.
    \end{itemize}
    \textbf{Étape 2 : Vérifier la cohérence (congruence) entre canaux.}
    \begin{itemize}
        \item \textbf{Congruence :} Verbal = Non-verbal $\to$ Communication claire, confiance.
        \item \textbf{Incongruence (Double bind) :} Verbal $\neq$ Non-verbal (ex: « Je vais bien » + visage triste, voix monocorde). Le récepteur privilégie \textbf{le non-verbal} (plus difficile à contrôler volontairement). Source de malaise, méfiance.
    \end{itemize}
    \textbf{Étape 3 : Analyser la boucle de rétroaction (Feedback).} Communication = processus circulaire (Émetteur $\to$ Message $\to$ Récepteur $\to$ Décodage $\to$ Réponse $\to$ Nouvel Émetteur). Rôle des \textbf{neurones miroirs} dans la synchronisation (mimétisme postural, alignement linguistique).
    \textbf{Étape 4 : Identifier les barrières (Bruit).} Bruit physique, physiologique (surdité, fatigue), psychologique (préjugés, émotions fortes, stéréotypes), culturel (codes différents : regard, distance, gestuelle).
\end{methode}

\begin{exoresolu}[Analyse d'un malentendu interculturel en classe]
M. Ngono, professeur de SVT, demande à son élève, Marie, venue d'une zone rurale où le regard direct vers un aîné est perçu comme un manque de respect : « Marie, regarde-moi quand je te parle ! ». Marie baisse les yeux, ne répond pas et reste silencieuse. M. Ngono s'énerve : « Tu me manques de respect ! Sors de la classe ! ».
\tcblower
\textbf{Solution détaillée}

\textbf{1. Décodage de la communication de M. Ngono (Émetteur 1)}
\begin{itemize}
    \item \textbf{Verbal :} Injonction directe (« Regarde-moi »), question rhétorique (« Tu me manques de respect »), ordre (« Sors »).
    \item \textbf{Paraverbal :} Voix forte, intonation montante (colère/autorité), débit rapide.
    \item \textbf{Non-verbal :} Regard direct fixé sur Marie (exigence de contact oculaire), posture redressée, bras croisés ou mains sur hanches (dominance), distance interpersonnelle réduite (intrusion espace personnel/intime).
    \item \textbf{Intention consciente :} Vérifier l'attention, exiger la soumission aux normes scolaires urbaines (regard = respect/honnêteté).
\end{itemize}

\textbf{2. Décodage de la communication de Marie (Récepteur 1 $\to$ Émetteur 2)}
\begin{itemize}
    \item \textbf{Contexte culturel (A priori / Schéma cognitif) :} Dans sa culture d'origine, \textbf{baisser les yeux} face à un aîné/autorité = \cle{marque de respect}, soumission, humilité. Lever les yeux = défi, insolence, égalité inappropriée.
    \item \textbf{Traitement de l'injonction :} M. Ngono demande « Regarde-moi ». Pour Marie, obéir = \textbf{transgresser sa norme culturelle profonde} (manquer de respect). Ne pas obéir = \textbf{transgresser la norme scolaire} (désobéissance).
    \item \textbf{Réponse Non-verbale :} \textbf{Baisse les yeux} (comportement de respect selon son code), immobilité, silence.
    \item \textbf{Réponse Verbale :} \textbf{Silence} (marque de déférence, ne pas contredire l'aîné).
    \item \textbf{État émotionnel :} Peur, confusion, sentiment d'injustice (double contrainte : si je regarde = irrespectueux ; si je ne regarde pas = puni).
\end{itemize}

\textbf{3. Analyse de l'incongruence et de l'escalade (Boucle de rétroaction défaillante)}
\begin{center}
\begin{tabularx}{\linewidth}{|l|X|X|}\hline
\textbf{Niveau} & \textbf{M. Ngono (Code Urbain/Scolaire)} & \textbf{Marie (Code Rural/Traditionnel)} \\ \hline
\textbf{Regard direct} & \textbf{Attendu} : Attention, honnêteté, respect. & \textbf{Interdit} : Défi, arrogance, manque de respect. \\ \hline
\textbf{Regard baissé} & \textbf{Interprété} : Culpabilité, mensonge, désintérêt, défi passif. & \textbf{Produit} : Respect, soumission, écoute. \\ \hline
\textbf{Silence} & \textbf{Interprété} : Révolte mutine, refus de communiquer. & \textbf{Produit} : Déférence, attente de la fin de la réprimande. \\ \hline
\end{tabularx}
\end{center}

\textbf{4. Mécanismes neuronaux du malentendu}
\begin{itemize}
    \item \textbf{M. Ngono :} Perception du non-verbal de Marie (regard bas, silence) $\to$ \textbf{STS / Amygdale} (détection signal social « anormal / menaçant ») $\to$ \textbf{TPJ/mPFC} (Attribution d'intention : « Elle me défie » / « Elle ment ») $\to$ Colère (Orthosympathique) $\to$ Sanction.
    \item \textbf{Marie :} Perception de la colère (visage rouge, voix forte, intrusion spatiale) $\to$ \textbf{Amygdale} (Peur) $\to$ \textbf{Cortex préfrontal} (Inhibition de l'action « lever les yeux » car interdite par la norme interne) $\to$ Figement (Freezing) / Silence.
\end{itemize}

\textbf{5. Conclusion et résolution}
Le conflit naît d'une \textbf{incompatibilité des codes non-verbaux} (proxémie, occulesique) non explicités. M. Ngono a commis une \textbf{erreur d'attribution fondamentale} : il a attribué le comportement de Marie à sa personnalité/disposition (« elle est insolente ») plutôt qu'à la situation/culture (« elle suit sa norme de respect »).
\textbf{Solution pédagogique :} M. Ngono doit expliciter \textbf{le code scolaire} (« Ici, au lycée, regarder le prof dans les yeux prouve qu'on écoute et qu'on n'a pas peur ») sans dévaloriser le code de Marie (« Chez toi, baisser les yeux est une marque de respect, c'est bien aussi »). Cela permet à Marie de \textbf{changer de cadre de référence} (code-switching) sans perte d'identité.
\end{exoresolu}

\courssub{Méthode 4 : Expliquer la régulation hormonale des comportements sociaux (ocytocine, vasopressine, testostérone, cortisol)}
\begin{methode}[Démarche pour expliquer la modulation hormonale des comportements sociaux]
    \textbf{Étape 1 : Identifier l'hormone clé et sa source.}
    \begin{itemize}
        \item \textbf{Ocytocine (OT) :} Synthétisée dans \textbf{noyaux paraventriculaire (PVN) et supraoptique (SON) de l'hypothalamus} $\to$ stockée/libérée par \textbf{neurohypophyse} (axones) + libération centrale (dendritique/somatique dans SNC). Rôle : \cle{attachement}, \cle{confiance}, \cle{reconnaissance sociale}, \cle{réduction anxiété/stress} (inhibe amygdale), \cle{comportement maternel}.
        \item \textbf{Vasopressine (AVP/ADH) :} Même origine hypothalamique. Rôle : \cle{territorialité}, \cle{agressivité défensive}, \cle{liaison de couple} (mâle), \cle{mémoire sociale} (chez rongeurs), régulation eau.
        \item \textbf{Testostérone (T) :} Gonades (testicules), surrénales. Rôle : \cle{dominance}, \cle{competition}, \cle{libido}, modulation de l'amygdale (réactivité à la menace/colère), \textbf{effet « challenge »} (montée avant compétition). Interaction avec cortisol (Dual Hormone Hypothesis : T favorise dominance seulement si Cortisol bas).
        \item \textbf{Cortisol :} Corticosurrénales (Axe HHS). Stress chronique $\to$ \textbf{retrait social}, anxiété, altération mémoire hippocampique, inhibition axe reproducteur.
    \end{itemize}
    \textbf{Étape 2 : Préciser le mode d'action (Récepteurs).} Récepteurs couplés aux protéines G (OTR, V1aR, V1bR, Récepteur aux androgènes intracellulaire $\to$ génomique lent + membranaire rapide). \textbf{Densité et distribution des récepteurs} = spécificité de l'action (ex: densité V1aR dans le septum ventral $\to$ monogamie chez le campagnol des prairies).
    \textbf{Étape 3 : Décrire l'interaction avec le système nerveux.}
    \begin{itemize}
        \item \textbf{Action centrale :} Projections hypothalamiques vers noyau accumbens (récompense), amygdale (peur), hippocampus (mémoire), cortex préfrontal (contrôle). Modulent la \textbf{salience} des stimuli sociaux.
        \item \textbf{Action périphérique $\to$ Centrale :} Signaux viscéraux (cœur, intestin) via nerf vague $\to$ NTS $\to$ Hypothalamus/Locus cœruleus (interoception).
    \end{itemize}
    \textbf{Étape 4 : Contextualiser (Gène $\times$ Environnement).} Polymorphismes gènes récepteurs (ex: \textit{OXTR} rs53576, \textit{AVPR1A} RS3) $\times$ Expériences précoces (soins maternels, stress) $\to$ Phénotype comportemental variable. \textbf{Épigénétique} (méthylation ADN récepteurs).
\end{methode}

\begin{exoresolu}[Le « Jeu de la Confiance » (Trust Game) et l'ocytocine]
Dans une expérience d'économie comportementale (Jeu de la Confiance), l'Investisseur (A) reçoit 10 000 FCFA. Il peut en donner une partie au Gestionnaire (B). La somme donnée est triplée. B peut alors rendre une partie à A.
Deux groupes : Groupe 1 reçoit une pulvérisation nasale d'\textbf{ocytocine (24 UI)} ; Groupe 2 reçoit un \textbf{placebo} (sérum physiologique). Résultats : Le transfert moyen d'argent de A vers B est significativement plus élevé dans le Groupe OT que dans le Groupe Placebo. En revanche, dans une version « risque non-social » (lotterie contre ordinateur), l'OT n'augmente pas les mises.
\tcblower
\textbf{Solution détaillée}

\textbf{1. Identification des variables}
\begin{itemize}
    \item \textbf{Variable indépendante :} Administration intranasale d'Ocytocine (OT) vs Placebo (double aveugle).
    \item \textbf{Variable dépendante :} Montant transféré par l'Investisseur (A) au Gestionnaire (B) = \textbf{indicateur comportemental de la confiance} (prise de risque social).
    \item \textbf{Condition contrôle :} Jeu de loterie (risque non-social, même probabilité de gain/perte).
\end{itemize}

\textbf{2. Voie d'administration et accès au SNC}
\begin{itemize}
    \item \textbf{Voie intranasale :} Contournement de la barrière hémato-encéphalique (BHE) via le \textbf{nerf olfactif} (bulbe olfactif $\to$ cortex piriforme/amygdale) et le \textbf{nerf trijumeau} $\to$ diffusion dans le LCR et le tissu cérébral (aire préoptique, amygdale, noyau accumbens) en $\sim$ 30-45 min. Pic plasmatique aussi mais effet central direct majeur.
\end{itemize}

\textbf{3. Mécanisme neurobiologique de l'effet de l'OT sur la confiance}
\begin{itemize}
    \item \textbf{Cible 1 : Amygdale (Noyau central / Basolatéral).} L'OT se lie aux \textbf{Récepteurs à l'Ocytocine (OTR)} $\to$ \textbf{Inhibition de l'activité de l'amygdale} (via interneurones GABAergiques). $\to$ \textbf{Réduction de la peur de la trahison / de l'aversion au risque social} (réduction de l'anxiété anticipatoire). L'amygdale code normalement la « menace » sociale (risque de perte, trahison).
    \item \textbf{Cible 2 : Système de récompense (VTA $\to$ Noyau Accumbens / Striatum ventral).} L'OT module la transmission \textbf{dopaminergique} (D2) dans le noyau accumbens. $\to$ \textbf{Augmentation de la valeur subjective (valence positive) de l'interaction sociale} / de la coopération. L'acte de « faire confiance » devient plus gratifiant (récompense sociale).
    \item \textbf{Cible 3 : Cortex Préfrontal (mPFC / DLPFC).} Amélioration de l'\textbf{inférence mentale (Théorie de l'Esprit)} : meilleure modélisation des intentions de B (« Il va me rendre de l'argent car c'est rationnel pour lui »).
    \item \textbf{Cible 4 : Insula antérieure.} Réduction de l'activité (insula = dégoût, aversion à l'inéquité, sentiment de « se faire avoir »).
\end{itemize}

\textbf{4. Interprétation de l'absence d'effet sur la loterie (Risque non-social)}
\begin{itemize}
    \item La loterie n'implique \textbf{pas d'agent social} (pas d'intention, pas de réciprocité possible).
    \item L'amygdale et le réseau de la Théorie de l'Esprit (TPJ, mPFC) sont peu sollicités.
    \item L'OT module spécifiquement la \textbf{salience et la valence des stimuli sociaux} (visages, intentions, interactions), pas le traitement mathématique pur du risque (probabilités, montants) qui relève davantage du cortex préfrontal dorsolatéral et de l'insula (calcul risque/bénéfice abstrait).
    \item \textbf{Conclusion :} L'OT ne rend pas « idiot » ou « joueur » de manière générale ; elle \textbf{sélectivement facilite l'approche sociale} en réduisant la peur (amygdale) et en augmentant la récompense attendue (striatum) \textbf{uniquement dans un contexte interindividuel}.
\end{itemize}

\textbf{5. Conclusion}
L'expérience démontre que l'\textbf{ocytocine favorise spécifiquement la confiance interpersonnelle} (prise de risque social) en agissant sur le \textbf{circuit amygdale-récompense-préfrontal}. Elle « désactive » le frein de la peur de l'autre (amygdale) et « active » le moteur de la récompense sociale (striatum), sans affecter l'appréciation du risque pur (lotterie). Cela confirme son rôle d'\textbf{hormone de l'affiliation} et non simplement de « l'amour » ou de la naïveté.
\end{exoresolu}

\courssub{Méthode 5 : Analyser la diversité des comportements humains (Facteurs biologiques, psychologiques, sociaux)}
\begin{methode}[Démarche pour analyser la diversité comportementale et fonder le respect]
    \textbf{Étape 1 : Identifier les trois niveaux déterminants (Modèle Biopsychosocial).}
    \begin{itemize}
        \item \textbf{Biologique (Génotype/Phénotype) :} Génétique (polymorphismes \textit{MAOA}, \textit{5-HTTLPR}, \textit{DRD4}, \textit{OXTR}), Neurobiologie (tempérament, réactivité amygdalienne, sensibilité dopaminergique), Hormonal (testostérone prénatale/pubertaire, cycle menstruel), Pathologie (TSA, TDAH, troubles de l'humeur).
        \item \textbf{Psychologique (Développement/Apprentissage) :} Attachement (sécure/insécure), Traumatismes, Styles cognitifs, Personnalité (Big Five), Régulation émotionnelle, Estimé de soi.
        \item \textbf{Social / Culturel / Environnemental :} Normes culturelles (individualisme/collectivisme, distance pouvoir), Socialisation (famille, école, pairs), Statut socio-économique, Discriminations, Événements de vie.
    \end{itemize}
    \textbf{Étape 2 : Comprendre les interactions (Épigénétique / Plasticité).}
    \begin{itemize}
        \item \textbf{Gène $\times$ Environnement (G$\times$E) :} Ex: Allèle court \textit{5-HTTLPR} $\times$ Maltraitance enfance $\to$ Risque dépression majoré. Allèle long $\to$ Résilience.
        \item \textbf{Épigénétique :} Environnement (stress, soin, nutrition) $\to$ Méthylation ADN / Histones $\to$ Expression gènes (ex: \textit{NR3C1} récepteur glucocorticoïde, \textit{BDNF}) $\to$ Phénotype comportemental stable mais réversible.
        \item \textbf{Plasticité cérébrale :} Expériences $\to$ Remodelage synaptique (élagage, myélinisation) tout au long de la vie.
    \end{itemize}
    \textbf{Étape 3 : Distinguer « Différence » et « Pathologie / Déviance ».}
    \begin{itemize}
        \item \textbf{Neurodiversité :} Variation naturelle du fonctionnement cérébral (TSA, TDAH, HPI, Dyslexie) $\neq$ maladie à guérir, mais différence à accommoder (inclusion).
        \item \textbf{Relativisme culturel :} Comportement « normal » ici = « déviant » là-bas (ex: expression colère, deuil, regard, toucher). Éviter l'\textbf{ethnocentrisme}.
    \end{itemize}
    \textbf{Étape 4 : Fonder l'attitude de respect (Savoir-être).}
    \begin{itemize}
        \item \textbf{Epistémique :} Reconnaître l'opacité de l'esprit d'autrui (je ne connais pas son génotype, son histoire, sa culture). Humilité cognitive.
        \item \textbf{Morale :} Principe de dignité (Kant : traiter autrui comme une fin, jamais comme un moyen). Droits de l'homme.
        \item \textbf{Pragmatique :} La diversité = richesse adaptative pour le groupe (résilience collective, innovation, division du travail social).
    \end{itemize}
\end{methode}

\begin{exoresolu}[Cas d'un élève « perturbateur » : Analyse biopsychosociale et attitude éducative]
Kevin, 16 ans, 1ère A. Comportements : Agitation constante, difficultés à rester assis, réponses impulsives en classe, conflits fréquents avec pairs (bagarrades), notes en baisse. Dossier : Père absent, mère dépressive, placement en foyer à 8 ans (négligences), sommeil irrégulier, consommation occasionnelle de cannabis. Diagnostic médical récent : \textbf{TDAH (Trouble Déficit de l'Attention/Hyperactivité) non traité}.
\tcblower
\textbf{Solution détaillée}

\textbf{1. Analyse multiniveau des déterminants (Modèle Biopsychosocial)}

\begin{center}
\begin{tabularx}{\linewidth}{|l|X|}\hline
\textbf{Niveau} & \textbf{Facteurs identifiés chez Kevin} \\ \hline
\textbf{Biologique} & \begin{itemize}\item \textbf{Génétique :} Hérédité forte TDAH (père ?). Polymorphismes \textit{DRD4-7R} (récepteur D4 faible affinité $\to$ recherche de nouveauté), \textit{DAT1} (recapture DA rapide $\to$ signal DA court). \item \textbf{Neurodéveloppement :} Retard maturation cortex préfrontal (contrôle inhibiteur, planification), hypoactivité striatale (récompense), dysfonctionnement réseau \textbf{Mode Par Défaut (DMN)} $\leftrightarrow$ Réseau Exécutif (difficulté à basculer « repos » $\to$ « tâche »). \item \textbf{Toxicologique :} Cannabis $\to$ agoniste CB1 $\to$ perturbation maturation préfrontale (myélinisation, élagage), amplification impulsivité/psychose risque. \item \textbf{Sommeil :} Privation $\to$ dysfonction préfrontale accrue, irritabilité. \end{itemize} \\ \hline
\textbf{Psychologique} & \begin{itemize}\item \textbf{Attachement :} Insécure désorganisé (négligence + placement). Modèle interne : « Je ne vaux rien », « Les adultes sont inaccessibles/dangereux ». \item \textbf{Régulation émotionnelle :} Déficitaire (alexithymie, impulsivité). Colère = seule émotion accessible/exprimée. \item \textbf{Estime de soi :} Très basse. Comportement perturbateur = \textbf{stratégie de protection} (mieux passer pour « méchant/bête » que pour « nul »), recherche d'attention (même négative). \item \textbf{Traumatisme :} Stress toxique précoce $\to$ recalibration axe HHS (réactivité au stress altérée : hyper ou hypo-réactif). \end{itemize} \\ \hline
\textbf{Social / Environnemental} & \begin{itemize}\item \textbf{Famille :} Monoparental fragile, charge mentale mère, absence figure paternelle (modèle identification, autorité). \item \textbf{Institutionnel :} Ruptures (placements), étiquette « élève difficile », manque de continuité éducative. \item \textbf{Pairs :} Rejet par pairs « normatifs » $\to$ affiliation groupe déviant (renforcement comportements antisociaux, cannabis). \item \textbf{École :} Exigences attentionnelles (cours magistraux longs) inadaptées au profil TDAH $\to$ échec scolaire $\to$ désengagement. \end{itemize} \\ \hline
\end{tabularx}
\end{center}

\textbf{2. Interactions Gène $\times$ Environnement (Épigénétique)}
\begin{itemize}
    \item \textbf{Vulnérabilité génétique (TDAH) + Adversité précoce (négligence, placement)} $\to$ \textbf{Effet cumulatif}. Le stress toxique précoce a pu méthyler le promoteur du gène \textit{NR3C1} (récepteur glucocorticoïde) dans l'hippocampe $\to$ \textbf{feedback négatif du cortisol altéré} $\to$ réactivité au stress prolongée $\to$ aggravation symptômes TDAH (impulsivité, labilité émotionnelle).
    \item \textbf{Cannabis + Prédisposition génétique (ex: COMT Val/Met)} $\to$ Risque majoré de troubles psychotiques ou aggravation troubles attentionnels.
\end{itemize}

\textbf{3. Erreurs d'attribution fréquentes (Pièges à éviter pour l'enseignant)}
\begin{itemize}
    \item \textbf{Erreur d'attribution dispositionnelle (Fondamentale) :} « Kevin est paresseux / méchant / irrespectueux / cherche les ennuis ». $\to$ \textbf{Faux}. Il \emph{ne peut pas} (déficit contrôle inhibiteur) plus qu'il \emph{ne veut pas}.
    \item \textbf{Effet Pygmalion négatif :} Attentes basses $\to$ interactions négatives $\to$ confirmation prophétique.
    \item \textbf{Négligence du handicap invisible :} TDAH = handicap cognitif (fonctions exécutives), pas simple « turbulence ».
\end{itemize}

\textbf{4. Attitude de respect et ajustements raisonnables (Savoir-être / Savoir-faire professionnel)}
\begin{itemize}
    \item \textbf{Stance :} « Kevin a un fonctionnement différent (neurodiversité) + un fardeau lourd. Mon rôle n'est pas de le « corriger » mais de \textbf{lui donner les béquilles} pour réussir ».
    \item \textbf{Aménagements (PAP - Plan d'Accompagnement Personnalisé) :}
    \begin{itemize}
        \item Place stratégique (devant, près prof, loin distractions).
        \item Fractionnement tâches, consignes courtes/écrites/orales, temps majoré évaluations.
        \item Droit au mouvement (distributeur feuilles, effacer tableau, « pause active »).
        \item Renforcement positif immédiat (jeton, sourire, note effort) $>$ Punition retardée.
    \end{itemize}
    \item \textbf{Relationnel :} Alliance thérapeutique/éducative. Un adulte ressource stable, bienveillant, prévisible (figure d'attachement de substitution sécurisante). Ne pas prendre les attaques personnellement (dépersonnalisation).
    \item \textbf{Orientation :} Bilan neuropsychologique complet $\to$ Traitement médicamenteux (Méthylphénidate : bloqueur recapture DA/NA $\to$ stimulation préfrontale) si indiqué + Psychothérapie (TCC, compétences sociales) + Soutien social (assistante sociale).
\end{itemize}

\textbf{Conclusion}
Le comportement de Kevin n'est pas un « choix » moral mais l'\textbf{expression phénotypique} d'une interaction complexe entre une \textbf{vulnérabilité neurodéveloppementale (TDAH)}, un \textbf{traumatisme précoce} ayant sculpté son cerveau et son psychisme (attachement, axe stress), et un \textbf{environnement actuel} (école, pairs, famille) inadapté à ses besoins spécifiques. Le respect d'autrui, ici, impose de \textbf{sortir du jugement moral} pour entrer dans la \textbf{compréhension clinique et l'action pédagogique adaptée} (inclusion), seule voie pour restaurer son agency (pouvoir d'agir) et sa dignité.
\end{exoresolu}

\courssub{Méthode 6 : Modéliser une interaction sociale par la Théorie des Jeux (Dilemme du Prisonnier / Coopération)}
\begin{methode}[Démarche pour modéliser une interaction sociale (Approche éthologique/économique)]
    \textbf{Étape 1 : Formaliser la situation en Matrice de Gains (Payoff Matrix).}
    \begin{itemize}
        \item Identifier les \textbf{acteurs} (2 joueurs minimum).
        \item Identifier les \textbf{stratégies} possibles pour chacun (ex: Coopérer C / Trahir D).
        \item Quantifier les \textbf{gains (Payoffs)} pour chaque combinaison de stratégies : $R$ (Récompense mutuelle), $T$ (Tentation trahir), $S$ (Pigeon/Sucker), $P$ (Punition mutuelle).
        \item Condition \textbf{Dilemme du Prisonnier (PD)} : $T > R > P > S$ et $2R > T + S$ (la coopération mutuelle est optimal pour le groupe).
    \end{itemize}
    \textbf{Étape 2 : Déterminer l'Équilibre de Nash (Stratégie dominante).}
    \begin{itemize}
        \item En \textbf{one-shot (coup unique)} : \textbf{Trahir (D)} est la stratégie dominante (équilibre de Nash) $\to$ Résultat $(P, P)$ sous-optimal.
        \item En \textbf{itéré (répété)} : Possibilité de \textbf{réciprocité} (Stratégie \textbf{Tit-for-Tat / Donnant-Donnant} : Coopère au 1er coup, puis imite le coup adverse). Conditions : Probabilité de rencontre future ($w$) suffisamment élevée ($w > (T-R)/(T-P)$).
    \end{itemize}
    \textbf{Étape 3 : Intégrer les mécanismes proximaux (Neurobiologie/Hormones).}
    \begin{itemize}
        \item \textbf{Coopération :} Activation \textbf{Striatum ventral / VTA} (Récompense), \textbf{Ocytocine} (Confiance), \textbf{mPFC/TPJ} (Théorie de l'Esprit, calcul réputation).
        \item \textbf{Trahison / Punition :} \textbf{Amgydale} (Colère/Peur), \textbf{Insula} (Dégoût/Aversion inéquité), \textbf{DLPFC} (Contrôle cognitif pour inhiber impulsion égoïste ou punir).
        \item \textbf{Testostérone / Cortisol :} Modulent la sensibilité au statut/dominance vs affiliation.
    \end{itemize}
    \textbf{Étape 4 : Analyser l'évolution (Sélection de parentèle / Réciprocité indirecte / Sélection de groupe).} Pourquoi la coopération existe-t-elle malgré le dilemme ? \cle{Règle de Hamilton} ($rB > C$), \cle{Réciprocité directe} (Axelrod), \cle{Réciprocité indirecte} (Réputation), \cle{Punition altruiste}.
\end{methode}

\begin{exoresolu}[Analyse d'une situation de « travail de groupe » en classe]
Quatre élèves (A, B, C, D) doivent réaliser un exposé commun noté sur 20 (note de groupe). Le travail total demande 10h. Si tous travaillent (Coopération), chacun fournit 2.5h, note = 18/20. Si un seul « profite » (Trahison = 0h), les 3 autres font 3.33h chacun, note = 16/20 (fatigue, moins de temps). Le profiteur a 16/20 pour 0h. Si 2 profitent, note = 12/20 pour les 2 travailleurs (5h chacun). Si personne ne travaille, note = 0.
\tcblower
\textbf{Solution détaillée}

\textbf{1. Modélisation : Dilemme du Prisonnier à N joueurs (Bien public / Public Goods Game)}
Simplification à 2 joueurs (Moi vs « Le Reste du groupe ») pour analyser l'incitation individuelle.
\begin{itemize}
    \item \textbf{Stratégies :} $C$ = Travailler (Coopérer), $D$ = Ne rien faire (Trahir/Profiter).
    \item \textbf{Gains (Utilité = Note - Coût Effort)}. Supposons : Note 18 = 18 pts utilité, Note 16 = 16, Note 12 = 12, Note 0 = 0. Coût effort : 2.5h = -2 pts, 3.33h = -3 pts, 5h = -5 pts, 0h = 0 pts.
\end{itemize}

\begin{center}
\begin{tabular}{|c|c|c|}\hline
\textbf{Moi $\backslash$ Groupe} & \textbf{Groupe Coopère (3 travaillent)} & \textbf{Groupe Trahit (0-1 travaillent)} \\ \hline
\textbf{Je Coopère (Travaille)} & $R = 18 - 2,5 = \mathbf{15,5}$ & $S = 12 - 5 = \mathbf{7}$ (ou $0-5=-5$) \\ \hline
\textbf{Je Trahis (Profite)} & $T = 16 - 0 = \mathbf{16}$ & $P = 0 - 0 = \mathbf{0}$ \\ \hline
\end{tabular}
\end{center}

\textbf{Vérification conditions Dilemme :} $T(16) > R(15,5) > P(0) > S(7 \text{ ou } -5)$. \textbf{Condition respectée.}
\textbf{Stratégie dominante :} \textbf{Trahir (Profiter)}. Quel que soit le choix du groupe, mon gain est plus élevé si je ne travaille pas ($16 > 15,5$ et $0 > 7$ ou $-5$).
\textbf{Équilibre de Nash :} (Trahir, Trahir) $\to$ Gain $(0, 0)$. \textbf{Tragédie des communs :} Le groupe échoue (note 0) alors que la coopération mutuelle donnait 15,5 à chacun.

\textbf{2. Pourquoi la coopération émerge-t-elle pourtant souvent ? (Mécanismes proximaux et ultimes)}

\begin{itemize}
    \item \textbf{Répétition (Itération) :} Le groupe se revoit pour d'autres devoirs, examens, vie lycéenne. $w$ (probabilité future interaction) $\approx 1$. Stratégie \textbf{Donnant-Donnant (Tit-for-Tat)} : Je commence par travailler. Si l'autre travaille $\to$ Je continue. Si l'autre profite $\to$ J'arrête (punition) au devoir suivant. \textbf{Condition de stabilité :} $w > \frac{T-R}{T-P} = \frac{16-15,5}{16-0} = \frac{0,5}{16} \approx 0,03$. Très facile à satisfaire $\to$ Coopération stable possible.
    \item \textbf{Réciprocité indirecte / Réputation :} « Profiteur » étiqueté $\to$ Exclu des futurs groupes, perte amitiés, mauvaise réputation profs. Coût social $C_{rep} \gg 0$. Gain réel Trahir = $16 - C_{rep}$. Si $C_{rep} > 0,5$, $R > T$ $\to$ Coopération devient dominante.
    \item \textbf{Punition altruiste (Sanction) :} Les travailleurs peuvent sanctionner le profiteur (conflit, signalement prof, exclusion chat groupe). Coût pour punisseur ($c$), coût pour puni ($p$). Si $p$ élevé, Trahir devient coûteux. \textbf{Base neuronale :} Punition active \textbf{Striatum dorsal} (récompense « justice ») + \textbf{DLPFC} (contrôle).
    \item \textbf{Parentèle / Amitié forte :} Si $r$ (parenté/affinité) élevé. Règle de Hamilton : $r \times B > C$. $B$ = gain pour l'ami (1,5 pts utilité), $C$ = mon coût effort (0,5 pts). $r > 0,5/1,5 = 1/3$. Vrai pour amis proches/frères $\to$ Coopération favorisée.
    \item \textbf{Neurobiologie de la coopération conditionnelle :}
    \begin{itemize}
        \item \textbf{Confiance initiale :} \textbf{Ocytocine} $\to$ Amygdale $\downarrow$ $\to$ Risque social $\downarrow$ $\to$ Premier pas (C).
        \item \textbf{Détection violation :} \textbf{Insula antérieure} + \textbf{Amgydale} $\to$ Signal « Injustice / Trahison » $\to$ Colère/Dégoût.
        \item \textbf{Décision punition/arrêt :} \textbf{DLPFC} intègre signal émotionnel + calcul futur ($w$) $\to$ Décision rationnelle de rétorsition (Tit-for-Tat).
        \item \textbf{Récompense mutuelle :} Réciprocité $\to$ \textbf{Striatum ventral (Dopamine)} $\to$ Renforcement apprentissage « Coopérer paie ».
    \end{itemize}
\end{itemize}

\textbf{3. Conclusion}
Le travail de groupe est un \textbf{Dilemme du Prisonnier itéré}. L'incitation individuelle immédiate est la tricherie (Free-rider). La coopération est maintenue par : (1) l'\textbf{ombre du futur} (répétition, stratégie Tit-for-Tat), (2) la \textbf{réputation} (réciprocité indirecte), (3) la \textbf{punition altruiste} (normes de groupe), (4) l'\textbf{affinité} (sélection de parentèle/amitié). Les \textbf{bases neuro-hormonales} (Ocytocine/Confiance, Dopamine/Récompense, Insula/Amgydale/Sanction, DLPFC/Contrôle) implémentent ces stratégies évolutives. L'enseignant optimise la coopération en : \textbf{répétant les interactions} (groupes stables), \textbf{individualisant la note} (part note individuelle $\to$ réduit $T$), \textbf{rendant visible l'apport} (transparence $\to$ réputation), \textbf{valorisant l'entraide} (norme sociale).
\end{exoresolu}

\courssub{Méthode 7 : Argumenter scientifiquement sur l'origine des différences sexuelles des comportements sociaux}
\begin{methode}[Démarche rigoureuse pour aborder les différences sexuelles (Nature / Culture)]
    \textbf{Étape 1 : Distinguer les niveaux d'analyse (Tinbergen).}
    \begin{itemize}
        \item \textbf{Causalité prochaine (Mécanisme) :} Hormones (testostérone, œstradiol, ocytocine), Neuroanatomie (SDN-POA, INAH3, corps calleux, amygdale), Gènes (chromosomes sexuels $SRY$, échappement inactivation X).
        \item \textbf{Ontogenèse (Développement) :} Organisation prénatale (masculinisation/déféminisation cerveau par T $\to$ Estradiol), Activational pubertaire/adulte, Socialisation différentielle (genres).
        \item \textbf{Fonction (Adaptation) :} Stratégies reproductives différentielles (Investissement parental : Trivers), Sélection sexuelle (concurrence intra-sexe, choix inter-sexe).
        \item \textbf{Phylogénie (Évolution) :} Héritage primate, hominidés.
    \end{itemize}
    \textbf{Étape 2 : Quantifier l'effet (Taille d'effet / Effect Size $d$ de Cohen).}
    \begin{itemize}
        \item La plupart des différences psychologiques/comportementales ont des \textbf{tailles d'effet faibles à modérées} ($d < 0,5$). \textbf{Chevauchement massif des distributions} (ex: Agressivité physique $d \approx 0,5-0,8$ ; Capacités verbales $d \approx -0,1$ ; Rotation mentale 3D $d \approx 0,7$).
        \item \textbf{Variabilité intra-sexe $>$ Variabilité inter-sexe.} Connaître le sexe d'une personne ne permet \textbf{pas} de prédire son comportement individuel avec certitude.
    \end{itemize}
    \textbf{Étape 3 : Identifier les interactions Biologie $\times$ Environnement.}
    \begin{itemize}
        \item \textbf{Hypothèse du cerveau mosaïque :} Pas de « cerveau mâle » ou « femelle » type, mais mosaïque de traits.
        \item \textbf{Effets d'attente / Stéréotypes :} \emph{Stereotype Threat} (menace du stéréotype), \emph{Pygmalion}. L'environnement social \textbf{créé} une partie des différences observées (plasticité).
    \end{itemize}
    \textbf{Étape 4 : Éviter le déterminisme biologique (Naturalistic Fallacy).}
    \begin{itemize}
        \item « Naturel » $\neq$ « Bon / Inévitable / Immuable ». La culture, la loi, l'éducation modulent l'expression des prédispositions.
        \item \textbf{Éthique :} Égalité des droits/opportunités $\neq$ Identité des individus. Respect de la diversité (y compris identité de genre $\neq$ sexe biologique).
    \end{itemize}
\end{methode}

\begin{exoresolu}[Débat : « Les garçons sont naturellement plus agressifs que les filles »]
Un élève affirme : « C'est biologiquement prouvé, la testostérone rend les garçons violents, c'est dans la nature, on ne peut rien y changer. » En vous appuyant sur les connaissances du cours, construisez une argumentation scientifique nuancée pour répondre.
\tcblower
\textbf{Solution détaillée}

\textbf{1. Validation partielle du fait brut (Niveau populationnel / Statistique)}
\begin{itemize}
    \item \textbf{Donnée épidémiologique :} Oui, les statistiques criminelles, scolaires (bagarrades) et expérimentales (paradigmes de provocation) montrent une \textbf{prévalence plus élevée de l'agression physique directe} chez les garçons/hommes (taille d'effet $d \approx 0,5$ à $0,8$, modérée à forte).
    \item \textbf{Base biologique proximale :} Corrélation positive (modérée) entre \textbf{taux de testostérone (T)} circulante (surtout prénatale/pubertaire) et agressivité/dominance. Rôle de la T sur \textbf{amygdale} (réactivité menace $\uparrow$) et \textbf{cortex préfrontal} (contrôle inhibiteur $\downarrow$ lors de forte T). Rôle de l'\textbf{axe HHS} (cortisol bas + T haut = profil « psychopathe »/antisocial).
\end{itemize}

\textbf{2. Nuances critiques : Le déterminisme n'est pas écrit}

\begin{center}
\begin{tabularx}{\linewidth}{|l|X|X|}\hline
\textbf{Argument « Nature »} & \textbf{Contre-arguments / Complexité} & \textbf{Preuves / Mécanismes} \\ \hline
\textbf{Testostérone = Agression} & \textbf{Faux. T = Recherche de Statut / Dominance.} & T augmente la motivation pour le statut. La \textbf{forme} (agression physique vs leadership vs compétition sportive vs altruisme public) dépend du \textbf{contexte social} et des \textbf{normes de genre}. Expérience : T administrée $\to$ comportement \textbf{prosocial} si cela augmente le statut (ex: jeu de l'ultimatum équitable). \\ \hline
\textbf{Cerveau « mâle » violent} & \textbf{Cerveau Mosaïque.} & Pas de dimorphisme sexuel catégoriel pour « circuit agressivité ». Différences moyennes volume amygdale / connexions PFC-amygdale, mais \textbf{chevauchement immense}. Plasticité : Entraînement méditation / régulation émotionnelle modifie ces circuits chez les deux sexes. \\ \hline
\textbf{Universel culturel} & \textbf{Variabilité culturelle massive.} & Taux d'homicides M/F ratio varie de 3:1 à 10:1 selon pays. Sociétés « pacifiques » (ex: Semai Malaisie, Inuit traditionnels) : socialisation \textbf{active de la non-violence chez les garçons} $\to$ quasi disparition différence. \textbf{Socialisation différentielle} : « Les garçons ne pleurent pas », valorisation force, jeux bagarre, modèles médiatiques violents. \\ \hline
\textbf{Agression = Physique} & \textbf{Biais de définition.} & Filles/Femmes : \textbf{Agression relationnelle/indirecte} (exclusion, rumeurs, cyberharcèlement) \textbf{aussi fréquente, voire plus}. Différence = \textbf{Style} (physique vs relationnel), pas forcément \textbf{Niveau} d'hostilité. Liée à socialisation (interdiction agressivité physique filles) + stratégie (force physique inférieure). \\ \hline
\textbf{Immuable} & \textbf{Plasticité / Éducation.} & Programmes \textbf{SEL (Apprentissage Socio-Émotionnel)} réduisent agressivité physique garçons \textbf{et} relationnelle filles. Éducation non-genrée, modèles non-violents, sanction cohérente $\to$ changement normes $\to$ changement comportement. \\ \hline
\end{tabularx}
\end{center}

\textbf{3. Synthèse argumentée (Modèle Biopsychosocial Intégré)}
L'affirmation « C'est dans la nature, on ne peut rien y changer » commet la \textbf{faute naturaliste} (confondre \emph{est} et \emph{doit être}, et surestimer la fixité biologique).
\begin{enumerate}
    \item \textbf{Prédisposition (Biais) :} La biologie (T prénatale/pubertaire, circuits amygdale/PFC) crée une \textbf{probabilité accrue} (biais statistique) pour l'agression physique chez le mâle (stratégie évolutive ancestrale : concurrence mâles pour accès femelles/ressources).
    \item \textbf{Expression (Réalisation) :} Cette prédisposition s'exprime \textbf{uniquement si l'environnement le permet/encourage} (socialisation « garçons durs », absence figures paternelles régulatrices, médias violents, inégalités, stress).
    \item \textbf{Modulation (Contrôle) :} Le \textbf{cortex préfrontal} (maturité $\sim$ 25 ans) permet l'\textbf{inhibition top-down} des impulsions amygdaliennes. L'éducation, la régulation émotionnelle, les normes de groupe développent ce frein.
    \item \textbf{Diversité :} La variabilité \textbf{entre individus d'un même sexe} est bien plus grande que la différence \textbf{moyenne entre sexes}. Réduire un garçon à « futur violent » est une \textbf{erreur probabiliste grave} (stéréotype) qui nuit à son développement (prophétie auto-réalisatrice).
\end{enumerate}

\textbf{Conclusion}
La testostérone ne « rend » pas violent ; elle \textbf{module la sensibilité aux défis de statut}. La violence est un \textbf{comportement appris et choisi} (même inconsciemment) dans un répertoire permis par la biologie mais \textbf{sélectionné par la culture}. Dire « c'est naturel » pour justifier l'inaction éducative est une \textbf{erreur scientifique (réductionnisme)} et une \textbf{faute morale (abdication de responsabilité collective)}. Le respect d'autrui passe par la reconnaissance de sa \textbf{liberté et de sa plasticité}, au-delà de ses déterminants biologiques.
\end{exoresolu}

\begin{attention}[Les 5 erreurs qui coûtent des points au Probatoire]
    \textbf{1. Confondre « Corrélation » et « Causalité » (surtout Hormones/Comportement).}
    \textit{Erreur :} « La testostérone cause l'agressivité. » / « L'ocytocine cause l'amour. »
    \textit{Correct :} « La testostérone \textbf{module} la réactivité aux défis de statut / la dominance ; l'agressivité n'est qu'une expression possible selon le contexte. » « L'ocytocine \textbf{facilite} l'affiliation et la confiance en réduisant l'anxiété sociale (amygdale), elle ne détermine pas le contenu de la relation. » \textbf{Mots-clés :} Modulation, Facilitation, Contexte-dépendance, Interaction Gène$\times$Environnement.

    \textbf{2. Réductionnisme biologisant (Négliger le psychosocial).}
    \textit{Erreur :} Expliquer un comportement complexe (harcelement, délinquance, réussite scolaire) \textbf{uniquement} par le cerveau, les gènes ou les hormones.
    \textit{Correct :} Utiliser systématiquement le \textbf{Modèle Biopsychosocial} : Prédisposition biologique $\times$ Histoire développementale (attachement, traumatismes) $\times$ Contexte social actuel (normes, pairs, inégalités). Citer l'\textbf{épigénétique} comme interface.

    \textbf{3. Essentialiser les différences sexuelles / Ignorer le chevauchement.}
    \textit{Erreur :} « Les filles sont empathiques, les garçons sont logiques. » / « C'est inné. »
    \textit{Correct :} « En moyenne, une différence statistiquement significative existe pour [trait X] ($d = 0,XX$), mais les distributions se \textbf{chevauchent à plus de 70\%}. La variabilité intra-sexe est majeure. L'effet du genre (socialisation) est souvent supérieur à l'effet du sexe (biologie). » Ne jamais essentialiser un individu.

    \textbf{4. Oublier la distinction Proximal / Ultime (Tinbergen) dans l'argumentation évolutionniste.}
    \textit{Erreur :} « Nous sommes égoïstes car nos gènes sont égoïstes (Dawkins) $\to$ donc la morale est impossible. »
    \textit{Correct :} Distinguer : \textbf{Cause ultime (Fonction) :} Pourquoi ce comportement a été sélectionné ? (ex: Réciprocité $\to$ fitness inclusif). \textbf{Cause prochaine (Mécanisme) :} Comment ça marche ici/maintenant ? (ex: Ocytocine $\to$ Amygdale $\to$ Confiance). L'explication ultime ne \textbf{ne justifie pas} moralement le comportement proximal (Naturalistic Fallacy).

    \textbf{5. Imprécision anatomique / vocabulaire flou (Perte de rigueur).}
    \textit{Erreurs classiques :} « Le cerveau limbique » (vague $\to$ préciser : Amygdale, Hippocampe, Hypothalamus, Cingulaire). « Les hormones » (vague $\to$ préciser : Adrénaline vs Cortisol vs Ocytocine vs Testostérone ; voie sanguine vs projection centrale). « Le système nerveux » (vague $\to$ Orthosympathique vs Parasympathique vs Somatique ; Cortex Préfrontal vs Amygdale). \textbf{Exigence Probatoire :} Nommer la \textbf{structure}, la \textbf{molécule}, le \textbf{récepteur}, la \textbf{voie} (nerf vague, faisceau arqué, tractus cortico-spinal).
\end{attention}

\newpage

\coursec{Exercices}

\courssub{Je vérifie mes connaissances}

\begin{exercice}[QCM : les bases physiologiques des comportements]
Pour chaque question, choisis la (ou les) bonne(s) réponse(s).
\begin{enumerate}
\item Un comportement interindividuel est :
\begin{enumerate}
\item un comportement qui se déroule à l'intérieur d'un seul individu ;
\item un comportement qui met en relation au moins deux individus ;
\item un comportement exclusivement inné.
\end{enumerate}
\item L'amygdale est une structure du système limbique impliquée dans :
\begin{enumerate}
\item la digestion ; \item les émotions, notamment la peur ; \item la vision des couleurs.
\end{enumerate}
\item Le message hormonal, comparé au message nerveux, est :
\begin{enumerate}
\item plus rapide et plus bref ; \item plus lent mais plus durable ; \item toujours transporté par les neurones.
\end{enumerate}
\item L'adrénaline est sécrétée par :
\begin{enumerate}
\item l'hypophyse ; \item la thyroïde ; \item les glandes surrénales.
\end{enumerate}
\item L'aire de Broca, située dans l'hémisphère gauche chez la plupart des droitiers, est impliquée dans :
\begin{enumerate}
\item la production du langage ; \item l'équilibre ; \item la mémoire des visages.
\end{enumerate}
\end{enumerate}
\end{exercice}

\begin{exercice}[Vrai ou faux ? Justifie ta réponse]
Réponds par \emph{vrai} ou \emph{faux}, puis justifie chaque réponse en une ou deux phrases.
\begin{enumerate}
\item La perception d'autrui fait intervenir uniquement les organes des sens, sans aucun rôle du cerveau.
\item Une émotion s'accompagne toujours de réactions physiologiques (accélération du rythme cardiaque, sudation...).
\item Le système nerveux autonome (sympathique et parasympathique) agit sous le contrôle de la volonté.
\item L'ocytocine intervient dans les comportements d'attachement, par exemple entre la mère et son nouveau-né.
\item Tous les individus réagissent de la même manière à une même situation sociale.
\end{enumerate}
\end{exercice}

\begin{exercice}[Définitions et structures]
\begin{enumerate}
\item Définis les termes suivants : \cle{comportement interindividuel} ; \cle{perception} ; \cle{émotion} ; \cle{communication}.
\item Cite trois structures du \cle{système limbique} et précise le rôle de chacune dans les comportements.
\item Cite deux modes de communication interindividuelle chez l'Homme et donne un exemple de chacun dans la vie quotidienne camerounaise.
\end{enumerate}
\end{exercice}

\begin{exercice}[Comparaison message nerveux / message hormonal]
Recopie et complète le tableau comparatif suivant :
\begin{center}
\begin{tabularx}{\linewidth}{|l|X|X|}
\hline
\rowcolor{popblueL} \textbf{Critère} & \textbf{Message nerveux} & \textbf{Message hormonal} \\
\hline
Nature du message & \ldots & \ldots \\
\hline
Voie de transport & \ldots & \ldots \\
\hline
Vitesse de transmission & \ldots & \ldots \\
\hline
Durée d'action & \ldots & \ldots \\
\hline
Exemple impliqué dans une émotion & \ldots & \ldots \\
\hline
\end{tabularx}
\end{center}
\end{exercice}

\courssub{Je m'entraîne}

\begin{exercice}[Exercice à données : le cœur qui s'emballe]
Lors d'une cérémonie de remise de bulletins dans un lycée de Yaoundé, on mesure le rythme cardiaque d'un élève de Première A à trois moments :
\begin{center}
\begin{tabular}{|l|c|c|c|}
\hline
\rowcolor{popblueL} \textbf{Moment} & Avant l'annonce des résultats & Pendant l'annonce de son nom & 20 minutes après \\
\hline
\textbf{Rythme cardiaque (battements/min)} & 72 & 118 & 75 \\
\hline
\end{tabular}
\end{center}
\begin{enumerate}
\item Calcule le pourcentage d'augmentation du rythme cardiaque entre le premier et le deuxième moment.
\item Explique ces variations en mobilisant le rôle du \cle{système nerveux autonome} (sympathique) et de l'\cle{adrénaline}.
\item Pourquoi le rythme cardiaque revient-il à une valeur proche de la normale 20 minutes après ? Cite le système nerveux impliqué.
\end{enumerate}
\end{exercice}

\begin{exercice}[Schéma à compléter : du visage menaçant à la fuite]
Un élève aperçoit soudain, dans la cour du lycée, un chien agressif qui se dirige vers lui ; il s'enfuit aussitôt. Reproduis et complète la chaîne suivante en plaçant aux bons endroits les éléments de la liste : \emph{glandes surrénales ; rétine (récepteur) ; amygdale ; nerfs sensitifs ; cortex cérébral ; adrénaline ; muscles des membres}.
\[ \text{Stimulus visuel} \longrightarrow \boxed{\phantom{XXXX}} \longrightarrow \boxed{\phantom{XXXX}} \longrightarrow \text{thalamus} \longrightarrow \boxed{\phantom{XXXX}} \longrightarrow \boxed{\phantom{XXXX}} \]
\[ \longrightarrow \text{hypothalamus} \longrightarrow \boxed{\phantom{XXXX}} \longrightarrow \boxed{\phantom{XXXX}} \longrightarrow \boxed{\phantom{XXXX}} \longrightarrow \text{fuite} \]
\begin{enumerate}
\item Légende ton schéma en précisant, pour chaque élément, s'il relève du message nerveux ou du message hormonal.
\item Explique pourquoi la réaction de fuite est d'abord très rapide, puis accompagnée d'une activation générale de l'organisme (cœur, respiration, muscles).
\end{enumerate}
\end{exercice}

\begin{exercice}[Tableau à construire : les hormones des comportements sociaux]
Construis un tableau à quatre colonnes (\textbf{Hormone} ; \textbf{Glande productrice} ; \textbf{Rôle comportemental principal} ; \textbf{Exemple de situation de vie courante}) et complète-le pour les quatre hormones suivantes : \cle{adrénaline}, \cle{cortisol}, \cle{ocytocine}, \cle{testostérone}. Pour chaque exemple de situation, choisis un contexte de la vie quotidienne au Cameroun (examen, accouchement, compétition sportive, stress prolongé...).
\end{exercice}

\begin{exercice}[Analyse de document : un cas d'aphasie]
Au Centre Hospitalier Universitaire de Yaoundé, un patient de 58 ans est admis après un accident vasculaire cérébral (AVC) ayant endommagé une région du lobe frontal de l'\textbf{hémisphère gauche}. Le médecin note : « Le patient comprend parfaitement ce qu'on lui dit et exécute correctement les ordres simples, mais il ne parvient plus à produire un discours articulé : ses phrases sont réduites à quelques mots, prononcés avec difficulté. »
\begin{enumerate}
\item Nomme le trouble du langage présenté par ce patient.
\item Quelle aire cérébrale est vraisemblablement lésée ? Précise sa localisation et sa fonction.
\item Déduis de cette observation le rôle des aires du langage dans la communication interindividuelle.
\item Explique pourquoi une altération de la communication peut modifier les relations sociales de ce patient.
\end{enumerate}
\end{exercice}

\courssub{Je me prépare au Probatoire}

\begin{exercice}[Sujet type Probatoire : la peur lors d'une épreuve orale]
Lors d'un exposé oral devant la classe, une élève de Première A présente les signes suivants : mains moites, cœur qui bat vite (110 battements/min contre 70 au repos), voix tremblante, rougeur du visage. On lui prélève un échantillon sanguin juste après l'exposé : le taux d'adrénaline et de cortisol y est nettement élevé.
\begin{enumerate}
\item Définis les termes \cle{émotion} et \cle{comportement interindividuel}, puis montre que la situation décrite en est une illustration.
\item Identifie, parmi les signes observés, ceux qui relèvent d'une réponse nerveuse rapide et ceux qui relèvent d'une réponse hormonale. Justifie.
\item Décris, sous forme de chaîne fonctionnelle ordonnée, le trajet de l'information depuis la perception du public (regards des camarades) jusqu'à la sécrétion d'adrénaline, en citant au moins quatre structures nerveuses et une glande.
\item Le cortisol est appelé « hormone du stress ». Explique l'intérêt de sa sécrétion pour l'organisme, puis le danger d'une sécrétion prolongée (stress chronique des examens, insomnies...).
\item Propose deux conseils fondés sur des arguments physiologiques pour aider cette élève à mieux maîtriser son stress avant un oral.
\end{enumerate}
\end{exercice}

\begin{exercice}[Sujet type Probatoire : diversité des réactions à une provocation]
Pendant un match de football inter-classes dans un lycée de Douala, un joueur commet une faute sur un adversaire. Deux coéquipiers du joueur fautif réagissent différemment à la même provocation verbale de l'arbitre : le premier (Âmandou, 17 ans) reste calme et discute ; le second (Basile, 17 ans) s'énerve, crie et manque de frapper l'arbitre. L'enquête de l'infirmier scolaire révèle que Basile a mal dormi depuis plusieurs nuits, vit des tensions familiales et consomme parfois des boissons énergisantes avant les matchs.
\begin{enumerate}
\item Rappelle ce qu'est un comportement interindividuel et cite les deux grands systèmes de régulation qui en constituent le support physiologique.
\item En t'appuyant sur le rôle du système limbique (amygdale, cortex préfrontal), explique comment un même stimulus peut déclencher deux réactions différentes.
\item À partir du document, identifie deux \textbf{facteurs biologiques} et deux \textbf{facteurs psychosociaux} pouvant expliquer la différence de comportement entre les deux adolescents. Justifie chaque choix.
\item Explique le rôle possible de la testostérone et du cortisol dans la réaction de Basile.
\item En tant qu'élève éduqué à la citoyenneté, propose deux attitudes concrètes de respect et de gestion pacifique des conflits face à la diversité des comportements.
\end{enumerate}
\end{exercice}

\begin{exercice}[Sujet type Probatoire : communication et attachement — étude de documents]
\textbf{Document 1.} Dans une maternité de Bafoussam, on observe que les mères qui pratiquent le contact peau à peau avec leur nouveau-né dans l'heure qui suit la naissance présentent un taux sanguin d'ocytocine significativement plus élevé, et que le lien d'attachement mère-enfant s'établit plus rapidement (allaitement facilité, apaisement mutuel).

\textbf{Document 2.} On mesure chez des élèves la fréquence cardiaque moyenne dans trois situations de communication : conversation amicale (68 battements/min), dispute (105 battements/min), silence seul dans une salle (70 battements/min).
\begin{enumerate}
\item Définis la \cle{communication interindividuelle} et cite, en les illustrant par des exemples camerounais, deux formes de communication non verbale.
\item À partir du document 1, dégage le rôle de l'ocytocine dans les comportements sociaux. Précise la glande qui la libère dans la circulation sanguine.
\item Explique pourquoi le contact peau à peau peut être considéré comme un stimulus à l'origine d'un message hormonal.
\item Exploite le document 2 : construis un histogramme représentant la fréquence cardiaque moyenne en fonction de la situation, puis interprète les variations observées en faisant intervenir le système nerveux autonome et l'adrénaline.
\item Conclus en montrant, à l'aide des deux documents, que les relations interindividuelles reposent sur des mécanismes à la fois nerveux et hormonaux.
\end{enumerate}
\end{exercice}

\newpage

\coursec{Corrigés des exercices}

\courssub{Je vérifie mes connaissances}

\begin{corrige}[Exercice 1 — QCM : les bases physiologiques des comportements]
\begin{enumerate}
\item \textbf{b.} Un comportement interindividuel met obligatoirement en relation au moins deux individus (émetteur et récepteur). C'est la définition même de l'interindividualité.
\item \textbf{b.} L'amygdale est une structure clé du système limbique, centre de traitement des émotions, particulièrement impliquée dans la reconnaissance de la peur et le déclenchement des réponses de défense (fuite, combat).
\item \textbf{b.} Le message hormonal (transport sanguin) est plus lent à mettre en place (secondes à minutes) que l'influx nerveux (millisecondes), mais son action est plus diffuse et plus durable (minutes à heures).
\item \textbf{c.} L'adrénaline (ou épinéphrine) est sécrétée par la médullosurrénale (partie centrale des glandes surrénales), sous stimulation du système nerveux sympathique.
\item \textbf{a.} L'aire de Broca (aire motrice du langage, circonvolution frontale inférieure, hémisphère dominant gauche chez 90\% des droitiers) programme les mouvements articulatoires nécessaires à la production du langage parlé.
\end{enumerate}
\end{corrige}

\begin{corrige}[Exercice 2 — Vrai ou faux ? Justifie ta réponse]
\begin{enumerate}
\item \textbf{Faux.} La perception n'est pas une simple réception sensorielle. Elle nécessite l'intégration et l'interprétation des informations sensorielles par le cortex cérébral (aires sensitives primaires et aires d'association) et le système limbique (reconnaissance affective, mémoire). Sans cerveau, il n'y a pas de perception consciente d'autrui.
\item \textbf{Vrai.} Par définition, une émotion est une réponse brutale, brève et intense à un stimulus, caractérisée par trois composantes indissociables : une composante \textbf{physiologique} (activation du système nerveux autonome : tachycardie, sudation, boule au ventre...), une composante \textbf{comportementale} (fuite, cri, paralysie, expression faciale) et une composante \textbf{subjective} (ressenti conscient : peur, colère, joie).
\item \textbf{Faux.} Le système nerveux autonome (ou végétatif) fonctionne de manière \textbf{involontaire} et inconsciente. Il échappe au contrôle direct de la volonté (on ne peut pas décider volontairement d'arrêter son cœur ou de dilater ses pupilles). Il est régulé par l'hypothalamus et le tronc cérébral.
\item \textbf{Vrai.} L'ocytocine, libérée par la neurohypophyse, est l'hormone majeure de l'attachement social. Elle favorise le lien mère-enfant (accouchement, allaitement, contact peau à peau), la confiance, l'empathie et les comportements coopératifs.
\item \textbf{Faux.} La diversité des comportements humains est une règle biologique et psychosociale. Face à un même stimulus, les réactions diffèrent selon le \textbf{génotype} (tempérament, sensibilité hormonale), l'\textbf{expérience passée} (apprentissage, traumatismes), le \textbf{contexte culturel} (normes, valeurs), l'\textbf{état physiologique du moment} (fatigue, stress, faim) et la \textbf{maturation cérébrale} (cortex préfrontal non mature à l'adolescence).
\end{enumerate}
\end{corrige}

\begin{corrige}[Exercice 3 — Définitions et structures]
\begin{enumerate}
\item \textbf{Définitions :}
\begin{itemize}
\item \cle{Comportement interindividuel} : Ensemble des actions, réactions et interactions observables qui se produisent entre au moins deux individus et qui influencent mutuellement leur état physiologique, émotionnel ou social.
\item \cle{Perception} : Processus cognitif par lequel le système nerveux central (cortex sensoriel, aires d'association, système limbique) sélectionne, organise et interprète les informations brutes issues des récepteurs sensoriels pour leur donner un sens conscient (reconnaître un visage, identifier une émotion dans une voix).
\item \cle{Émotion} : Réponse globale, transitoire et intense de l'organisme face à un stimulus signifiant, associant une activation neurovégétative (système nerveux autonome), une expression comportementale (fuite, combat, approche) et un ressenti subjectif (peur, colère, joie, tristesse).
\item \cle{Communication} : Processus d'échange d'informations entre un émetteur et un récepteur via un code partagé (langage verbal, gestes, mimiques, phéromones), permettant d'influencer le comportement ou l'état physiologique de l'autre.
\end{itemize}
\item \textbf{Trois structures du système limbique et leurs rôles :}
\begin{center}
\begin{tabularx}{\linewidth}{|l|X|}
\hline
\rowcolor{popblueL} \textbf{Structure} & \textbf{Rôle dans les comportements} \\
\hline
Amygdale & Détection de la valeur émotionnelle des stimuli (surtout menace/peur), déclenchement des réponses autonomes et endocrines (via hypothalamus), mémorisation émotionnelle. \\
\hline
Hippocampe & Mémorisation contextuelle et épisodique (où, quand, quoi), mise en relation du stimulus actuel avec les souvenirs passés pour moduler la réponse émotionnelle. \\
\hline
Hypothalamus & Centre de commande neuroendocrine et autonome : traduit l'émotion en réponse physiologique (activation sympathique/parasympathique, libération hormonale hypophysaire). \\
\hline
Cortex cingulaire & Interface cognition-émotion : régulation de l'affect, prise de décision sociale, détection des erreurs, empathie. \\
\hline
\end{tabularx}
\end{center}
(Any three correct structures with accurate roles are accepted.)
\item \textbf{Deux modes de communication interindividuelle chez l'Homme (exemples camerounais) :}
\begin{itemize}
\item \textbf{Communication verbale (langage oral) :} Discussion entre élèves en langue officielle (français/anglais) ou en langue locale (bassa, bamileke, fulfulde) lors d'un travail de groupe ; exposé oral devant la classe ; négociation des prix au marché (Marché Central de Yaoundé, Marché Congo à Douala).
\item \textbf{Communication non verbale :} \emph{Gestes et mimiques} : salut respectueux (inclinaison du buste, poignée de main douce) face à un aîné ou un enseignant ; regard baissé par politesse ; \emph{Paraverbal} : ton de la voix, débit, silences ; \emph{Proxémie} : distance de conversation plus courte entre amis qu'avec un inconnu.
\end{itemize}
\end{enumerate}
\end{corrige}

\begin{corrige}[Exercice 4 — Comparaison message nerveux / message hormonal]
\begin{center}
\begin{tabularx}{\linewidth}{|l|X|X|}
\hline
\rowcolor{popblueL} \textbf{Critère} & \textbf{Message nerveux} & \textbf{Message hormonal} \\
\hline
Nature du message & Influx nerveux : variation de potentiel membranaire (onde de dépolarisation) + libération de neurotransmetteurs à la synapse (signal électrochimique). & Hormone : molécule messagère chimique (peptide, stéroïde, amine) sécrétée dans le sang. \\
\hline
Voie de transport & Axones des neurones (voie filaire, point à point) + fente synaptique. & Circulation sanguine (voie diffuse, broadcast) + liquide interstitiel. \\
\hline
Vitesse de transmission & Très rapide : de l'ordre de la milliseconde à la centaine de mètres/seconde (fibres myélinisées). & Lente : de quelques secondes à plusieurs minutes pour atteindre la cible via le sang. \\
\hline
Durée d'action & Brève : quelques millisecondes à secondes (désensibilisation rapide des récepteurs, recapture neurotransmetteurs). & Longue : minutes à heures, voire jours (demi-vie plasmatique, effets géniques pour stéroïdes). \\
\hline
Exemple impliqué dans une emotion & Influx sympathique (noradrénaline) sur le cœur (tachycardie immédiate) ; influx moteur vers muscles (fuite). & Adrénaline (médullosurrénale) : amplification et prolongation de la tachycardie, bronchodilatation, glycogénolyse. Cortisol (corticosurrénale) : adaptation au stress prolongé. \\
\hline
\end{tabularx}
\end{center}
\end{corrige}

\courssub{Je m'entraîne}

\begin{corrige}[Exercice 5 — Exercice à données : le cœur qui s'emballe]
\begin{enumerate}
\item \textbf{Calcul du pourcentage d'augmentation :}
\[
\text{Augmentation} = \frac{\text{FC}_2 - \text{FC}_1}{\text{FC}_1} \times 100 = \frac{118 - 72}{72} \times 100 = \frac{46}{72} \times 100 \approx \num{63,9}\% \approx \textbf{64\%}
\]
\item \textbf{Explication physiologique :}
L'annonce des résultats constitue un \textbf{stimulus stressogène} (enjeu social, évaluation). La perception visuelle/auditive (nom appelé) est traitée par le thalamus $\to$ amygdale (détection menace/peur) $\to$ hypothalamus. L'hypothalamus active le \textbf{système nerveux sympathique} (voie courte, rapide) : fibres pré-ganglionnaires $\to$ ganglions chaînes latéro-vertébraux $\to$ fibres post-ganglionnaires libérant de la \textbf{noradrénaline} sur le nœud sinusal (cœur) $\to$ augmentation immédiate de la fréquence cardiaque (FC).
Simultanément, le système sympathique innerve la \textbf{médullosurrénale} (glandes surrénales) qui libère massivement de l'\textbf{adrénaline} dans le sang. L'adrénaline se lie aux récepteurs $\beta_1$-adrénergiques du nœud sinusal et du myocarde, \textbf{potentant et prolongeant} la tachycardie (FC = 118 bpm).
\item \textbf{Retour à la normale (20 min après) :}
La disparition du stimulus stressant et les rétrocontrôles négatifs (barorécepteurs carotidien/aortique détectent l'hypertension artérielle) activent le \textbf{système nerveux parasympathique} (nerf vague / X). Celui-ci libère de l'\textbf{acétylcholine} sur le nœud sinusal (récepteurs muscariniques M2), freinant la dépolarisation spontanée : \textbf{bradycardie relative} (retour à 75 bpm). Parallèlement, l'adrénaline circulante est métabolisée (foie, reins) et sa concentration chute.
\end{enumerate}
\end{corrige}

\begin{corrige}[Exercice 6 — Schéma à compléter : du visage menaçant à la fuite]
\begin{enumerate}
\item \textbf{Chaîne fonctionnelle complétée :}
\[
\text{Stimulus visuel} \longrightarrow \fbox{\textbf{Rétine (récepteur)}} \longrightarrow \fbox{\textbf{Nerfs sensitifs (nerf optique)}} \longrightarrow \text{Thalamus} \longrightarrow \fbox{\textbf{Amygdale}} \longrightarrow \fbox{\textbf{Cortex cérébral (préfrontal, visuel)}}
\]
\[
\longrightarrow \text{Hypothalamus} \longrightarrow \fbox{\textbf{Système nerveux sympathique}} \longrightarrow \fbox{\textbf{Glandes surrénales (médullosurrénale)}} \longrightarrow \fbox{\textbf{Adrénaline (sang)}} \longrightarrow \fbox{\textbf{Muscles des membres}} \longrightarrow \text{Fuite}
\]

\textbf{Légende (Nature du message) :}
\begin{itemize}
\item Rétine, Nerfs sensitifs, Thalamus, Amygdale, Cortex cérébral, Hypothalamus, Système nerveux sympathique (jusqu'à la glande) : \textbf{Message nerveux} (influx électrique + neurotransmetteurs).
\item Glandes surrénales (organe sécréteur), Adrénaline (dans le sang), Cibles vasculaires/musculaires : \textbf{Message hormonal} (transport sanguin).
\item Muscles des membres : \textbf{Effecteur} (réponse motrice finale).
\end{itemize}
\item \textbf{Explication de la biphasie de la réaction :}
\begin{itemize}
\item \textbf{Phase 1 : Très rapide (voie basse « courte »)}. L'information visuelle va du thalamus \textbf{directement vers l'amygdale} (sans passer par le cortex visuel conscient). L'amygdale active instantanément l'hypothalamus et le tronc cérébral $\to$ système sympathique $\to$ réaction de fuite immédiate (survie). Temps : $\sim$ 10-50 ms.
\item \textbf{Phase 2 : Activation générale durable (voie haute « longue » + hormonale)}. L'information passe aussi par le thalamus $\to$ cortex visuel $\to$ cortex préfrontal (évaluation consciente : « c'est un chien, il est agressif, je dois fuir »). L'hypothalamus commande la médullosurrénale $\to$ \textbf{adrénaline dans le sang}. L'adrénaline assure une \textbf{mobilisation générale et prolongée} : cœur (FC, force), poumons (bronchodilatation), foie (glycogénolyse $\to$ glucose), muscles (afflux sanguin), inhibition digestion. Cela prépare l'organisme à l'effort intense et soutenu (lutte ou fuite).
\end{itemize}
\end{enumerate}
\end{corrige}

\begin{corrige}[Exercice 7 — Tableau à construire : les hormones des comportements sociaux]
\begin{center}
\begin{tabularx}{\linewidth}{|l|X|X|X|}
\hline
\rowcolor{popblueL} \textbf{Hormone} & \textbf{Glande productrice} & \textbf{Rôle comportemental principal} & \textbf{Exemple de situation de vie courante au Cameroun} \\
\hline
\cle{Adrénaline} & Médullosurrénale (glandes surrénales) & Médiatrice immédiate de la réaction de \textbf{lutte-fuite} : mobilisation énergétique d'urgence, vigilance accrue, activation cardiovasculaire et respiratoire. & \textbf{Examen oral du Probatoire :} montée brutale du stress à l'entrée dans la salle, cœur qui bat la chamade, mains moites, voix tremblante. \\
\hline
\cle{Cortisol} & Corticosurrénale (zone fasciculée des glandes surrénales) & Hormone du \textbf{stress chronique / adaptation prolongée} : néoglucogenèse (énergie durable), modulation immunitaire, régulation de la réactivité de l'amygdale et du cortex préfrontal. & \textbf{Période de révision intense (3 semaines avant le Bac) :} insomnies, irritabilité, difficulté de concentration, prise de poids abdominale, immunité affaiblie (maladies fréquentes). \\
\hline
\cle{Ocytocine} & Neurohypophyse (synthèse dans noyaux paraventriculaire et supraoptique de l'hypothalamus) & Hormone de l'\textbf{attachement, de la confiance et du lien social} : lien mère-enfant, couple, cohésion de groupe, réduction de l'anxiété sociale (inhibe amygdale). & \textbf{Accouchement à la maternité de Bafoussam :} contact peau-à-peau immédiat, mise au sein facilitée, apaisement mutuel mère-nouveau-né ; réconciliation après un conflit familial autour d'un repas partagé. \\
\hline
\cle{Testostérone} & Testicules (cellules de Leydig) \textit{(aussi corticosurrénale chez la femme)} & Régulation de la \textbf{dominance, compétitivité, agressivité instrumentale, libido, comportements reproducteurs} ; module la réactivité de l'amygdale et l'inhibition préfrontale. & \textbf{Finale de football inter-quartiers (Douala) :} engagement physique intense, recherche de statut/dominance, célébration virile de la victoire ; tensions lors d'un contrôle d'identité perçu comme injuste. \\
\hline
\end{tabularx}
\end{center}
\end{corrige}

\begin{corrige}[Exercice 8 — Analyse de document : un cas d'aphasie]
\begin{enumerate}
\item \textbf{Trouble : Aphasie de Broca} (ou \textbf{aphasie motrice / expressive / non-fluente}). Caractérisée par une production langagière laborieuse, agrammaticale (télégraphique), avec articulation difficile, mais une \textbf{compréhension préservée}.
\item \textbf{Aire lésée : Aire de Broca} (aire 44 et 45 de Brodmann), située dans la \textbf{partie opéculaire et triangulaire de la circonvolution frontale inférieure de l'hémisphère gauche} (hémisphère dominant pour le langage chez 90\% des droitiers). \textbf{Fonction :} programmation motrice du langage (planification des séquences articulatoires), traitement de la syntaxe et de la grammaire, coordination des muscles de la phonation (langue, lèvres, larynx, diaphragme).
\item \textbf{Rôle des aires du langage dans la communication :} La communication verbale repose sur une \textbf{organisation fonctionnelle spécialisée et latéralisée} :
\begin{itemize}
\item Aire de \textbf{Wernicke} (cortex temporo-pariétal postérieur) : décodage, compréhension du sens (intacte ici).
\item Aire de \textbf{Broca} (cortex frontal antérieur) : codage, programmation motrice de l'expression (lésée ici).
\item Faisceau arqué : connexion entre les deux (intègre compréhension et production).
\end{itemize}
Une lésion sélective de Broca dissocie compréhension (intacte) et production (altérée), prouvant que ces deux opérations reposent sur des substrats neuronaux distincts mais interconnectés.
\item \textbf{Conséquences sur les relations sociales :}
\begin{itemize}
\item \textbf{Isolement relationnel :} difficulté à exprimer besoins, opinions, émotions $\to$ frustration, retrait.
\item \textbf{Malentendus :} l'interlocuteur peut sous-estimer les capacités intellectuelles du patient (préservation de l'intelligence).
\item \textbf{Perte de statut social/professionnel :} impossible d'enseigner, de diriger, de négocier (marché, administration).
\item \textbf{Dépendance accrue :} besoin d'un tiers pour communiquer (médecine, banque, famille).
\item \textbf{Risque dépressif :} deuil de l'identité communicante.
\end{itemize}
\end{enumerate}
\end{corrige}

\courssub{Je me prépare au Probatoire}

\begin{corrige}[Exercice 9 — Sujet type Probatoire : la peur lors d'une épreuve orale]
\textbf{Barème indicatif : 20 points}
\begin{enumerate}
\item \textbf{(4 pts) Définitions et illustration.}
\begin{itemize}
\item \cle{Émotion} : Réponse globale, transitoire, intense à un stimulus signifiant, associant composante physiologique (autonome/hormonale), comportementale (expression, fuite/combat) et subjective (ressenti).
\item \cle{Comportement interindividuel} : Interaction observable entre au moins deux individus (émetteur/récepteur) modifiant leur état respectif.
\item \textbf{Illustration :} La situation est un comportement interindividuel (élève face au public/enseignant, enjeu d'évaluation sociale). Elle déclenche une \textbf{émotion de peur/anxiété sociale} (stimulus = regards, évaluation) avec signes physiologiques (tachycardie, sudation, rougeur, voix tremblante) et comportementaux (difficulté à parler, posture figée).
\end{itemize}
\item \textbf{(4 pts) Identification signes nerveux vs hormonaux.}
\begin{center}
\begin{tabularx}{\linewidth}{|l|X|X|}
\hline
\rowcolor{popblueL} \textbf{Signe observé} & \textbf{Voie} & \textbf{Justification} \\
\hline
Voix tremblante, mains moites, rougeur visage & \textbf{Nerveuse (Sympathique)} & Latence très courte (< 1 s). Activation directe : fibres sympathiques $\to$ glandes sudoripares (sudation), vaisseaux cutanés (vasodilatation active $\to$ rougeur), muscles laryngés (tremblements). \\
\hline
FC 110 bpm (maintien), Taux sanguin Adrénaline/Cortisol élevés & \textbf{Hormonale} & Latence plus longue (min). Adrénaline (médullosurrénale) maintient FC élevée. Cortisol (corticosurrénale) mobilise glucose pour effort prolongé. Dosage sanguin = preuve message hormonal. \\
\hline
\end{tabularx}
\end{center}
\item \textbf{(5 pts) Chaîne fonctionnelle (perception $\to$ adrénaline).}
\[
\text{Regards public (Stimulus)} \xrightarrow{\text{Récepteurs}} \text{Rétine} \xrightarrow{\text{Message nerveux}} \text{Nerf optique} \to \text{Thalamus} \to \textbf{Amygdale} \to \textbf{Hypothalamus} \to \textbf{Système nerveux sympathique (Moelle épinière, Ganglions)} \to \textbf{Médullosurrénale (Glande)} \xrightarrow{\text{Sang (Message hormonal)}} \textbf{Adrénaline} \to \text{Cibles (Cœur, Vaisseaux, Foie, Muscles)}
\]
\textit{Structures nerveuses citées (min 4) : Rétine, Nerf optique, Thalamus, Amygdale, Hypothalamus, Moelle épinière, Ganglions sympathiques. Glande : Médullosurrénale (Surrénale).}
\item \textbf{(4 pts) Cortisol : intérêt et danger.}
\begin{itemize}
\item \textbf{Intérêt (adaptation aiguë) :} Néoglucogenèse hépatique $\to$ hyperglycémie (carburant cerveau/muscles) ; lipolyse ; action anti-inflammatoire/immunomodulatrice (prévient excès) ; facilite la mémorisation de l'événement stressant (survie) ; maintient la réactivité vasculaire à l'adrénaline.
\item \textbf{Danger (stress chronique / examens prolongés) :} \textbf{Immunodépression} (infections fréquentes) ; \textbf{Troubles métaboliques} (résistance à l'insuline, diabète type 2, obésité androïde) ; \textbf{Neurotoxicité hippocampique} (atrophie $\to$ troubles mémoire, apprentissage) ; \textbf{Cardiovasculaire} (HTA, athérosclérose) ; \textbf{Psychique} (anxiété généralisée, dépression, burn-out) ; \textbf{Sommeil} (insomnie, fragmentation).
\end{itemize}
\item \textbf{(3 pts) Deux conseils physiologiques.}
\begin{enumerate}
\item \textbf{Respiration lente, profonde, abdominale (cohérence cardiaque) :} Stimule les \textbf{barorécepteurs} et le \textbf{nerf vague (parasympathique)} $\to$ freine l'activité sympathique $\to$ baisse FC, PA, calme l'amygdale via rétrocontrôle. \textit{Argument : activation directe du frein physiologique.}
\item \textbf{Préparation répétée / Simulation (entraînement) :} Renforce les connexions \textbf{Cortex Préfrontal $\to$ Amygdale} (voie inhibitrice GABAergique) $\to$ meilleure \textbf{régulation cognitive de l'émotion} (réévaluation : « c'est un exercice, pas un danger vital »). \textit{Argument : plasticité synaptique, contrôle top-down.}
\end{enumerate}
\textit{Autres acceptables : Sommeil suffisant (régule cortisol), Activité physique régulière (baisse réactivité sympathique basale), Visualisation positive (active circuits récompense/dopamine).}
\end{enumerate}
\textbf{Points de vigilance (Probatoire) :}
\begin{itemize}
\item Distinguer clairement \emph{voie nerveuse} (rapide, ciblée, neurotransmetteur) et \emph{voie hormonale} (lente, diffuse, sang).
\item Citer \textbf{nommément} les structures : Amygdale, Hypothalamus, Système sympathique, Médullosurrénale. « Cerveau » ou « système nerveux » seul = point partiel.
\item Pour le cortisol, opposer \textbf{aigu (adaptatif)} vs \textbf{chronique (pathogène)}.
\item Conseils doivent être \textbf{argumentés physiologiquement} (pas juste « respire » mais « active le parasympathique »).
\end{itemize}
\end{corrige}

\begin{corrige}[Exercice 10 — Sujet type Probatoire : diversité des réactions à une provocation]
\textbf{Barème indicatif : 20 points}
\begin{enumerate}
\item \textbf{(3 pts) Définition et systèmes.}
\begin{itemize}
\item \cle{Comportement interindividuel} : Interaction entre au moins deux individus (ici : joueur/arbitre, coéquipiers/arbitre) impliquant perception, émotion, communication, action.
\item \textbf{Deux grands systèmes de régulation :} \textbf{Système nerveux} (central + autonome) et \textbf{Système endocrinien (hormonal)}.
\end{itemize}
\item \textbf{(5 pts) Rôle système limbique (Amygdale vs Cortex Préfrontal).}
\begin{itemize}
\item \textbf{Amygdale :} Détecte la provocation verbale comme \textbf{menace sociale/agression}. Déclenche instantanément la réponse de peur/colère (voie basse thalamus-amygdale) : activation sympathique, libération adrénaline, préparation combat/fuite.
\item \textbf{Cortex Préfrontal (CPF, zone dorsolatérale et ventromédiane) :} \textbf{Régulation top-down (contrôle inhibiteur).} Évalue le contexte (« c'est un match, l'arbitre a autorité, frapper = exclusion/prudence »), inhibe l'amygdale (via GABA), permet une réponse adaptée (dialogue, self-control).
\item \textbf{Différence Amandou vs Basile :} Chez \textbf{Amandou}, le CPF mature et non fatigué inhibe efficacement l'amygdale $\to$ réponse calme (dialogue). Chez \textbf{Basile}, l'amygdale est \textbf{hyper-réactive} (sensibilisée par stress chronique, cortisol) et/ou le CPF est \textbf{fonctionnellement déconnecté/inhibé} (fatigue, cortisol, caféine, immaturité) $\to$ désinhibition de la colère $\to$ réaction explosive.
\end{itemize}
\item \textbf{(5 pts) Facteurs biologiques et psychosociaux (Document).}
\begin{center}
\begin{tabularx}{\linewidth}{|l|X|X|}
\hline
\rowcolor{popblueL} \textbf{Type} & \textbf{Facteur identifié} & \textbf{Justification physiologique / psychologique} \\
\hline
\textbf{Biologique 1} & \textbf{Manque de sommeil (mal dormi plusieurs nuits)} & Privation de sommeil $\to$ \textbf{hypofonction du Cortex Préfrontal} (métabolisme glucidique réduit, connexions affaiblies) + \textbf{hyperréactivité de l'amygdale} (couplage amygdale-CPF rompu) + élévation du \textbf{cortisol} matinal. Perte du contrôle inhibiteur. \\
\hline
\textbf{Biologique 2} & \textbf{Consommation boissons énergisantes (caféine, taurine)} & Caféine = antagoniste récepteurs \textbf{adénosine A1/A2A} $\to$ libération \textbf{noradrénaline/dopamine} $\to$ excitation centrale, anxiété, tachycardie, \textbf{potentiation de la réponse sympathique} à la provocation. \\
\hline
\textbf{Psychosocial 1} & \textbf{Tensions familiales (stress chronique)} & Stress chronique $\to$ élévation basale \textbf{cortisol} $\to$ \textbf{sensibilisation de l'amygdale} (plasticité synaptique LTP) + \textbf{atrophie dendritique CPF/hippocampe}. Baisse du seuil de déclenchement de la colère (irritabilité trait). \\
\hline
\textbf{Psychosocial 2} & \textbf{Contexte match (enjeu, public, identité de groupe)} & Présence de pairs (facilitation sociale), enjeu de statut/dominance, norme de « virilité » (ne pas se laisser faire) $\to$ activation circuits \textbf{récompense/dopamine} (victoire) et \textbf{menace sociale} (honneur). Pression normative du groupe. \\
\hline
\end{tabularx}
\end{center}
\item \textbf{(4 pts) Rôle Testostérone et Cortisol chez Basile.}
\begin{itemize}
\item \textbf{Testostérone (T) :} Adolescent 17 ans = pic pubertaire T. T module l'amygdale (↑ réactivité menace) et inhibe le CPF (↓ contrôle). Favorise \textbf{dominance, compétitivité, agressivité réactive} face à un défi social (provocation arbitre = menace statut). Interaction T $\times$ Cortisol (dual hormone hypothesis) : T élevé + Cortisol élevé (stress chronique) = \textbf{maximise l'agressivité/dominance} (le cortisol ne freine plus la T par feedback négatif central altéré).
\item \textbf{Cortisol (C) :} Stress chronique (famille, sommeil) $\to$ C basale élevée. Aigu (match) $\to$ pic C. C facilite la consolidation mémorielle de la menace, mobilise l'énergie (glucose) pour l'action (frappe), mais altère le jugement (CPF) et la flexibilité comportementale.
\end{itemize}
\item \textbf{(3 pts) Attitudes citoyennes (Respect / Gestion pacifique).}
\begin{enumerate}
\item \textbf{Communication Non Violente (CNV) / Dialogue :} « J'entends ta frustration, mais frapper l'arbitre nous coûtera le match et des sanctions. Parlons-en calmement à la mi-temps. » (Active CPF, désamorce amygdale pair).
\item \textbf{Respect de l'autorité et des règles (Fair-play) :} Accepter la décision de l'arbitre (règle du jeu acceptée collectivement), serrer la main de l'adversaire après le match. (Intériorisation normes sociales, inhibition impulsivité via CPF).
\end{enumerate}
\textit{Autres : Médiation par le capitaine/coach, Techniques de respiration (cohérence cardiaque) sur le banc, Soutien émotionnel au coéquipier (oxytocine).}
\end{enumerate}
\textbf{Points de vigilance :}
\begin{itemize}
\item Ne pas confondre « facteur biologique » (génétique, physiologie immédiate, substance) et « facteur psychosocial » (environnement, culture, relations, histoire de vie).
\item L'interaction \textbf{Gène $\times$ Environnement} (épigénétique, plasticité) est la clé de la diversité.
\item Citer les structures cérébrales (Amygdale, CPF) et les molécules (Testostérone, Cortisol, Caféine/Adénosine) avec leur mécanisme d'action.
\end{itemize}
\end{corrige}

\begin{corrige}[Exercice 11 — Sujet type Probatoire : communication et attachement — étude de documents]
\textbf{Barème indicatif : 20 points}
\begin{enumerate}
\item \textbf{(4 pts) Définition et communication non verbale.}
\begin{itemize}
\item \cle{Communication interindividuelle} : Processus dynamique d'échange d'informations entre un émetteur et un récepteur via un canal (sensoriel) et un code partagé, intentionnel ou non, modifiant le comportement ou l'état physiologique des partenaires.
\item \textbf{Deux formes non verbales (exemples camerounais) :}
\begin{enumerate}
\item \textbf{Kinésique (Gestes, Mimiques, Posture) :} \textit{Ex:} Salut respectueux à un aîné (inclinaison buste, deux mains pour serrer la main, regard baissé) ; \textit{Ex:} Sourire et hochements de tête lors d'une discussion amicale au « buvette » pour marquer l'accord/écoute.
\item \textbf{Proxémie (Gestion de l'espace) / Haptique (Contact) :} \textit{Ex:} Distance rapprochée, contact épaule/bras entre amis d'enfance (confiance) vs distance formelle avec un supérieur hiérarchique ; \textit{Ex:} Portage du bébé en écharpe (pagne) par la mère ou la grande sœur (contact corporel continu, sécurité).
\end{enumerate}
\end{itemize}
\item \textbf{(3 pts) Rôle Ocytocine (Doc 1).}
\begin{itemize}
\item \textbf{Rôle :} Hormone clé de l'\textbf{attachement} et de la \textbf{confiance sociale}. Facilite la reconnaissance du nouveau-né, déclenche les comportements maternels (allaitement, protection, apaisement), réduit l'anxiété (inhibe amygdale), renforce la synchronie mère-enfant (rythmes cardiaques, affectifs).
\item \textbf{Glande libératrice :} \textbf{Neurohypophyse} (lobe postérieur de l'hypophyse), sur commande des neurones neurosecréteurs de l'\textbf{hypothalamus} (noyaux paraventriculaire et supraoptique).
\end{itemize}
\item \textbf{(3 pts) Contact peau à peau = Stimulus hormonal.}
Le contact peau à peau active les \textbf{mécanorécepteurs cutanés} (disques de Merkel, corpuscules de Ruffini, fibres C-tactiles afférentes non myélinisées) $\to$ influx nerveux via moelle épinière $\to$ thalamus $\to$ \textbf{hypothalamus} (noyaux PVN/SON) $\to$ dépolarisation des terminaisons axonales dans la \textbf{neurohypophyse} $\to$ libération massive d'\textbf{ocytocine dans le sang} (message hormonal). C'est une \textbf{transduction sensoriel $\to$ neuroendocrine}.
\item \textbf{(5 pts) Exploitation Doc 2 : Histogramme et Interprétation.}
\begin{popfigure}
\begin{tikzpicture}
\begin{axis}[
    ybar,
    width=\linewidth,
    height=8cm,
    ylabel={Fréquence cardiaque moyenne (bpm)},
    symbolic x coords={Conversation amicale, Silence seul, Dispute},
    xtick=data,
    ymin=0, ymax=120,
    bar width=2.5cm,
    nodes near coords,
    nodes near coords align={vertical},
    every node near coord/.append style={font=\small, popdark},
    tick label style={font=\small},
    label style={font=\small},
    ymajorgrids=true,
    grid style={popline},
    axis lines*=left,
    enlarge x limits=0.3,
]
\addplot[popblue, fill=popblueL] coordinates {
    (Conversation amicale, 68)
    (Silence seul, 70)
    (Dispute, 105)
};
\end{axis}
\end{tikzpicture}
\legende{Histogramme de la fréquence cardiaque moyenne selon la situation de communication (Données Document 2).}
\end{popfigure}

\textbf{Interprétation :}
\begin{itemize}
\item \textbf{Silence seul (70 bpm) \& Conversation amicale (68 bpm) :} FC proche de la normale au repos. Dominance du \textbf{système nerveux parasympathique} (nerf vague) $\to$ libération d'acétylcholine sur nœud sinusal $\to$ freinage vagal (homéostasie, sécurité, lien social positif). La conversation amicale (lien positif) peut même augmenter légèrement le tonus vagal (effet apaisant ocytocine/ventral vagal).
\item \textbf{Dispute (105 bpm) :} Augmentation massive (+50\% vs repos). Activation aiguë du \textbf{système nerveux sympathique} (perception menace/colère $\to$ amygdale $\to$ hypothalamus) $\to$ noradrénaline cardiaque + \textbf{adrénaline} (médullosurrénale) $\to$ récepteurs $\beta_1$ $\to$ chronotropie et inotropie positives. Préparation combat/fuite. Le contexte interindividuel conflictuel bascule l'équilibre autonomique vers la sympathicotonie.
\end{itemize}
\item \textbf{(5 pts) Conclusion : Nerveux + Hormonal dans les relations.}
Les deux documents démontrent la \textbf{double régulation neuro-hormonale} des relations interindividuelles :
\begin{itemize}
\item \textbf{Doc 1 (Attachement positif) :} Stimulus tactile (nerveux) $\to$ Hypothalamus $\to$ \textbf{Ocytocine (hormonal)} $\to$ Comportement maternel, apaisement, lien durable. \textit{Message hormonal pour la consolidation du lien social à long terme.}
\item \textbf{Doc 2 (Conflit / Stress aigu) :} Stimulus verbal/visuel (nerveux) $\to$ Amygdale/Hypothalamus $\to$ \textbf{Système Nerveux Sympathique (nerveux rapide)} + \textbf{Adrénaline/Cortisol (hormonal)} $\to$ Mobilisation énergétique, vigilance, réaction immédiate. \textit{Message nerveux pour l'urgence, hormonal pour l'amplification et la durée.}
\end{itemize}
\textbf{Synthèse :} Les relations interindividuelles ne sont pas de simples échanges cognitifs ; elles sont \textbf{incarnées}. Le \textbf{système nerveux} assure la \textbf{rapidité, la précision et l'intentionnalité} (perception, décision, langage, geste). Le \textbf{système hormonal} assure la \textbf{diffusion, la durée, la modulation métabolique et l'ancrage mémoriel/émotionnel} (attachement, stress, dominance, confiance). Leur \textbf{intégration au niveau hypothalamique et limbique} permet l'adaptation fine de l'individu à son environnement social.
\end{enumerate}
\textbf{Points de vigilance (Probatoire) :}
\begin{itemize}
\item L'histogramme doit être \textbf{propre} : axes légendés avec unités, barres distinctes, titres.
\item L'interprétation \textbf{doit nommer les systèmes} (Sympathique / Parasympathique) et les \textbf{médiateurs} (Acétylcholine, Noradrénaline, Adrénaline).
\item La conclusion doit \textbf{synthétiser les deux documents} en opposant/complétant les deux types de relations (positif/négatif) et les deux modes de régulation (nerveux/hormonal).
\item Vocabulaire précis : \textit{transduction, neurohypophyse, tonus vagal, sympathicotonie, homéostasie, plasticité.}
\end{itemize}
\end{corrige}

\newpage

\coursec{Fiche bilan}

\begin{popfigure}
\begin{center}\tcbox[colback=popblue,colframe=popblue,coltext=white,arc=9pt,boxsep=4pt,left=6pt,right=6pt]{\bfseries\small Comportements Interindividuels}\end{center}
\vspace{-2pt}
\begin{tcbraster}[raster columns=3,raster equal height=rows,raster column skip=6pt,raster row skip=6pt,enhanced,blankest]
\begin{tcolorbox}[enhanced,colback=poporange,colframe=poporange,coltext=white,boxrule=0.9pt,arc=3pt,left=4pt,right=4pt,top=3pt,bottom=3pt,boxsep=1pt,halign=left]\bfseries\scriptsize 1. Généralités Définitions Types de relations\end{tcolorbox}
\begin{tcolorbox}[enhanced,colback=popgreen,colframe=popgreen,coltext=white,boxrule=0.9pt,arc=3pt,left=4pt,right=4pt,top=3pt,bottom=3pt,boxsep=1pt,halign=left]\bfseries\scriptsize 2. Perception d'autrui Système nerveux Intégration\end{tcolorbox}
\begin{tcolorbox}[enhanced,colback=poppurple,colframe=poppurple,coltext=white,boxrule=0.9pt,arc=3pt,left=4pt,right=4pt,top=3pt,bottom=3pt,boxsep=1pt,halign=left]\bfseries\scriptsize 3. Émotions Bases nerveuses Bases hormonales Expression\end{tcolorbox}
\begin{tcolorbox}[enhanced,colback=poppink,colframe=poppink,coltext=white,boxrule=0.9pt,arc=3pt,left=4pt,right=4pt,top=3pt,bottom=3pt,boxsep=1pt,halign=left]\bfseries\scriptsize 4. Communication Verbal / Non-verbal Langage Signaux\end{tcolorbox}
\begin{tcolorbox}[enhanced,colback=popgold,colframe=popgold,coltext=white,boxrule=0.9pt,arc=3pt,left=4pt,right=4pt,top=3pt,bottom=3pt,boxsep=1pt,halign=left]\bfseries\scriptsize 5. Régulation Hormonale Ocytocine Testostérone Cortisol\end{tcolorbox}
\begin{tcolorbox}[enhanced,colback=popblue!60,colframe=popblue!60,coltext=white,boxrule=0.9pt,arc=3pt,left=4pt,right=4pt,top=3pt,bottom=3pt,boxsep=1pt,halign=left]\bfseries\scriptsize 6. Diversité \& Respect Facteurs bio/psycho Tolérance\end{tcolorbox}
\end{tcbraster}

\legende{Carte mentale récapitulative de la leçon NA03 : Maîtrise des comportements interindividuels (1ère A).}
\end{popfigure}

\begin{bilan}

\textbf{Définitions clés}
\begin{itemize}
  \item \cle{Comportement} : Ensemble des réactions observables d'un individu face à une situation (stimulus interne ou externe).
  \item \cle{Relation interindividuelle} : Interaction dynamique entre au moins deux individus, impliquant perception, émotion et communication.
  \item \cle{Perception sociale} : Processus cognitif par lequel on interprète les caractéristiques, intentions et émotions d'autrui (rôle de l'amygdale, cortex préfrontal, neurones miroirs).
  \item \cle{Émotion} : Réponse physiologique et psychologique brève à un stimulus, préparant l'action (composantes : subjective, neurovégétative, motrice/expressive).
  \item \cle{Communication} : Échange d'informations entre émetteur et récepteur via un code partagé (verbal : langage ; non-verbal : mimiques, posture, prosodie, phéromones).
\end{itemize}

\textbf{Mécanismes physiologiques à connaître par cœur}
\begin{center}
\begin{tabularx}{\linewidth}{|l|X|X|}\hline
\rowcolor{popblueL}
\textbf{Processus} & \textbf{Structures nerveuses principales} & \textbf{Messagers hormonaux clés} \\ \hline
Perception d'autrui & Rétine $\to$ Thalamus $\to$ Cortex visuel ; Amygdale (valence émotionnelle) ; Cortex préfrontal (jugement social) ; \textbf{Neurones miroirs} (aire de Broca, pariétale inf.) & -- \\ \hline
Émotion (Peur/Colère) & \textbf{Amygdale} (déclenchement) ; Hypothalamus ; Tronc cérébral ; Cortex préfrontal (régulation) & \textbf{Adrénaline}, Noradrénaline (médullosurrénale) ; Cortisol (corticosurrénale, axe HHS) \\ \hline
Attachement / Confiance & Système de récompense (Aire tegmentale ventrale, Noyau accumbens) & \textbf{Ocytocine} (neurohypophyse) ; Dopamine \\ \hline
Agressivité / Dominance & Amygdale ; Hypothalamus ; Cortex préfrontal (inhibition) & \textbf{Testostérone} (androgène) ; Sérotonine (bas taux $\leftrightarrow$ impulsivité) \\ \hline
Stress social & Axe \textbf{Hypothalamo-Hypophyso-Surrénalien (HHS)} & \textbf{CRH} $\to$ \textbf{ACTH} $\to$ \textbf{Cortisol} \\ \hline
\end{tabularx}
\end{center}

\textbf{Méthodes express}
\begin{enumerate}
  \item \textbf{Analyser un comportement :} Identifier le stimulus $\to$ Décrire la réponse (motrice, verbale, physiologique) $\to$ Identifier les structures nerveuses (voie sensorielle, intégration, voie motrice) et hormonales impliquées $\to$ Conclure sur la fonction adaptative.
  \item \textbf{Interpréter une situation de communication :} Distinguer verbal / non-verbal $\to$ Repérer les indices d'émotion (mimiques universelles d'Ekman : joie, colère, peur, dégoût, tristesse, surprise) $\to$ Évaluer la congruence message verbal / non-verbal.
  \item \textbf{Argumenter sur la diversité :} Citer facteurs biologiques (génétique, hormones, maturation cérébrale) ET psychosociaux (culture, éducation, expériences, normes sociales) $\to$ Refuser le déterminisme unique $\to$ Prôner le respect (savoir-être).
\end{enumerate}

\end{bilan}

\begin{aretenir}[Checklist avant le Probatoire]
$\square$ Savoir définir : comportement, perception sociale, émotion, communication (verbal/non-verbal), empathie (neurones miroirs).\\
$\square$ Citer les 6 émotions universelles (Ekman) et leur expression faciale caractéristique.\\
$\square$ Décrire la voie nerveuse de la peur : Stimulus $\to$ Thalamus $\to$ \textbf{Amygdale} (voie courte/rapide) ET Cortex (voie longue/précise) $\to$ Hypothalamus $\to$ Réponse neurovégétative (Adrénaline).\\
$\square$ Expliquer le rôle de l'\textbf{ocytocine} (attachement, confiance, lien mère-enfant, lien social) et de la \textbf{testostérone} (agressivité, dominance, compétition) en précisant : \emph{modulation, pas déterminisme strict}.\\
$\square$ Schématiser l'axe \textbf{HHS} (CRH $\to$ ACTH $\to$ Cortisol) et donner 2 effets du cortisol chronique sur le cerveau (hippocampe, mémoire) et le système immunitaire.\\
$\square$ Différencier \textbf{système nerveux} (rapide, ponctuel, influx nerveux) et \textbf{système hormonal} (lent, diffus, sanguin, durée longue) dans la régulation des comportements.\\
$\square$ Citer 3 facteurs biologiques (ex: génétique, taux hormonal, maturation cortex préfrontal) et 3 facteurs psychosociaux (ex: culture, famille, pairs, médias) de la diversité des comportements.\\
$\square$ Savoir analyser un document (texte, image, courbe hormonale) pour en extraire le support physiologique d'une relation interindividuelle (ex: pic d'ocytocine lors d'un câlin, pic de cortisol lors d'un conflit).\\
$\square$ Rédiger une conclusion éthique : la biologie explique les potentialités, l'environnement module l'expression $\to$ Respect de la singularité, lutte contre les stéréotypes et discriminations.\\
$\square$ Connaître l'expérience type : \emph{Test du miroir (reconnaissance de soi)}, \emph{Expérience du visage immobile (Still Face - Tronick)}, ou \emph{Jeu de l'ultimatum (neuroéconomie)}.\\
$\square$ Maîtriser le vocabulaire précis : \cle{neurotransmetteur}, \cle{hormone}, \cle{récepteur}, \cle{synapse}, \cle{axe HHS}, \cle{neurones miroirs}, \cle{empathie cognitive/affective}.
\end{aretenir}

\findecours

\end{document}
